\documentclass[11pt]{amsart}
\usepackage{amsmath,amssymb,amsthm,mathtools,mathrsfs}
\usepackage[margin=1.08in]{geometry}
\usepackage{booktabs}
\usepackage{array}
\usepackage{longtable}
\usepackage{tabularx}
\usepackage{hyperref}
\hypersetup{hidelinks}
\usepackage[final,expansion=false]{microtype}
\setlength{\parskip}{4pt plus 1pt}
\emergencystretch=2.5em

\newtheorem{theorem}{Theorem}[section]
\newtheorem{proposition}[theorem]{Proposition}
\newtheorem{lemma}[theorem]{Lemma}
\newtheorem{corollary}[theorem]{Corollary}
\theoremstyle{definition}
\newtheorem{definition}[theorem]{Definition}
\newtheorem{problem}[theorem]{Problem}
\newtheorem{conjecture}[theorem]{Conjecture}
\theoremstyle{remark}
\newtheorem{remark}[theorem]{Remark}

\newcommand{\C}{\mathbb C}
\newcommand{\R}{\mathbb R}
\newcommand{\N}{\mathbb N}
\newcommand{\Z}{\mathbb Z}
\newcommand{\Q}{\mathbb Q}
\newcommand{\A}{\mathbb A}
\newcommand{\dd}{\,d}
\newcommand{\RH}{\mathrm{RH}}
\newcommand{\QW}{\mathcal Q_W}
\newcommand{\Div}{\operatorname{Div}}
\newcommand{\divv}{\operatorname{div}}
\newcommand{\Prin}{\operatorname{Prin}}
\newcommand{\Tr}{\operatorname{Tr}}
\newcommand{\Rel}{\operatorname{Re}}
\newcommand{\Spec}{\operatorname{Spec}}
\newcommand{\Gr}{\mathcal G}
\newcommand{\Pic}{\operatorname{Pic}}
\newcommand{\GL}{\operatorname{GL}}
\newcolumntype{Y}{>{\raggedright\arraybackslash}X}

\begin{document}

\title[An arithmetic Lefschetz programme over $\Z$]
{An Arithmetic Lefschetz Programme over $\Z$:\\
A Metrized Bivariant Correspondence Square,\\
its Determinant Pairing and its Corner Trace\\[4pt]
\large Conditional row-(d) evaluation draft}

\author{Alain Connes}
\thanks{Independent Researcher. \texttt{alain@connes.org}}

\date{\today}

\begin{abstract}
We construct the first three ingredients of Weil's argument over $\Z$ in an
explicit periodic--cohomological--nuclear setting.  The noncollapsed
spherical square carries intrinsic periodic section cohomology, a
coefficient-one determinant and derived prime contact.  Witt Frobenius gives
a faithful multiplicative correspondence family whose dynamic contact is
$[\Z\xrightarrow{\Phi_n(1)}\Z]$ and has degree $\Lambda(n)$.  Its monoidal
realization in Meyer's Poisson quotient yields the complete nuclear
Lefschetz character, including the finite, polar and archimedean terms.

For the fourth ingredient we identify the primitive operator, prove
full-space positivity for $0<T\leq\log2$, and construct at every arithmetic
threshold the finite Gamma--Tate amplitude, conservative Tate terminal,
positive unreset Gamma storage and exact correlated rational leakage.  The
reciprocal half-integer band theorem pays every connected leakage cluster
and reduces disconnected-band recombination to an exact constant/mean-zero
Schur block.  We isolate its remaining content as the Logarithmic Schur
Angle Conjecture.  Thus the geometric and operator-theoretic construction of
row (d) is complete, while its global positivity is conditional on that
single conjecture.  Rows (a)--(c) and the certified initial range of row (d)
are unconditional; the global Hodge inequality and the Riemann Hypothesis
remain open unless the conjecture is proved.
\end{abstract}


\maketitle

\section{Where this comes from}

\subsection{The two companion papers}

The companion paper \emph{Branch Selection and Global Polarization in the
Connes--CCM Route to RH: Finite Prime--Gamma Inertia, Tate Boundary
Geometry, and the Divisor-Sensitivity Obstruction} develops the finite
semilocal chain of Connes, Consani and Moscovici \cite{CCM2025} together
with Suzuki's screw-function realization \cite{Suzuki2025,SuzukiScrew}, and
reduces $\RH$ to a strict inequality on a cofinal family of finite rows.  Its
own contribution is a branch-selection theorem and an inertia calculus ---
Haynsworth additivity, Schur complements, scalar innovations, an adaptive
Krylov hierarchy --- carried out on the finite prime--Gamma coordinate, and
a Tate boundary-geometry analysis of the archimedean fibre.  It closes with
an unconditional theorem, requiring nothing beyond Hardy's 1914 result: the
ordinary absolutely continuous completion destroys exactly the resonant data
the argument needs.  Its concluding section names two open programmes: a
geometric intersection product with a Hodge index theorem, and a
divisor-sensitive resonant completion.

The second companion, \emph{Obstructions to Weil Positivity: A Catalogue of
No-Go Theorems, Falsifiers and Countermodels in the Prime--Gamma Doob
Coordinate}, does the complementary work: it records, with proofs, the
mechanisms which cannot supply the missing inequality.

Both converge on the same reading.  The obstruction is not a gap in the
analysis.  It is that a source object with the structure of a continuum is
repeatedly asked to factor through a target with a finiteness constraint ---
atomic support against a dense group, point spectrum against absolutely
continuous spectrum, a divisible group against a finitely generated one.
The present paper takes the first of the two named programmes.

\subsection{What this paper attempts}

Weil's proof of the Riemann Hypothesis for function fields
\cite{Weil1948,Weil1952} rests on four ingredients over $\mathbb F_q$:

\begin{enumerate}
\item[(a)] the surface $C\times_{\mathbb F_q}C$;
\item[(b)] the Frobenius correspondences $\Gamma_n\subset C\times C$;
\item[(c)] the identity $\Gamma_n\cdot\Delta=N_n$ between intersection
numbers and the zeta counting function;
\item[(d)] the Hodge index theorem, in the Castelnuovo--Severi form.
\end{enumerate}

Together these give $|a_n|\le 2g\sqrt q$.  The programme of this paper is to
ask what each of (a)--(d) would be over $\Z$, and to build them.  The
ambient framework is that of Connes and Consani
\cite{CC-ArithmeticSite,CC-ScalingSite,CC-RRstrategy,CC-RRSpecZ}, whose
Riemann--Roch strategy adapts Weil's argument to characteristic zero along
the lines of Mattuck--Tate and Grothendieck
\cite{MattuckTate,Grothendieck}, and whose Scaling Site supplies the
correspondences $\Psi_\lambda$ we use throughout.

Two structural facts fix the shape of the programme before any construction
begins.

\begin{itemize}
\item \textbf{(b) is not imported.}  Deninger's foliated
dynamical systems for arithmetic schemes \cite{Deninger2024} produce, for
$\Spec\Z$, a flow whose periodic orbits are in bijection with the primes,
of length $\log p$.  Morishita \cite{Morishita2026} constructs a continuous,
Galois-equivariant, flow-anti-equivariant map from these systems to the
Connes--Consani adelic spaces, sending the periodic packet at $p$ onto the
corresponding circle.  But that map is built from rational Witt vectors, and
at a point over $p$ the residue field is $\overline{\mathbb F}_p$, whose
multiplicative group $\mu_{(p)}$ has order prime to $p$.  Every such map
therefore annihilates the $p$-adic component --- which is precisely the
transverse direction whose kernel trace produces the local factor
$|1-u|_p^{-1}$.  The bridge transports which orbits exist and nothing about
how they meet.  Section~\ref{sec:correspondences} instead builds row~(b)
inside row~(a), using the Witt Frobenius and the cotangent contact already
native to the arithmetic curve.

\item \textbf{(d) cannot be imported from the explicit formula.}  One may
compute the signature of the corner pairing directly and find that it is
equivalent to $\RH$; that computation is recorded in \S3.5.  An equivalence
derived from the explicit formula does no work, because Weil's Hodge index
is an \emph{input} from geometry which then \emph{implies} the zeta bound.
It cannot be obtained by reading the desired sign back from the explicit
formula.  A proof must instead come either from an independent geometric
index theorem or from a source-defined positive factorization whose
construction does not use that sign.
\end{itemize}

Accordingly the burden falls first on (a).  Section~\ref{sec:geometric-a}
constructs and audits the intrinsic periodic--Deligne--nuclear spherical
package, including its total coefficient cohomology, normalized determinant,
independent code certificate and reduced derived contact.
Section
\ref{sec:correspondences} constructs row~(b) on that carrier.  The
intervening sections study the analytic corner trace, and
Section~\ref{sec:rowc} constructs the nuclear cohomological completion that
proves row~(c).  The fourth-row construction and its final open inequality are stated in
Section~\ref{ss:gammafouriertate}.

\subsection{The four rows over $\mathbb F_q$ and over $\Z$}
\label{ss:map}

The following table is the map the rest of the paper fills in.  The left
half is Weil's, complete since 1948; the right half is what the present
programme has over $\Z$, with its status stated as sharply as the results
below allow.  It is included because the distance between the two halves,
and its precise shape, is the actual subject of this work.

\begingroup
\footnotesize
\setlength{\tabcolsep}{4pt}
\renewcommand{\arraystretch}{1.25}
\noindent
\begin{tabularx}{\textwidth}{@{}c p{0.205\textwidth} p{0.36\textwidth} Y@{}}
\toprule
 & \textbf{Over $\mathbb F_q$ (Weil)} & \textbf{Over $\Z$ (this programme)}
 & \textbf{Status over $\Z$} \\
\midrule
(a) &
The surface $C\times_{\mathbb F_q}C$: smooth, projective, with finitely
generated Néron--Severi group. &
The noncollapsed spherical square $\mathscr Y_{\mathbb S}$; invertible
Arakelov line objects; intrinsic periodic coefficient cohomology on the
canonical integer-radius sector with
K\"unneth, regular section-moduli cotangent and coefficient-one determinant;
an independent spherical code certificate; reduced derived prime contact; and a
rank-three module over the nuclear Dirichlet algebra
(\S\ref{sec:geometric-a}). &
\emph{Constructed in the intrinsic periodic--Deligne--nuclear category.}
The total cohomology is that of the geometrically induced periodic
coefficient object on the square; it is not, and by the proved growth
obstruction cannot be, the raw functor of all bounded spherical sections.
\\
\addlinespace
(b) &
The Frobenius correspondences $\Gamma_n\subset C\times C$: honest algebraic
cycles, with an intersection theory already available. &
The faithful Witt cohomological correspondences $\Gamma_n$ on
$\mathscr Y_{\mathbb S}$: constant Witt coefficient dynamics, nuclear
realization, metrized divisor line, and the derived
lambda-characteristic contact
in one constrained graph-image category (\S\ref{sec:correspondences}). &
\emph{Complete as cohomological correspondences.}  This is relative to
row~(a) in the stated category.  The contact is derived from $F_n$, agrees
in $\operatorname{Perf}_{IDN}$ with row~(a), and has determinant mass
$\Lambda(n)$; no
undecorated Chow lift is asserted. \\
\addlinespace
(c) &
$\Gamma_n\cdot\Delta=N_n$: an intersection number \emph{equals} the counting
function of the zeta function. &
The nuclear diagonal character of the completed Frobenius family equals
the two finite contact sums plus the Fourier--archimedean term
(\S\ref{sec:rowc}). &
\emph{Complete as a cohomological/nuclear Lefschetz identity.}  An ordinary
Cartier/Chow promotion is not asserted. \\
\addlinespace
(d) &
Castelnuovo--Severi: the intersection form on $NS$ has exactly one positive
eigenvalue.  An \emph{input} from geometry, which then implies the bound. &
The corner pairing has $n_+=1+\#P$, whence $n_+=1\iff\RH$ (\S3.5).
Its nuclear form has infinite rank and cannot factor through the
two-dimensional ruling quotient
(\S\ref{sec:rowdgate}).  The associated full primitive operator is proved
positive for every $0<T\le\log2$.  At every arithmetic threshold the
Gamma--Euler--Tate source, reciprocal bands, connected-cluster payment and
constant/mean-zero Schur reduction are constructed. &
\emph{Construction complete; global positivity conditional.}  The theorem
through $T=\log2$ is unconditional.  For all later thresholds the sole
remaining input is Conjecture~\ref{conj:newdLogSchurAngle}, the logarithmic
Schur-angle estimate.  Conditional on it, unit source domination, threshold
propagation and row~(d) follow.  The conjecture itself, and hence
unconditional global row~(d), are not claimed proved.
 \\
\bottomrule
\end{tabularx}
\endgroup

\medskip

Two asymmetries in this table are worth naming, because together they
locate the whole difficulty.

First, row (a) supplies a literal noncollapsed carrier, genuine divisor
lines, an intrinsic periodic coefficient object, its total derived global
sections, regular section-moduli cotangent, coefficient-one normalized
determinant and reduced derived contact in one category.  The universal
framed code supplies an independent spherical determinant certificate.
Its finiteness is local
freeness over a nuclear arithmetic algebra, not finite abelian rank; the
latter cannot realize the full bilinear kernel $\Lambda(mn)$ by
Theorem~\ref{thm:finiteNSnogo}.  In Weil's column (c) is a theorem in the
ordinary intersection theory of a surface and (d) is an input.  Over $\Z$,
row~(c) is proved below as a categorical trace theorem in the nuclear
cohomological completion of the row-(b) correspondences.  It is not promoted
back to an ordinary Cartier/Chow intersection on $\mathscr Y_{\mathbb S}$.  Row~(d)
remains an open geometric input.

Second, the preliminary Dirichlet assembly of
Theorem~\ref{thm:assembly} is one-sided: the sum runs over positive shells
$k\ge1$.  The completed row-(c) object of Theorem~\ref{thm:rowc} is instead
the summable nuclear supercharacter of the two Poisson quotients and is
intrinsically two-sided.  These
must not be confused with the failed operation of merely symmetrizing a
logarithmic derivative.  Indeed, writing
$G=\log\Gamma_{\mathbb R}$, the functional equation
$\xi(s)=\xi(1-s)$ gives, after cancelling the polar contributions,
\begin{equation}
\label{eq:collapse}
 \frac{\zeta'}{\zeta}(s)+\frac{\zeta'}{\zeta}(1-s)
 \;=\;-G'(s)-G'(1-s)
 \;=\;\log\pi-\tfrac12\psi\!\big(\tfrac s2\big)
 -\tfrac12\psi\!\big(\tfrac{1-s}2\big),
\end{equation}
an elementary function of $\Gamma$ alone.  Symmetrizing the logarithmic
derivative of $\zeta$ under $s\mapsto1-s$ therefore cancels every pole at
every zero of $\xi$: the poles of the two summands there have opposite
residues.  Everything downstream of that bare symmetrization is blind to the
zeros.  The Poisson quotient super-representation does something different:
it retains the two graded quotient representations before taking their
supercharacter.  This distinction
accounts for the negative principal-difference calculation below while permitting
Theorem~\ref{thm:rowc}; it also explains why the spectral signature cannot be
used as the geometric input required in row~(d).

\subsection{A standing constraint}
\label{ss:constraint}

Throughout, no definition may use a zero of $\xi$, a Li coefficient, the
sign of the Weil quadratic form, or a positive part extracted from that
form.  Zeros may appear in a theorem which is a consequence of a zero-free
definition; they may not appear in the definition itself.  This constraint
is what makes the difference between a construction and a restatement, and
several of the results below are negative precisely because a candidate
construction violated it.


% BEGIN INLINED SOURCE: row-a-deligne-nuclear
\section{The Deligne--nuclear spherical square}
\label{sec:geometric-a}

This section constructs the object used as row~(a).  There are two points at
which the characteristic-zero problem forces a departure from the literal
wording of the classical row.  First, the set of all bounded sections of a
spherical tensor product is much too large to have a surface
Riemann--Roch dimension.  Second, the full bilinear contact kernel
$\Lambda(mn)$ has arbitrarily large prime rank and therefore cannot factor
through a finitely generated abelian N\'eron--Severi group carrying that
kernel.  Both assertions will be proved below.  The
replacements are, respectively, intrinsic periodic coefficient cohomology
on the same square and finite local freeness over a nuclear Dirichlet
algebra.  A negabinary code supplies an independent spherical determinant
certificate for the first replacement.  These are
not hypotheses added to hide the two failures: they are constructions on the
same spherical carrier, and the two failures explain why these are the
strongest mutually compatible requirements.

\subsection{The categorical contract}

Write $\mathbb S$ for the sphere $\Gamma$-ring and $HR$ for the
Eilenberg--Mac~Lane $\Gamma$-ring of a commutative ring $R$.  We use the
closed symmetric monoidal category of Segal $\Gamma$-sets and its derived
category of simplicial $\Gamma$-sets.  The Day smash product is denoted by
$\wedge_{\mathbb S}$.  These conventions are those of
\cite[\S\S2,7]{CC-AbsoluteAlgebra}; in particular, the assembly morphism
\begin{equation}\label{eq:assembly}
 \alpha_{M,N}:M\wedge_{\mathbb S}N\longrightarrow \widetilde M\circ N
\end{equation}
is functorial in both variables.

\begin{definition}\label{def:strongrowa}
An \emph{intrinsic periodic--Deligne--nuclear arithmetic square} consists of:
\begin{enumerate}
\item a noncollapsed spherical square $\mathscr Y_{\mathbb S}$;
\item invertible divisor objects with internal principal descent and a
functorial effective cone;
\item on the canonical positive integer-radius sector, a geometrically
induced coefficient object on the square, intrinsic derived global sections,
and perfect cotangents of the regular section moduli, with one-ruling linear
growth and mixed K\"unneth quadratic growth, followed by its unique
continuous numerical extension;
\item the resulting continuous determinant line, its spherical code
certificate, the reduced derived prime-contact line and their
quotient-defined Green factor in one symmetric monoidal category;
\item an explicit finite locally free correspondence module over a nuclear
arithmetic coefficient algebra, with exact composition and von Mangoldt
contact.
\end{enumerate}
Completeness below means completeness of these five specified structures
and their comparison maps.  It does not identify coefficient cohomology
with the functor of all raw bounded spherical sections; that identification
is impossible by Theorem~\ref{thm:rawsmashnogo}.  Nor does it mean that the
two incompatible properties ruled out in
Theorems~\ref{thm:rawsmashnogo} and~\ref{thm:finiteNSnogo} have been assumed.
\end{definition}

The definition is independent of the zeros of $\xi$ and of every positivity
statement.  It is therefore eligible as input to the later rows rather than
a reformulation of their desired conclusion.

\subsection{The literal spherical carrier and its divisor lines}

Let $X=\overline{\Spec\Z}$ with the sheaf $\mathcal O_X$ of
$\mathbb S$-algebras constructed by Connes and Consani
\cite[Proposition~6.3]{CC-AbsoluteAlgebra}.  Thus $X$ has the usual finite
points and the archimedean closed point, and $\mathcal O_X$ is obtained by
imposing the ordinary absolute-value bound at infinity.  Define
\begin{equation}\label{eq:sphericalsquare}
 \mathscr Y_{\mathbb S}:=X\times_{\mathbb S}X,
 \qquad
 \mathcal O_{\mathscr Y}:=\operatorname{pr}_1^{-1}\mathcal O_X
       \wedge_{\mathbb S}\operatorname{pr}_2^{-1}\mathcal O_X.
\end{equation}
This formula is also a construction of the product: on every pair of open
charts it takes the Day smash of the two section $\Gamma$-rings, uses the
smashes of the restriction maps, and then sheafifies levelwise.  Thus the
relative smash in \eqref{eq:sphericalsquare} is, by definition, the
sheafification of that presheaf; the sheaf axiom is not being assumed for an
un-sheafified Day product.

\begin{proposition}[Noncollapse]\label{prop:sphericalnoncollapse}
The two projections in \eqref{eq:sphericalsquare} do not identify
$\mathscr Y_{\mathbb S}$ with $X$.
\end{proposition}

\begin{proof}
At the generic stalk, filtered colimits commute with the Day colimits and
the structure $\Gamma$-ring is
$H\Z\wedge_{\mathbb S}H\Z$.  Connes and Consani prove that this smash is not
isomorphic to $H\Z$ \cite[Proposition~7.4]{CC-AbsoluteAlgebra}.  More
explicitly, after assembly the two-variable Laurent polynomial class
$(t_1-1)(t_2-1)$ restricts to zero under both coordinate marginals but is
nonzero before either restriction.  For $H\Z$ the pair of marginals is
bijective.  An isomorphism of the square with either factor would induce an
isomorphism on the generic stalk and contradict that distinction.
\end{proof}

An Arakelov divisor on $X$ is written $D=D_{\mathrm f}+a\{\infty\}$.
Connes--Consani associate to it a sheaf $\mathcal O_X(D)$ of
$\mathcal O_X$-modules, and multiplication by $q^{-1}$ gives the canonical
isomorphism
\begin{equation}\label{eq:principalmodule}
 \mathcal O_X(D)\xrightarrow{\ \sim\ }
 \mathcal O_X(D+\divv(q)),\qquad q\in\Q^\times.
\end{equation}
Their multiplication theorem says that
$\mathcal O_X(D)\wedge_{\mathcal O_X}\mathcal O_X(E)\simeq
\mathcal O_X(D+E)$ whenever the relevant archimedean radii are rational;
for arbitrary radii its strict nonunital version has the same exact tensor
law \cite[Proposition~6.4 and Remark~6.5]{CC-AbsoluteAlgebra}.

We shall use the integral, prime-generated, rational-radius sector
$\Div_{\mathrm{ext}}(X)$.  It contains all finite prime rulings needed for
the contact formula.  For $D,E$ in this sector put
\begin{equation}\label{eq:externalLine}
 \mathcal L(D,E):=\operatorname{pr}_1^*\mathcal O_X(D)
 \wedge_{\mathcal O_{\mathscr Y}}
 \operatorname{pr}_2^*\mathcal O_X(E).
\end{equation}

\begin{proposition}[Divisor tensor law and principal descent]
\label{prop:divisorlines}
The object $\mathcal L(D,E)$ is invertible and there are coherent
isomorphisms
\begin{equation}\label{eq:lineTensorLaw}
 \mathcal L(D,E)\wedge\mathcal L(D',E')
 \simeq\mathcal L(D+D',E+E'),
 \qquad \mathcal L(D,E)^{-1}\simeq\mathcal L(-D,-E).
\end{equation}
The assignment descends through principal divisors.
\end{proposition}

\begin{proof}
Associativity and symmetry of relative smash reduce the first isomorphism
to the one-dimensional multiplication theorem in each factor.  Taking
$(D',E')=(-D,-E)$ gives the unit and proves invertibility.  For principal
changes $(q,r)$, smash the two explicit maps \eqref{eq:principalmodule};
their composites and tensor products are multiplication by the corresponding
products of rational numbers.  Hence the unit, associativity, symmetry and
descent diagrams commute before passage to isomorphism classes.
\end{proof}

Thus the divisor classes used below are represented by actual invertible
objects on the literal square.  No freely declared prime lattice is being
substituted for a Picard category.

For precision, let $\Pic_{\mathrm{ext}}^{\mathrm{fact}}$ be the
\emph{factorized external Picard groupoid}: an object retains the two
one-dimensional divisor modules together with their external product
\eqref{eq:externalLine}, and a morphism is a pair of module isomorphisms.
Its forgetful functor lands in the Picard groupoid of the square.  Retaining
the factorization is essential: no assertion that these objects exhaust the
full Picard groupoid of $\mathscr Y_{\mathbb S}$ is used.

\begin{proposition}[External classes and effectivity]
\label{prop:externaleffectivity}
The degree identifies the isomorphism classes of the one-dimensional
factorized divisor modules in the rational-radius sector with
$\log\Q_{>0}\subset\R$.  Consequently
\begin{equation}\label{eq:rationalPicardClasses}
 \pi_0\Pic_{\mathrm{ext}}^{\mathrm{fact}}
 \simeq(\log\Q_{>0})^2.
\end{equation}
Its degree image has closure $\R^2$; we denote this completed numerical
degree space by $N^1_{\deg}\simeq\R^2$.  The actual effective monoid maps to
$(\log\Q_{>0}\cap\R_{\geq0})^2$, whose closure is
$\R_{\geq0}^2$.  Tensor product adds the coordinates, and the code
transition maps preserve the actual effective monoid and its closure.
\end{proposition}

\begin{proof}
For $D=D_{\mathrm f}+a\{\infty\}$ in the rational-radius sector,
$a\in\log\Q_{>0}$.  Multiply by the unique positive rational
whose valuations cancel $D_{\mathrm f}$.  The result is
$\deg(D)\{\infty\}$, and its degree still belongs to
$\log\Q_{>0}$.  Conversely, if $a\{\infty\}$ and
$b\{\infty\}$ differ by a principal divisor with zero finite part, the
rational function is $\pm1$, so $a=b$.  This is also the explicit content of
the divisor-module classification
\cite[Proposition~6.4(ii)]{CC-AbsoluteAlgebra}.

For $c\{\infty\}$, the level-one global sections are
$\{n\in\Z:|n|\leq e^c\}$.  A nonzero section exists exactly when $c\geq0$.
Apply this independently to the two retained factors.  Since
$\log\Q_{>0}$ is dense in $\R$, taking closures gives the asserted numerical
space and cone.  The tensor and transition assertions follow from
\eqref{eq:lineTensorLaw} and from zero-extension of digit coordinates.
\end{proof}

\subsection{Why all bounded sections cannot be the surface cohomology}

For $m\geq0$, let
\[
 H_m=(H\Z)_{[-m,m]}
\]
be the bounded $\Gamma$-submodule: at level $k_+$ it consists of integer
$k$-tuples whose $\ell^1$ norm is at most $m$.  These are precisely the
archimedean bounded-section modules appearing after principal normalization
of an integral divisor.  Put
\[
 B_s=(1,-2,4,\ldots,(-2)^{s-1})
\]
Let $M_s$ be the largest absolute value of a subset sum of $B_s$, and put
\begin{equation}\label{eq:digitL1radius}
 L_s:=\|B_s\|_1=\sum_{j=0}^{s-1}2^j=2^s-1.
\end{equation}
For a tolerant $\mathbb S$-module, $\dim_{\mathbb S}$ denotes the least
cardinality of a level-one subset from which every level-one element is an
admissible zero--one sum, as in the dimension definition of
\cite[Appendix~A]{CC-RRSpecZ}.

\begin{lemma}[Negabinary separation and admissible length]
\label{lem:negabinaryradius}
All subset sums of $B_s$ are distinct and form an interval of $2^s$
consecutive integers.  More precisely, for $k\geq0$,
\begin{equation}\label{eq:MrExact}
 M_{2k}=\frac{2(4^k-1)}3,
 \qquad
 M_{2k+1}=\frac{4^{k+1}-1}3.
\end{equation}
The largest length for which the digit vector itself is an admissible
element of $H_m$ is
\begin{equation}\label{eq:rmdefinition}
 r(m):=\max\{s\geq0:L_s\leq m\}
      =\left\lfloor\log_2(m+1)\right\rfloor,
\end{equation}
and $r(m)=\log_2m+O(1)$ as $m\to\infty$.
\end{lemma}

\begin{proof}
Suppose that two subsets have the same sum and subtract their incidence
vectors.  If $j$ is the least index with nonzero coefficient, division by
$(-2)^j$ gives an odd number equal to an even number, a contradiction.
Thus all subset sums are distinct.  If $I_s$ is their set, then
$I_{s+1}=I_s\cup(I_s+(-2)^s)$.  Starting from $I_0=\{0\}$, the two intervals
are adjacent at every step, with their order alternating with the parity of
$s$; induction proves contiguity.  The largest positive sum is the sum of
the even powers and the absolute value of the most negative sum is the sum
of the absolute values of the odd powers.  Summing the two geometric series
gives \eqref{eq:MrExact}.  Independently, the definition of the bounded
module uses the norm of the whole digit vector, not the largest subset
sum; equation~\eqref{eq:digitL1radius} therefore gives the exact formula
\eqref{eq:rmdefinition} and its asymptotic.  Notice that $M_s\leq m$ alone
would not suffice: for example $M_4=10$ whereas $L_4=15$.
\end{proof}

\begin{theorem}[Raw spherical K\"unneth obstruction]
\label{thm:rawsmashnogo}
For $m,n\geq1$, the Connes--Consani generator dimension satisfies
\begin{equation}\label{eq:rawSmashBound}
 \dim_{\mathbb S}(H_m\wedge_{\mathbb S}H_n)
 \geq 2^{r(n)}-1,
\end{equation}
and symmetrically with $m,n$ exchanged.  If
$n_t=\lfloor e^{tb}\rfloor$ with $b>0$, the right side is
$\exp(tb+O(1))$.  Hence raw bounded sections cannot have quadratic
$O(t^2)$ surface growth.
\end{theorem}

\begin{proof}
Use $1\in H_m(1_+)$ and $B_s\in H_n(s_+)$, where $s=r(n)$.  The latter
membership is now exactly $\|B_s\|_1=L_s\leq n$.  For each
subset $J\subset\{0,\ldots,s-1\}$ the Day-colimit presentation gives a
level-one element $z_J$ by retaining exactly the pairs $(1,j)$ with
$j\in J$.  Under \eqref{eq:assembly},
\[
 \alpha(z_J)=[b_J],\qquad b_J=\sum_{j\in J}(-2)^j,
\]
in the reduced free abelian group
$\Z[\Z]/\Z[0]$.  By Lemma~\ref{lem:negabinaryradius}, the $2^s-1$ classes
with $J\neq\varnothing$ are distinct basis vectors and hence linearly
independent.

If a level-one set $F$ generates the smash by admissible zero--one sums,
each $z_J$ is such a sum of elements of $F$.  Assembly sends admissible sums
to ordinary sums in its special additive target.  The subgroup generated by
$\alpha(F)$ therefore contains $2^s-1$ independent vectors, while its rank
is at most $\#F$.  This proves \eqref{eq:rawSmashBound}.  Lemma
\ref{lem:negabinaryradius} gives
$2^{r(n_t)}=\exp(tb+O(1))$, which dominates every quadratic polynomial.
\end{proof}

The theorem fixes the logical meaning of the next construction.  We do not
rename the full Day smash.  We linearize a distinguished, functorial family
of its genuine sections and take the cotangent space of that linearization.

\subsection{The initial negabinary section code}

The Day presentation is Boolean at this point: a map to $1_+$ either retains
or discards each input pair.  We require the radix to be negative so that
one Boolean tower contains sections of both signs without an additional
sign coordinate.  Call a negative integer $b$ an
\emph{expanding Boolean radix} when $|b|>1$, and order these radices by
absolute value.  This ordered set has the unique initial element $-2$.
Equivalently, $-2$ is the unique negative integral solution of
\begin{equation}\label{eq:absoluteRelation}
 1+1=X+X^2.
\end{equation}
Thus evaluation $X\mapsto-2$ defines a morphism from the quotient of the
free spherical algebra by \eqref{eq:absoluteRelation} to $H\Z$.  The other
integral solution $X=1$ repeats the unit and gives no expanding digit
system.  Hence the initial expanding negative integral Boolean-radix rule
selects $-2$
without an auxiliary prime, finite field or interpolation base.  The
radices $-3,-4,\ldots$ are not alternative outputs of this rule because
none is initial.  This is canonicity inside the explicitly stated ordered
class of radices; no assertion that the spherical carrier itself uniquely
selects that class is used.

Let $G_r=\mathbb F_2^r$.  For $r=r(m)$ define
\begin{equation}\label{eq:oneCode}
 c_m:G_r\longrightarrow H_m(1_+),\qquad
 (\epsilon_i)\longmapsto\sum_{i=0}^{r-1}\epsilon_i(-2)^i.
\end{equation}
Every represented integer has absolute value at most
$M_r\leq L_r\leq m$, so the map lands in $H_m(1_+)$; injectivity is the
first assertion of Lemma~\ref{lem:negabinaryradius}.  Let
\begin{equation}\label{eq:oneEnvelope}
 Q_m:=\Spec\operatorname{Sym}_{\mathbb F_2}(G_r^\vee)
       \simeq\mathbb A^r_{\mathbb F_2}.
\end{equation}
Its $\mathbb F_2$-points are precisely the code parameters.  We call it the
\emph{linear code envelope}; it is not claimed to represent all bounded
spherical sections over arbitrary coefficient rings.

For $r=r(m)$ and $s=r(n)$, and for a matrix
$\epsilon=(\epsilon_{ij})\in\mathbb F_2^{r\times s}$, let
$v_\epsilon:r_+\wedge s_+\to1_+$ retain $(i,j)$ exactly when
$\epsilon_{ij}=1$.  The Day presentation defines
\begin{equation}\label{eq:mixedSection}
 z_\epsilon=[r_+,s_+,v_\epsilon,B_r\wedge B_s]
       \in(H_m\wedge_{\mathbb S}H_n)(1_+).
\end{equation}
Here $B_r\in H_m(r_+)$ and $B_s\in H_n(s_+)$ by
\eqref{eq:digitL1radius}--\eqref{eq:rmdefinition}; this is the admissibility
condition required by the Day representative.

\begin{theorem}[Separation of the mixed code]
\label{thm:mixedcodeseparation}
The map $\epsilon\mapsto z_\epsilon$ is injective.  In particular the
$2^{r(m)r(n)}$ points of
\begin{equation}\label{eq:mixedEnvelope}
 Q_{m,n}:=\Spec\operatorname{Sym}_{\mathbb F_2}
       ((G_r\otimes_{\mathbb F_2}G_s)^\vee)
\end{equation}
label distinct genuine sections of the spherical mixed module.
\end{theorem}

\begin{proof}
Put $D_m=(\log m)\{\infty\}$ and $D_n=(\log n)\{\infty\}$.  The
external-product morphism sends \eqref{eq:mixedSection} to a global
section of $\mathcal L(D_m,D_n)$.  On restriction to the generic finite chart it
is the same Day representative in $H\Z\wedge_{\mathbb S}H\Z$.  It is
therefore enough to separate the restrictions: equality of two global
sections would imply equality on that chart.  There the assembly formula is
\begin{equation}\label{eq:mixedAssemblyFormula}
 \alpha(z_\epsilon)=
 \sum_{i=0}^{r-1}(-2)^i
 \left[\sum_{j=0}^{s-1}\epsilon_{ij}(-2)^j\right]
 \quad\text{in }\Z[\Z]/\Z[0].
\end{equation}
For a nonzero exponent $b$, its coefficient is the subset sum of the
$(-2)^i$ over the rows whose inner sum equals $b$.  Negabinary uniqueness
recovers that subset of row indices.  The exponent $b$ then recovers, again
by uniqueness, the selected columns in every one of those rows.  A nonempty
row cannot have inner sum zero, since zero already represents the empty
subset.  After all nonzero exponents are recovered, the remaining row
indices are exactly the empty rows.  Hence \eqref{eq:mixedAssemblyFormula}
determines every entry of $\epsilon$.  Equality before assembly would imply
equality after assembly, so the original sections are distinct.
\end{proof}

The precise status of the affine envelope is important.  For an
$\mathbb F_2$-algebra $R$, define the framed digit-deformation functor
\begin{equation}\label{eq:digitDeformationFunctor}
 \mathfrak P_{m,n}(R):=(G_r\otimes_{\mathbb F_2}G_s)
                       \otimes_{\mathbb F_2}R.
\end{equation}
It remembers the ordered digit-pair coordinates, while
Theorem~\ref{thm:mixedcodeseparation} supplies an injection from its
$\mathbb F_2$-points to genuine spherical sections.

\begin{proposition}[Universal framed linear envelope]
\label{prop:universalCodeEnvelope}
The scheme $Q_{m,n}$ represents $\mathfrak P_{m,n}$.  For every
finite-dimensional $\mathbb F_2$-space $W$, fibrewise linear natural
transformations from $\mathfrak P_{m,n}$ to $R\mapsto W\otimes R$ are in
canonical bijection with linear maps $G_r\otimes G_s\to W$.  Under this
bijection the identity is initial among framed linear realizations of the
digit parameters.  Its cotangent at zero is canonically
$(G_r\otimes G_s)^\vee$.
\end{proposition}

\begin{proof}
For every $R$,
\[
 Q_{m,n}(R)=\operatorname{Hom}_{\mathbb F_2\text{-alg}}
 (\operatorname{Sym}((G_r\otimes G_s)^\vee),R)
 \simeq(G_r\otimes G_s)\otimes R.
\]
The transformation assertion is the Yoneda correspondence for vector
schemes.  Taking the augmentation ideal at the zero vector gives the last
assertion.
\end{proof}

This universal property is deliberately framed: it is universal after the
ordered digit observables have been selected.  It is not a moduli theorem
for all bounded sections.  Indeed $c_m$ cannot be additive:
$c_m(e_0)=1$ while $2e_0=0$ in $G_r$ and $2\cdot1\ne0$ in the integral
section group.  Nor may one replace $Q_m$ by the finite Boolean scheme
$x_i^2=x_i$ without changing the tangent theory: that scheme is \`etale and
has zero cotangent.  Thus the nonzero object below is exactly the universal
\emph{linear relaxation of the framed code}; it is not claimed to be the
cotangent of the unlinearized section functor.

After a principal normalization $D\sim a\{\infty\}$ in the actual sector,
$a\in\log\Q_{>0}$.  For an integer scale $t\geq1$, the bounded module for
$tD$ is $H_{\lfloor e^{ta}\rfloor}$.  Proposition
\ref{prop:divisorlines} transports \eqref{eq:oneCode} and
\eqref{eq:mixedSection} back to actual global sections of
$\mathcal L(tD,tE)$.  Thus the code is attached to the divisor line on the
carrier, not merely to an unrelated finite vector space.

\subsection{The framed code cotangent and its K\"unneth theorem}

Let $0$ denote the zero code point.  Since the envelopes are affine spaces,
\begin{equation}\label{eq:cotangentSpaces}
 \Omega_m:=\mathfrak m_0/\mathfrak m_0^2\simeq G_{r(m)}^\vee,
 \qquad
 \Omega_{m,n}\simeq(G_{r(m)}\otimes G_{r(n)})^\vee.
\end{equation}
They are smooth at $0$, so there is no higher cotangent obstruction.  Let
$V_m$ and $V_{m,n}$ be the real vector spaces freely generated by the
distinguished coordinate bases of $\Omega_m$ and $\Omega_{m,n}$;
this is a basis lift, not a nonexistent scalar extension
$\mathbb F_2\to\R$.

\begin{theorem}[Cotangent K\"unneth]\label{thm:cotangentKunneth}
There are canonical basis-preserving isomorphisms
\begin{equation}\label{eq:cotangentKunneth}
 V_{m,n}\simeq V_m\otimes_{\R}V_n,
 \qquad
 \dim_{\R}V_{m,n}=r(m)r(n).
\end{equation}
If $m\leq m'$ and $n\leq n'$, zero-extension in the new high digit
coordinates gives split injections $V_{m,n}\hookrightarrow V_{m',n'}$,
and these injections commute with \eqref{eq:cotangentKunneth}.
\end{theorem}

\begin{proof}
The coordinate functions at zero of $Q_{m,n}$ are indexed by Cartesian
pairs $(i,j)$.  Sending the $(i,j)$ cotangent coordinate to the tensor of the
$i$th and $j$th one-dimensional cotangent coordinates gives the first map and its
inverse on bases.  The dimension formula follows.  Monotonicity of
$r(m)$ makes zero-extension well defined; coordinate restriction is the
split cotangent projection, so duality gives the asserted tangent
injection.  All maps are defined on the same Cartesian basis, which proves
their compatibility and transitivity.
\end{proof}

For actual effective rational-radius divisors $D,E$ of positive degrees
$a,b\in\log\Q_{>0}$ and integer $t\to\infty$, set
\begin{equation}\label{eq:scaledCode}
 m_t=\lfloor e^{ta}\rfloor,\qquad n_t=\lfloor e^{tb}\rfloor,\qquad
 \mathcal H_t^{\mathrm{cot}}(D,E):=V_{m_t,n_t}.
\end{equation}

\begin{corollary}[Coefficient-one framed-code dimension]
\label{cor:cotangentRRdimension}
The limit through the positive integers
\begin{equation}\label{eq:cotangentDimension}
 h_{\mathrm{cot}}(D,E):=(\log2)^2
 \lim_{t\to\infty}\frac{\dim_{\R}\mathcal H_t^{\mathrm{cot}}(D,E)}{t^2}
\end{equation}
exists and equals $ab$.
\end{corollary}

\begin{proof}
Equation \eqref{eq:rmdefinition} gives
\[
 r(m_t)=\frac{ta}{\log2}+O(1),\qquad
 r(n_t)=\frac{tb}{\log2}+O(1).
\]
Multiply the two expansions, use Theorem~\ref{thm:cotangentKunneth}, divide
by $t^2$ and multiply by $(\log2)^2$.
\end{proof}

The normalization is forced after imposing degree matching on this code.
In one ruling the only constant
$c$ for which $c\,r(\lfloor e^{ta}\rfloor)/t\to a$ for every $a>0$ is
$c=\log2$.  K\"unneth multiplicativity then forces the product constant
$(\log2)^2$ in \eqref{eq:cotangentDimension}.  This is the precise sense in
which the coefficient one is canonical for the chosen digit code.  It is
not a metric supplied by the determinant functor before the degree-matching
condition is imposed.  No interpolation base has been chosen.  The formula,
initially proved on $(\log\Q_{>0})_{>0}^2$, extends uniquely by continuity to
the completed numerical cone $\R_{>0}^2$.

Principal invariance is equally explicit.  Multiplication by the unique
positive rational cancelling the finite coefficients sends $D$ to its
purely archimedean normalization.
The remaining rational unit is $-1$; it changes the signs of represented
sections but fixes the zero code point and every digit index.  Degree,
$r(m_t)$, the tangent lattice and all transition maps therefore depend only
on the metrized divisor class.

For a finite prime-ruling divisor
\begin{equation}\label{eq:Dpr}
 x=\sum_p x_{p,1}e_{p,1}+\sum_p x_{p,2}e_{p,2},
 \qquad
 d_i(x)=\sum_p x_{p,i}\log p,
\end{equation}
the preceding construction applied to the total divisor in each ruling
gives the following formula on the effective monoid.  Since $d_1,d_2$ are
additive, this polynomial has a unique quadratic extension to the
Grothendieck group.  For a general class we define the based line by the
same metric formula; on the effective cone it agrees with the determinant
limit already constructed:
\begin{equation}\label{eq:cotangentRRform}
 F_{\mathrm{code}}(x)=d_1(x)d_2(x),\qquad
 B_{\mathrm{code}}(x,y)
 =d_1(x)d_2(y)+d_2(x)d_1(y).
\end{equation}

\subsection{The nuclear arithmetic coefficient algebra}

Let
\begin{equation}\label{eq:DirichletAlgebra}
 \mathcal C_{\R}:=
 \left\{a=(a_n)_{n\geq1}:q_k(a):=\sum_{n\geq1}|a_n|n^k<\infty
                  \text{ for every }k\geq0\right\}
\end{equation}
with Dirichlet convolution
$(a*b)_r=\sum_{mn=r}a_mb_n$.

\begin{proposition}[Nuclear Dirichlet algebra]\label{prop:nuclearDirichlet}
$\mathcal C_{\R}$ is a complete commutative unital nuclear Fr\'echet
algebra.  Its point masses satisfy
\begin{equation}\label{eq:DirichletComposition}
 \delta_m*\delta_n=\delta_{mn}.
\end{equation}
The von Mangoldt functional
\begin{equation}\label{eq:MangoldtFunctional}
 \ell(a)=\sum_{n\geq1}a_n\Lambda(n)
\end{equation}
is continuous.
\end{proposition}

\begin{proof}
Absolute convergence and $(mn)^k=m^kn^k$ give
\[
 q_k(a*b)\leq q_k(a)q_k(b),
\]
which proves joint continuity and the algebra law.  Completeness follows by
viewing \eqref{eq:DirichletAlgebra} as the projective limit of the Banach
spaces $\ell^1(\N,n^k)$.  After conjugating the bonding map
$\ell^1(\N,n^{k+2})\to\ell^1(\N,n^k)$ to ordinary $\ell^1$, it is the
diagonal operator with entries $n^{-2}$.  Its rank-one expansion has nuclear
norm at most $\sum_nn^{-2}<\infty$.  The projective-limit criterion for
countably normed spaces therefore proves nuclearity
\cite[Chapter~4]{PietschNuclear}.  Formula
\eqref{eq:DirichletComposition} is immediate.  Finally
$\Lambda(n)\leq\log n\leq n$ for $n\geq2$, so
$|\ell(a)|\leq q_1(a)$.
\end{proof}

The augmentation
\begin{equation}\label{eq:augmentation}
 \varepsilon:\mathcal C_{\R}\longrightarrow\R,\qquad
 \varepsilon(a)=a_1,
\end{equation}
is a continuous algebra map, since $(a*b)_1=a_1b_1$.

\subsection{One arithmetic perfect category}

Let $H_{\mathrm{sol}}(\mathcal C_{\R})$ denote the Eilenberg--Mac~Lane
$E_\infty$-algebra of the solidification of the condensed Fr\'echet algebra
$\mathcal C_{\R}$.  Solidification is symmetric monoidal by
\cite[Theorem~6.2]{ClausenScholze}; hence the continuous commutative
convolution algebra of Proposition~\ref{prop:nuclearDirichlet} determines a
commutative algebra object after solidification, and its connective
Eilenberg--Mac~Lane object is an $E_\infty$-algebra.  This is the meaning of
the notation; no discrete replacement of the Fr\'echet topology is made.
On $\mathscr Y_{\mathbb S}$ put
\begin{equation}\label{eq:threeStructures}
 \mathcal O_{\Z}=\mathcal O_{\mathscr Y}\wedge_{\mathbb S}H\Z,
 \quad
 \mathcal O_{\R}=\mathcal O_{\mathscr Y}\wedge_{\mathbb S}H\R,
 \quad
 \mathcal O_{\mathcal C}=\mathcal O_{\mathscr Y}
        \wedge_{\mathbb S}H_{\mathrm{sol}}(\mathcal C_{\R}).
\end{equation}
The maps $\Z\to\R$ and \eqref{eq:augmentation} induce maps to
$\mathcal O_{\R}$.  Write $\mathscr Y_{\Z}$,
$\mathscr Y_{\R}$ and $\mathscr Y_{\mathcal C}$ for the same carrier with
these three structure sheaves.  Define
\begin{equation}\label{eq:PerfDN}
 \operatorname{Perf}_{DN}(\mathscr Y_{\mathbb S})=
 \operatorname{Perf}(\mathcal O_{\Z})
 \times^h_{\operatorname{Perf}(\mathcal O_{\R})}
 \operatorname{Perf}(\mathcal O_{\mathcal C}).
\end{equation}

Concretely, an object is a triple $(M_{\Z},M_{\mathcal C},\alpha)$ of
perfect modules and an equivalence
\[
 \alpha:M_{\Z}\otimes_{\mathcal O_{\Z}}\mathcal O_{\R}
 \xrightarrow{\sim}
 M_{\mathcal C}\otimes_{\mathcal O_{\mathcal C}}\mathcal O_{\R}.
\]
Tensor products, duals and cofibres are taken componentwise and the real
comparison maps are tensored or mapped accordingly.  This description
proves directly that \eqref{eq:PerfDN} is stable and symmetric monoidal and
that its perfect objects are dualizable.  It is the arithmetic
integral--real gluing familiar from Deligne cohomology, with the nuclear
coefficient component added.

Only the full stable symmetric monoidal subcategory generated by invertible
divisor lines, finite free modules and their finite cofibres is used below.
For these objects perfectness is local and immediate from a finite free
resolution; tensor products and descent maps are computed componentwise.
Thus no exactness assertion for an infinite completed tensor product enters
any proof.  Relative affine objects such as
\eqref{eq:relativeCodeEnvelope} mean, by definition, the opposites of the
displayed commutative free algebras; their functor of points is the
free-algebra adjunction \eqref{eq:relativeEnvelopeUniversal}.

The line \eqref{eq:externalLine} extends scalars to the three components;
the two realifications agree tautologically.  Proposition
\ref{prop:divisorlines} therefore gives invertible objects, with the same
principal descent, in \eqref{eq:PerfDN}.

At scale $t$, let $V_{t,\Z}(D,E)$ be the free integral lattice on the
ordered digit-pair coordinate basis of \eqref{eq:scaledCode}, and let
$V_t(D,E)$ be its real span.  There is an integral affine envelope
\begin{equation}\label{eq:relativeCodeEnvelope}
 \mathcal Q_{t,\Z}(D,E)=
 \underset{\mathscr Y_{\Z}}{\operatorname{Spec}}
 \operatorname{Sym}_{\mathcal O_{\Z}}
 \bigl(\mathcal O_{\Z}\otimes_{\Z}V_{t,\Z}(D,E)\bigr).
\end{equation}
Its special fibre over $\mathbb F_2$ is the product of the carrier fibre
with the affine code envelope \eqref{eq:mixedEnvelope}; in particular its
parameter fibre is exactly $Q_{m_t,n_t}$.  Define
$\mathcal Q_{t,\mathcal C}$ by the same
formula over $\mathcal O_{\mathcal C}$ with $V_t$.

The free-algebra adjunction characterizes this construction: for an affine
$f:T\to\mathscr Y_{\Z}$,
\begin{equation}\label{eq:relativeEnvelopeUniversal}
 \operatorname{Hom}_{\mathscr Y_{\Z}}(T,\mathcal Q_{t,\Z})
 \simeq
 \operatorname{Hom}_{\mathcal O_T}
 (f^*(\mathcal O_{\Z}\otimes V_{t,\Z})^\vee,\mathcal O_T).
\end{equation}
Thus it is the relative vector bundle universally linearizing the framed
digit coordinates of Proposition~\ref{prop:universalCodeEnvelope}; it is
not identified with a moduli space of all sections.

If $e_t$ is the zero section, the polynomial-algebra calculation of the
relative cotangent complex gives
\[
 e_t^*L_{\mathcal Q_{t,\Z}/\mathscr Y_{\Z}}
   \simeq\mathcal O_{\Z}\otimes_{\Z}V_{t,\Z},\qquad
 e_t^*L_{\mathcal Q_{t,\mathcal C}/\mathscr Y_{\mathcal C}}
   \simeq\mathcal O_{\mathcal C}\otimes_{\R}V_t
\]
in degree zero: a symmetric algebra is free, so its cotangent complex is its
module of differentials and has no higher term.  Define
\begin{equation}\label{eq:perfectCotangentObject}
 \mathbf H_t^{\mathrm{cot}}(D,E)=
 \bigl(
 \mathcal O_{\Z}\otimes_{\Z}V_{t,\Z}(D,E),
 \mathcal O_{\mathcal C}\otimes_{\R}V_t(D,E),
 \alpha_t
 \bigr),
\end{equation}
where $\alpha_t$ identifies the two real digit bases.  Thus
\eqref{eq:perfectCotangentObject} is the relative cotangent at the zero
section of the universal framed linear envelopes.  It is determined by the
framed code through \eqref{eq:relativeEnvelopeUniversal}, but is not the
intrinsic cotangent of the full spherical section functor.  Each component
is finite free, hence perfect, and
Theorem~\ref{thm:cotangentKunneth} gives an isomorphism in the same category
\begin{equation}\label{eq:perfectKunneth}
 \mathbf H_t^{\mathrm{cot}}(D,E)
 \simeq\mathbf H_{t,1}^{\mathrm{cot}}(D)
 \otimes\mathbf H_{t,2}^{\mathrm{cot}}(E).
\end{equation}
The zero-extension maps make these objects a functorial pro-system on the
effective divisor category, and principal isomorphisms transport that
system.

\subsection{The auxiliary normalized code determinant}

Apply the determinant functor for perfect complexes
\cite{KnudsenMumford} to
\eqref{eq:perfectCotangentObject}.  On its real digit determinant, give the
ordered wedge generator the norm
\begin{equation}\label{eq:finiteRRmetric}
 \|1_t\|=\exp\bigl(- (\log2)^2\dim_{\R}V_t(D,E)\bigr).
\end{equation}
This is the degree-matched metric of the framed code.  The determinant
functor supplies the line and its basis; the displayed norm is the explicit
normalization selected in Corollary~\ref{cor:cotangentRRdimension}, not the
orthonormal determinant metric.
For a based normed line $(L,1)$ and $c>0$, write $L^{\langle c\rangle}$ for
the same based line with norm $\|1\|^c$; this is an explicit metric
operation, not a fractional algebraic tensor power.

\begin{proposition}[Continuous normalized code determinant]
\label{prop:continuousRRdet}
The normalized lines
$\det\mathbf H_t^{\mathrm{cot}}(D,E)^{\langle t^{-2}\rangle}$ converge in
the Picard groupoid of based normed real lines to
\begin{equation}\label{eq:RRline}
 \lambda_{\mathrm{code}}(D,E)
   =\bigl(\R\mathbf1,\ \|\mathbf1\|=e^{-ab}\bigr).
\end{equation}
The limit is principal-invariant, functorial under effective inclusions and
compatible with K\"unneth.
\end{proposition}

\begin{proof}
By the proof of Corollary~\ref{cor:cotangentRRdimension},
\[
 \frac{(\log2)^2}{t^2}\dim_{\R}V_t(D,E)=ab+O(t^{-1}).
\]
Taking minus the logarithm of \eqref{eq:finiteRRmetric} proves convergence
to \eqref{eq:RRline}.  The transition, principal and K\"unneth maps preserve
the ordered digit bases; the corresponding determinant diagrams therefore
commute at every finite stage and after passage to the limit.
\end{proof}


% BEGIN INLINED SOURCE: row-a-intrinsic-periodic
\subsection{Intrinsic periodic section cohomology}
\label{subsec:intrinsicperiodic}

The code cotangent above is a spherical certificate, but it cannot be the
cotangent of the functor of all bounded spherical sections: Theorem
\ref{thm:rawsmashnogo} gives incompatible growth.  We now construct the
cohomology object used in the final package from an actual moduli space of
sections.

Connes and Consani construct the arithmetic pullback
\[
 \pi:X_{\rm vis}\longrightarrow X
\]
of the Riemann sector along the extended Abel--Jacobi map.  Its fibre over a
finite prime $p$ is canonically the pointed periodic curve
\begin{equation}\label{eq:periodicfiber}
 C_p=\R_+^\times/p^\Z\simeq\R/(\log p)\Z,
\end{equation}
whose neutral point is the standard subgroup $\Z[1/p]\subset\R$
\cite[\S6.3]{CCJacobian}.  Form the finite-prime pair pullback
\begin{equation}\label{eq:periodicpairpullback}
 H_{\rm per}:=\coprod_{p,q}C_p\times C_q,
 \qquad
 \Pi_{\rm per}:H_{\rm per}\longrightarrow\mathscr Y_{\mathbb S}.
\end{equation}
Here the target of the underlying map is $X\times X$; its spherical
structure remains the one in \eqref{eq:sphericalsquare}.  Give the source
the disjoint-union topology of the periodic sites.  Inverse images preserve
covers, so \eqref{eq:periodicpairpullback} defines a geometric morphism of
the corresponding sites.

Write $o_p$ for the neutral point and
$D_p(a):=a[o_p]$ for $a\in\Z$.  This is an actual divisor, not merely a
degree label.  The full divisor-class classification is
$(\deg,\chi):\Div(C_p)/\Prin(C_p)\simeq
\R\times\Z/(p-1)\Z$ \cite[Theorem~5.6]{CC-ScalingSite}; below we need only
the integral neutral-point sector forced by the prime rulings.

We use an $\infty$-categorical enhancement, not an unspecified model
structure.  For a fixed $(p,q)$ let $\mathsf V^0_{p,q}$ be the essentially small
symmetric monoidal category generated by the actual pointed filtered
complete $\R_{\max}$-modules $H^0(C_p,D)$, $H^0(C_q,E)$, their reduced
external products and their principal isomorphisms.  Put
$\mathsf V_{p,q}:=\operatorname{Ind}(N\mathsf V^0_{p,q})$, where $N$ is the
categorical nerve.  Day extension of
tensor makes this a presentable symmetric monoidal category whose tensor
preserves colimits separately and is therefore closed; the original section modules embed fully
faithfully.  Equivalences in $\mathsf V_{p,q}$ are equivalences in this
presentable $\infty$-category; no additional class of weak equivalences is
chosen.

Define a $\mathsf V_{p,q}$-enriched category $\mathsf{Sec}_{p,q}$.  Its
objects are pairs $(D,E)$ in the principal-equivalence orbits of
$(D_p(a),D_q(b))$, for $(a,b)\in\Z^2$.  When both degree differences are
nonnegative, put
\begin{align}\label{eq:periodicEnrichedHom}
 \underline{\operatorname{Hom}}((D,E),(D',E'))
 &:={H^0(C_p,D'-D)}\widehat\boxtimes_{\rm red}{H^0(C_q,E'-E)}.
\end{align}
If either degree difference is negative, use the zero pointed module.
Composition is actual pointwise tropical multiplication of sections.  The
divisor tensor law proves associativity, the constant section is the unit,
and addition of divisors makes $\mathsf{Sec}_{p,q}$ symmetric monoidal.

Let
\begin{equation}\label{eq:periodicCoefficientCategory}
 \mathsf A_{p,q}:=\operatorname{PSh}_{\mathsf V_{p,q}}
 (\mathsf{Sec}_{p,q}^{\rm op},\mathsf V_{p,q})
\end{equation}
denote the $\mathsf V_{p,q}$-enriched presheaf $\infty$-category, with
enriched Day convolution, and write $y_{p,q}$ for enriched Yoneda.
Equivalences and colimits in $\mathsf A_{p,q}$ are computed objectwise.
Presentability, the objectwise formulas and generation by representables are
the enriched presheaf results of
\cite[Theorem~1.2 and Corollaries~1.4--1.5]{BermanEnriched}.
Principal multiplication gives canonical enriched equivalences between
representatives.  Put $\mathcal L_{p,q}(D,E):=y_{p,q}(D,E)$.  Then
\begin{align}\label{eq:periodicYoneda}
 \mathcal L_{p,q}(x)\otimes_{\rm Day}\mathcal L_{p,q}(y)
   &\simeq\mathcal L_{p,q}(x+y),\\
 \mathbf R\underline{\operatorname{Hom}}_{\mathsf A_{p,q}}
 (y_{p,q}(0),\mathcal L_{p,q}(x))
   &\simeq\underline{\operatorname{Hom}}_{\mathsf{Sec}_{p,q}}(0,x).
\end{align}
Here $\mathbf R\underline{\operatorname{Hom}}$ means the intrinsic enriched
mapping object of the displayed $\infty$-category; it is not a derived
functor computed in an undeclared model.  The second identity is the
$\infty$-categorical enriched Yoneda equivalence.  Representables are
compact projective because mapping from $y_{p,q}(x)$ is objectwise
evaluation, which preserves colimits.  Thus there is no replacement or
higher correction hidden in \eqref{eq:periodicYoneda}.

\begin{proposition}[Universal coefficient internalization]
\label{prop:periodicUniversalInternalization}
The category $\mathsf A_{p,q}$ is the free cocomplete
$\mathsf V_{p,q}$-linear coefficient category generated by the actual
periodic divisor-section category.  More precisely, every enriched strong
symmetric monoidal functor from $\mathsf{Sec}_{p,q}$ to a cocomplete
$\mathsf V_{p,q}$-linear category extends, uniquely up to a contractible
choice, to a cocontinuous lax symmetric monoidal functor from
$\mathsf A_{p,q}$.
\end{proposition}

\begin{proof}
Enriched left Kan extension along Yoneda gives the extension.  Every
presheaf is a weighted colimit of representables, so cocontinuity forces it
uniquely.  Day convolution is the left Kan extension of tensor on
$\mathsf{Sec}_{p,q}$, proving the monoidal assertion and its coherence.
\end{proof}

To place this coefficient category over the periodic carrier without a
type error, define
\begin{equation}\label{eq:periodicCoefficientTopos}
 \mathscr E_{p,q}:=\operatorname{Sh}(C_p\times C_q;\mathrm{Spc})
 \widehat\boxtimes\mathsf A_{p,q},\quad
 \mathcal O_{p,q}:=\mathbf1\boxtimes y_{p,q}(0),\quad
 \widetilde{\mathcal L}_{p,q}(D,E):=
 \mathbf1\boxtimes\mathcal L_{p,q}(D,E).
\end{equation}
The cocomplete tensor product exists by
\cite[Proposition~4.8.1.17]{LurieHA}.  The spatial factor is cartesian
monoidal, so its unit $\mathbf1$ is terminal and
$\mathbf R\!\operatorname{Hom}(\mathbf1,\mathbf1)$ is contractible; all
tropical tensor data remain in $\mathsf A_{p,q}$.  Define total coefficient
cohomology by
\begin{equation}\label{eq:periodicTotalCohomologyDefinition}
 \mathbf R\Gamma_{\rm per}(C_p\times C_q,M):=
 \mathbf R\underline{\operatorname{Hom}}_{\mathscr E_{p,q}}
 (\mathcal O_{p,q},M).
\end{equation}
For representable coefficient lines, the terminal sheaf factor and
\eqref{eq:periodicYoneda} identify this object with the actual global
periodic section module.  This definition makes no unproved
claim that a separate tropical sheaf cohomology has vanishing higher terms.

Write $\Pi_{p,q}:C_p\times C_q\to\mathscr Y_{\mathbb S}$ for the
$(p,q)$ component of \eqref{eq:periodicpairpullback} and set
\[
 \mathscr E_{\mathscr Y;p,q}:=
 \operatorname{Sh}(\mathscr Y_{\mathbb S};\mathrm{Spc})
 \widehat\boxtimes\mathsf A_{p,q}.
\]
The coefficient adjunction is the explicit tensor product
\begin{equation}\label{eq:periodicComponentAdjunction}
 \widetilde\Pi_{p,q}^{,*}:=Pi_{p,q}^*\widehat\boxtimes
 \operatorname{id}_{\mathsf A_{p,q}} :
 \mathscr E_{\mathscr Y;p,q}
 \
ightleftarrows\
 \mathscr E_{p,q} :
 \mathbf R\widetilde\Pi_{p,q,*}:=
 \mathbf R\Pi_{p,q,*}\widehat\boxtimes
 \operatorname{id}_{\mathsf A_{p,q}}.
\end{equation}
The left arrow is strong symmetric monoidal: inverse image for a geometric
morphism preserves finite products and all colimits, hence is strong for the
cartesian spatial tensor, while the identity is strong for Day convolution.
Its right adjoint exists by presentability and is therefore canonically lax
symmetric monoidal.

Let $\mathscr E_{\rm per}^{\rm fs}$ and
$\mathscr E_{\mathscr Y;\rm per}^{\rm fs}$ denote the finite-support direct
sums of the source and target categories in
\eqref{eq:periodicComponentAdjunction}.  Taking the componentwise adjunctions
gives the promised diagram
\begin{equation}\label{eq:periodicCoefficientProjection}
 \widetilde\Pi_{\rm per}^{,*}:
 \mathscr E_{\mathscr Y;\rm per}^{\rm fs}
 \ \rightleftarrows\
 \mathscr E_{\rm per}^{\rm fs}:
 \mathbf R\widetilde\Pi_{{\rm per},*},
 \qquad
 \widetilde\Pi_{\rm per}^{,*}\dashv
 \mathbf R\widetilde\Pi_{{\rm per},*}.
\end{equation}
This adds coefficient geometry over the fixed spherical carrier and does
not replace $\mathcal O_{\mathscr Y}$.  Total sections on the target are
derived enriched Hom from the componentwise unit
$\mathbf1_{\mathscr Y}\boxtimes y_{p,q}(0)$.  Adjunction identifies them
with \eqref{eq:periodicTotalCohomologyDefinition}; hence every source,
target and monoidal structure used below is explicit.

Let $\Pic_{\mathbb N}^{+}\subset\Pic_{\rm ext}^{\rm fact}$ be the monoid of
positive factorized classes whose two radii
$Q_i=\exp(d_i)$ are positive integers.  A class in this monoid has the
canonical, representative-independent coordinates
\begin{equation}\label{eq:canonicalPrimeCoordinates}
 a_{p,i}=v_p(Q_i),\qquad
 d_i=\sum_p a_{p,i}\log p,\qquad a_{p,i}\geq0.
\end{equation}
Indeed degree is principal-invariant and unique factorization recovers all
$a_{p,i}$ from it.  Thus a principal change of an Arakelov divisor cannot
change the periodic components used below.  This is the precise input
sector of the periodic cohomology; the resulting quadratic determinant is
then extended to the completed numerical degree cone by continuity.

For a finite-support effective ruling class in this monoid, written in its
canonical coordinates as
\[
 x=\sum_p a_{p,1}e_{p,1}+\sum_qb_{q,2}e_{q,2},
\]
form the finite coproduct of the internal representable lines and define
\begin{equation}\label{eq:intrinsicCoefficientObject}
 \mathcal F_{\rm per}(x)
 :=\mathbf R\widetilde\Pi_{{\rm per},*}
    \bigoplus_{p,q}'\widetilde{\mathcal L}_{p,q}
          (D_p(a_{p,1}),D_q(b_{q,2})).
\end{equation}
The prime means that only the support of $x$ is used.  Its spatial factor is
the unit sheaf, but its coefficient factor is not a precomputed numerical
vector space: the representable line is first constructed from the actual
periodic section category and only then pushed to the square.

\begin{theorem}[Internal periodic sections and K\"unneth]
\label{thm:intrinsicPeriodicKunneth}
There are natural equivalences
\begin{align}\label{eq:intrinsicPeriodicSections}
 \mathbf R\Gamma(\mathscr Y_{\mathbb S},\mathcal F_{\rm per}(x))
 &\simeq
 \bigoplus_{p,q}'
 H^0(C_p,D_p(a_{p,1}))
 \widehat\boxtimes_{\rm red}
 H^0(C_q,D_q(b_{q,2})).
\end{align}
They are principal-invariant, lax symmetric monoidal in $x$, and the
componentwise equivalence is the internal K\"unneth map.
\end{theorem}

No stabilization is used in this formula: both sides are objects of the
presentable coefficient $\infty$-category.  Stabilization begins only after
Theorem~\ref{thm:regularSectionCotangent}: its finite-dimensional real
cotangent frame $V$ is sent to the perfect complex $V[0]\in\operatorname{Perf}(\R)$,
and then to its integral and nuclear realizations.  Thus the degree-zero
perfect complex used later is defined explicitly and is not being inferred
from an unproved acyclicity assertion about the raw spherical structure
sheaf.

\begin{proof}
Adjunction for \eqref{eq:periodicCoefficientProjection}, with global
sections defined by $\mathbf R\!\operatorname{Hom}$ from the coefficient
unit, identifies the left side with
\eqref{eq:periodicTotalCohomologyDefinition} on $\mathscr E_{\rm per}^{\rm fs}$.
Its components are
open and closed, so finite-support global sections split componentwise.
Apply the second equivalence of \eqref{eq:periodicYoneda}.  By
\eqref{eq:canonicalPrimeCoordinates}, principal translation on the
Arakelov class leaves the prime coordinates unchanged; principal
translation inside a periodic component is an isomorphism of its divisor
line.  Day convolution
gives multiplication before direct image, and right direct image is lax
symmetric monoidal.  These constructions prove all compatibility diagrams.
\end{proof}

We next exhibit the intrinsic cotangent of the section object.  For $N>p$,
Connes--Consani's finite-depth module $\mathcal E_{N,p}$ has exactly
$d=N-p+1$ ordered extremal rays $\phi_0,\ldots,\phi_{d-1}$, these rays
generate the whole module, and their canonical coefficient retraction is
given by residuation
\begin{equation}\label{eq:periodicResiduation}
 \gamma_i(f)=\inf_z\bigl(f(z)-\phi_i(z)\bigr)
\end{equation}
\cite[Lemma~5.19, Appendix~A, Proposition~A.2 and Lemma~A.3]
{CC-ScalingSite}.  For the second prime use extremals
$\psi_0,\ldots,\psi_{e-1}$.  The functionally reduced external section
space is
\begin{equation}\label{eq:actualMixedSectionModuli}
 \mathfrak S_{d,e}
 =\left\{
  \max_{i,j}\{\phi_i(u)+\psi_j(v)+c_{ij}\}:c_{ij}\in\R_{\max}
  \right\}
 \subset C(C_p\times C_q,\R_{\max})
\end{equation}
with the uniform topology.

Call a section in \eqref{eq:actualMixedSectionModuli} \emph{regular} if
every pair $(i,j)$ is the unique maximizer at some point.  This is
intrinsic: it says that deleting the corresponding extremal ray changes
the section on an open set.  Denote the regular locus by
$\mathfrak S_{d,e}^{\rm reg}$.

\begin{theorem}[Regular section-moduli cotangent]
\label{thm:regularSectionCotangent}
The locus $\mathfrak S_{d,e}^{\rm reg}$ is nonempty and open.  It is a real
manifold of dimension $de$, and its cotangent bundle has the canonical
ordered frame
\begin{equation}\label{eq:intrinsicPeriodicCotangent}
 \Omega_{d,e}^{\rm per}
 \simeq\mathcal O_{\mathfrak S^{\rm reg}}
   \otimes_\R\R[\operatorname{Ext}(\mathcal E_{N,p})
                 \times\operatorname{Ext}(\mathcal E_{M,q})]^\vee.
\end{equation}
External product induces the cotangent K\"unneth isomorphism.  Principal
translations act on \eqref{eq:intrinsicPeriodicCotangent} by the identity.
\end{theorem}

\begin{proof}
Lemma~5.19(iii) of \cite{CC-ScalingSite} supplies offsets and points where
each one-dimensional extremal uniquely dominates.  Their Cartesian
products give a coefficient vector in \eqref{eq:actualMixedSectionModuli}
and a unique-dominance witness for every pair, proving nonemptiness.

For a regular coefficient vector, choose one witness for every pair.  The
minimum of the finitely many strict gaps is positive.  Perturbing every
coefficient by less than one quarter of that gap preserves all witnesses,
and evaluation there recovers each $c_{ij}$ continuously.  Conversely the
maximum map is $1$-Lipschitz in the coefficient sup norm.

The coefficients are intrinsic.  Applied to the pair extremal
$B_{ij}=\phi_i\boxplus\psi_j$, the residuation formula is
\[
 c_{ij}=\inf_{(u,v)}(f(u,v)-B_{ij}(u,v)).
\]
The inequality $f\geq B_{ij}+c_{ij}$ gives one direction, and equality at
an activity point gives the other.  The infimum is $1$-Lipschitz in $f$;
hence sufficiently small uniform perturbations preserve the positive gaps
and regularity.  The coefficient maps are therefore compatible manifold
charts and their differentials give \eqref{eq:intrinsicPeriodicCotangent}.
The extremals of an external tensor are exactly the Cartesian pairs:
unique-dominance witnesses prove that every pair is extremal, while the
pair functions generate the tensor image.  This proves K\"unneth.  A
principal translation adds a fixed function and leaves every coefficient
differential unchanged.
\end{proof}

For $a\in\Z_{>0}$ let $d_{p,r}(a)$ be the number of extremal rays of
$H^0(C_p,D_p(a))^{p^r}$.  Frobenius and the exact special-module theorem
give, for $ap^r>p$,
\begin{equation}\label{eq:periodicOneDimensionalLimit}
 d_{p,r}(a)=ap^r-p+1,
 \qquad p^{-r}d_{p,r}(a)\longrightarrow a
\end{equation}
\cite[Lemmas~5.18--5.19]{CC-ScalingSite}.  This is an exact model of the
actual filtered section module because $D_p(a)$ is the integral divisor at
the neutral point.  Explicitly, Lemma~5.18 identifies the $r$-fold
Frobenius pullback with the divisor of degree $ap^r$, and Lemma~5.19
identifies its special module with $\mathcal E_{ap^r,p}$, whose ordered
extremal count is $ap^r-p+1$.  Thus no asymptotic estimate is substituted
for the displayed finite-rank equality.

The logarithm identifies the quotient map
$\R_+^\times\to C_p$ with $\R\to\R/(\log p)\Z$.  Consequently there is a
unique translation-invariant measure $\mu_p$ whose pullback normalization
on the universal cover is Lebesgue measure, and
\begin{equation}\label{eq:canonicalPeriodicHaarLength}
 \mu_p(C_p)=\log p.
\end{equation}
This normalization is forced by the fixed logarithmic coordinate: replacing
$\mu_p$ by a scalar multiple would no longer pull back to Lebesgue measure.
In particular, the factors below are geometric fibre lengths, not adjustable
weights.  Define the arithmetic
continuous dimension by weighting the componentwise continuous dimensions
with the product of the two canonical fibre lengths.  Then, for effective
$x$,
\begin{equation}\label{eq:intrinsicContinuousDimension}
 \operatorname{cdim}_{\rm ar}\mathbf R\Gamma
    (\mathscr Y_{\mathbb S},\mathcal F_{\rm per}(x))
 =\sum_{p,q}(\log p)(\log q)a_{p,1}b_{q,2}
 =d_1(x)d_2(x).
\end{equation}
This is the dimension of internal global sections of
\eqref{eq:intrinsicCoefficientObject}, not the dimension of the raw Day
smash.

The claimed numerical extension is unique in the asymptotic degree cone,
not by a false density assertion about $\log\mathbb N$ itself.  For every
$a>0$ put $Q_t=\lfloor e^{ta}\rfloor$.  Then
$t^{-1}\log Q_t\to a$.  Applying this independently in the two rulings
shows that the normalized values on integer-radius classes satisfy
\[
 t^{-2}(\log Q_t)(\log R_t)\longrightarrow ab.
\]
Hence every point of $\R_{\geq0}^2$ is an asymptotic ray limit, and
continuity together with bidegree-one homogeneity forces the unique
extension $F(a,b)=ab$.

\subsection{The intrinsic periodic--Deligne--nuclear category}

Let $\operatorname{Perf}_{\rm per}$ be the stable idempotent completion of
the split exact category generated by the finite-depth cotangent frames
\eqref{eq:intrinsicPeriodicCotangent}, principal translations and K\"unneth
maps.  Its tensor product is induced by external product; Cartesian
products of extremal bases make it exact and symmetric monoidal.  Extend
scalars by
\begin{equation}\label{eq:periodicNuclearCategory}
 \operatorname{Perf}_{\rm per,\mathcal C}
 :=\operatorname{Perf}_{\rm per}
   \otimes_{\operatorname{Perf}(\R)}
   \operatorname{Perf}(\mathcal O_{\mathcal C}).
\end{equation}
The ordered-frame realization is the exact strong symmetric monoidal
functor
\begin{equation}\label{eq:periodicFrameFunctor}
 T_{\mathcal C}:\operatorname{Perf}_{\rm per,\mathcal C}
 \longrightarrow\operatorname{Perf}(\mathcal O_{\mathcal C}),
 \qquad
 P\longmapsto\mathcal O_{\mathcal C}\otimes_\R V_{\rm Ext}(P).
\end{equation}
Coefficient-only objects are obtained by tensoring the periodic unit with
an arbitrary perfect $\mathcal O_{\mathcal C}$-module.

Define
\begin{equation}\label{eq:PerfIDN}
 \operatorname{Perf}_{IDN}(\mathscr Y_{\mathbb S})
 :=\operatorname{Perf}_{DN}(\mathscr Y_{\mathbb S})
 \times^h_{\operatorname{Perf}(\mathcal O_{\mathcal C})}
 \operatorname{Perf}_{\rm per,\mathcal C},
\end{equation}
using nuclear-component evaluation on the first factor and
\eqref{eq:periodicFrameFunctor} on the second.

\begin{proposition}[Categorical gluing]
\label{prop:intrinsicCategoricalGluing}
The category \eqref{eq:PerfIDN} is stable and symmetric monoidal.  An object
is a Deligne--nuclear perfect object, a periodic perfect object, and an
equivalence identifying the nuclear component with its periodic frame
realization.  The coefficient-only construction gives compatible lifts of
the perfect divisor, determinant and contact objects used in rows
\textup{(b)} and \textup{(c)}; the reduced contact has zero periodic
component.  The completed nuclear correspondence module is not asserted to
be an object of this perfect pullback: it acts through the common
$\mathcal O_{\mathcal C}$ coefficient component.
\end{proposition}

\begin{proof}
Small stable idempotent-complete categories admit their tensor product and
stable completion.  Exact strong symmetric monoidal functors preserve
finite cofibres and tensor products.  Homotopy limits of stable categories
are stable; the componentwise tensor and duals preserve the comparison
equivalence.  This proves the first assertion.  The concrete description
is the definition of a homotopy pullback.  For a Deligne--nuclear perfect object
with nuclear component $M_{\mathcal C}$, use the periodic unit tensored
with $M_{\mathcal C}$; its image under $T_{\mathcal C}$ is
$M_{\mathcal C}$.  For a contact with zero nuclear component use the zero
periodic object.  These constructions apply to the named perfect objects.
An arbitrary complete nuclear module need not be dualizable; nevertheless
its scalar action on the common nuclear component is defined without
promoting it to a perfect object.
\end{proof}

At depth $(r,s)$, let $V_{p,q;r,s}^{\rm per}(x)$ be the real vector space
on the ordered extremal-pair frame and let $V_{p,q;r,s,\Z}^{\rm per}(x)$ be
its free integral lattice.  Together with its nuclear scalar extension and
the intrinsic cotangent object, these define
\begin{equation}\label{eq:intrinsicPerfectCohomology}
 \mathbf H_{p,q;r,s}^{\rm int}(x)
 \in\operatorname{Perf}_{IDN}(\mathscr Y_{\mathbb S}).
\end{equation}
Their finite-support direct sums give $\mathbf H_{r,s}^{\rm int}(x)$.
Theorem~\ref{thm:regularSectionCotangent} proves K\"unneth and principal
descent for these objects.

\begin{proposition}[Arithmetic scalar compatibility]
\label{prop:intrinsicArithmeticAction}
The intrinsic perfect systems carry endomorphisms $\rho_n^{\rm int}$ with
$\rho_m^{\rm int}\rho_n^{\rm int}=\rho_{mn}^{\rm int}$.  For $n>1$ the
integral and real components are zero, while the nuclear and periodic
scalar components are left convolution by $\delta_n$; for $n=1$ every
component is the identity.  These endomorphisms commute with principal
descent, finite-support sums and K\"unneth.
\end{proposition}

\begin{proof}
The augmentation \eqref{eq:augmentation} sends $\delta_n$ to zero for
$n>1$ and $\delta_1$ to one.  Hence the integral map just specified and
left convolution on the nuclear component have identical realifications,
so they define a morphism in $\operatorname{Perf}_{DN}$.  On the periodic
factor use the coefficient-only scalar map
$\operatorname{id}_P\otimes L_{\delta_n}$; its image under
$T_{\mathcal C}$ is the same nuclear convolution map, so the homotopy-pullback
comparison square commutes.  Finally
$L_{\delta_m}L_{\delta_n}=L_{\delta_m*\delta_n}=L_{\delta_{mn}}$.
All structural maps named in the statement are
$\mathcal O_{\mathcal C}$-linear and based, proving the compatibilities.
\end{proof}

For the comparison with the code, let $\operatorname{Sys}_{p,q}$ be the
stable category of bi-indexed sequences of based finite perfect real
objects, with pointwise exact triangles.  Let
$\operatorname{Sys}_{p,q}^{<2}$ be the thick subcategory of systems whose
total rank is $o(p^rq^s)$.  Shifts, retracts and finite cofibres preserve
this estimate, so it is thick.  We write
\begin{equation}\label{eq:subquadraticVerdierQuotient}
 \operatorname{Sys}_{p,q}^{(2)}
 :=\operatorname{Sys}_{p,q}/\operatorname{Sys}_{p,q}^{<2}
\end{equation}
for the resulting Verdier quotient.  Thus ``modulo subquadratic perfect
systems'' below refers to the explicitly defined category
\eqref{eq:subquadraticVerdierQuotient}.

Give the ordered extremal determinant generator on the $(p,q)$ component
the norm
\begin{equation}\label{eq:intrinsicFiniteMetric}
 \|\omega_{p,q;r,s}\|
 :=\exp\!\left(-(\log p)(\log q)p^{-r}q^{-s}
    \dim_\R V_{p,q;r,s}^{\rm per}(x)\right).
\end{equation}
Equations \eqref{eq:periodicOneDimensionalLimit} and
\eqref{eq:intrinsicContinuousDimension} prove convergence of the
finite-support determinant product to
\begin{equation}\label{eq:intrinsicRRline}
 \lambda_{\rm int}(x)
 =\bigl(\R\mathbf1,\ \|\mathbf1\|=e^{-d_1(x)d_2(x)}\bigr).
\end{equation}

For the system-level comparison define the decomposed spherical certificate
\begin{equation}\label{eq:decomposedCodeCertificate}
 \lambda_{\rm code}^{\rm dec}(x):=
 \bigotimes_{p,q}'
 \lambda_{\rm code}(a_{p,1}\log p,b_{q,2}\log q).
\end{equation}
The bilinear formula \eqref{eq:cotangentRRform} and distributivity give the
canonical based isometry
\begin{equation}\label{eq:decomposedTotalCodeIsometry}
 \lambda_{\rm code}^{\rm dec}(x)\simeq\lambda_{\rm code}(x);
\end{equation}
both distinguished generators have norm
$\exp(-\sum_{p,q}a_{p,1}b_{q,2}\log p\log q)$.

\begin{theorem}[Intrinsic determinant and spherical certificate]
\label{thm:intrinsicCodeComparison}
The line \eqref{eq:intrinsicRRline} is the continuous determinant of the
cotangent of the regular moduli of the intrinsic sections
\eqref{eq:intrinsicPeriodicSections}.  There is a canonical based isometry
\begin{equation}\label{eq:intrinsicCodeIsometry}
 \lambda_{\rm int}(x)\xrightarrow{\ \sim\ }\lambda_{\rm code}(x).
\end{equation}
For each fixed divisor and prime pair, the periodic cotangent system and the
corresponding component of the decomposed code certificate are isomorphic
in $\operatorname{Sys}_{p,q}^{(2)}$.
\end{theorem}

\begin{proof}
The first assertion is Theorem~\ref{thm:regularSectionCotangent} followed by
\eqref{eq:intrinsicFiniteMetric}.  On a $(p,q)$ component write
$a=a_{p,1}$ and $b=b_{q,2}$.  For the comparison choose the cofinal integer
code scales
\[
 t_{p,r}=\left\lfloor\frac{p^r\log2}{\log p}\right\rfloor,
 \qquad
 t_{q,s}=\left\lfloor\frac{q^s\log2}{\log q}\right\rfloor.
\]
By Lemma~\ref{lem:negabinaryradius}, the one-dimensional code ranks satisfy
\[
 \frac{r(\lfloor e^{t_{p,r}a\log p}\rfloor)}{p^r}\longrightarrow a,
\]
and likewise for $(q,s,b)$.  Their difference from the periodic ranks in
\eqref{eq:periodicOneDimensionalLimit} is therefore $o(p^r)$ and $o(q^s)$.
The difference of the mixed ranks is $o(p^rq^s)$.

Order both pair bases lexicographically and include each into the based
free object of the larger rank.  The two cokernels have subquadratic total
rank, so the resulting zigzag is an isomorphism in the stated stable
quotient.  Moreover
\[
 \frac{(\log2)^2}{t_{p,r}t_{q,s}}
 \big/\bigl((\log p)(\log q)p^{-r}q^{-s}\bigr)
 \longrightarrow1.
\]
The complementary determinant exponent tends to zero, and both determinant
norms tend to $e^{-ab\log p\log q}$.  Sending the distinguished limiting
generator to the distinguished limiting generator proves
the componentwise isometry; tensor over the finite support of $x$ and then
compose with \eqref{eq:decomposedTotalCodeIsometry} to obtain
\eqref{eq:intrinsicCodeIsometry}.
\end{proof}

Define
\begin{equation}\label{eq:intrinsicRRform}
 F_{\rm int}(x)=d_1(x)d_2(x),\qquad
 B_{\rm int}(x,y)=d_1(x)d_2(y)+d_2(x)d_1(y).
\end{equation}
The isometry \eqref{eq:intrinsicCodeIsometry} identifies this form with the
one computed by the auxiliary code, but its source is now intrinsic section
cohomology.

% END INLINED SOURCE: row-a-intrinsic-periodic


\subsection{Derived contact and the intrinsic Green factor}

We now derive the finite contact from the two prime rulings.  The canonical
section of $\mathcal O_X([p])$ pulls back to maps
\begin{equation}\label{eq:primeRulingSections}
 s_{p,1}:\mathcal L(-[p],0)\longrightarrow\mathcal O_{\mathscr Y},
 \qquad
 s_{p,2}:\mathcal L(0,-[p])\longrightarrow\mathcal O_{\mathscr Y}.
\end{equation}
After extension to $\mathcal O_{\Z}$, define the two derived zero-locus
objects by their cofibres
$\mathcal O_{V_{p,i}}:=\operatorname{cofib}(s_{p,i})$ and their derived
incidence by
\begin{equation}\label{eq:derivedPrimeIncidence}
 \mathcal I_p^{\mathbf L}:=
 \mathcal O_{V_{p,1}}\otimes_{\mathcal O_{\Z}}^{\mathbf L}
 \mathcal O_{V_{p,2}}.
\end{equation}
Let $\Delta:X\to\mathscr Y_{\mathbb S}$ be the diagonal.  Since
$\Delta^*\mathcal L(-[p],0)\simeq\mathcal O_X(-[p])$ and likewise in the
second factor, both pulled-back sections are multiplication by $p$ in an
integral local trivialization.

\begin{proposition}[Derived prime contact]
\label{prop:derivedPrimeContact}
The diagonal pullback of \eqref{eq:derivedPrimeIncidence} has the canonical
Koszul model
\begin{equation}\label{eq:fullDerivedPrimeContact}
 \mathbf L\Delta^*\mathcal I_p^{\mathbf L}
 \simeq
 (\mathcal O_X/p)\otimes_{\mathcal O_X}^{\mathbf L}
 (\mathcal O_X/p).
\end{equation}
Its degree-zero reduced contact is $\mathcal O_X/p$ and its integral
coefficient is $\mathbb F_p$.  The latter has the canonical perfect
$\Z$-resolution $[\Z\xrightarrow p\Z]$.  The degree $-1$ term of the full
contact is another copy of $\mathbb F_p$; it is retained in
\eqref{eq:fullDerivedPrimeContact} and is removed only by the explicitly
named reduced truncation.
\end{proposition}

\begin{proof}
Derived pullback is symmetric monoidal and takes the two maps
\eqref{eq:primeRulingSections} to the same Cartier map $p$.  Taking their
cofibres proves \eqref{eq:fullDerivedPrimeContact}.  Resolve either quotient
by $[\mathcal O_X\xrightarrow p\mathcal O_X]$.  Tensoring with
$\mathcal O_X/p$ makes the differential zero, so the derived tensor has
$H^0\simeq\mathcal O_X/p$ and
$H^{-1}\simeq\mathcal O_X/p$.  Finally
$0\to\Z\xrightarrow p\Z\to\mathbb F_p\to0$ is the canonical free
resolution of the reduced coefficient.
\end{proof}

Extending this reduced coefficient resolution to $\mathcal O_{\Z}$ gives
the perfect two-term complex
\begin{equation}\label{eq:localcontact}
 C_{p,\Z}=[\mathcal O_{\Z}\xrightarrow{p}\mathcal O_{\Z}].
\end{equation}
After real base change multiplication by $p$ is invertible, so this complex
is canonically contractible.  Hence
\begin{equation}\label{eq:DNcontact}
 C_p=(C_{p,\Z},0,C_{p,\Z}\otimes\mathcal O_{\R}\simeq0)
\end{equation}
is a perfect object of \eqref{eq:PerfDN}.  With the standard torsion metric,
the determinant generator of $[\Z\xrightarrow p\Z]$ has norm $p^{-1}$:
indeed the image lattice has index $p$ in the target lattice.  Therefore
\begin{equation}\label{eq:contactmass}
 -\log\|\det_{\mathrm{tor}}C_p\|=\log p.
\end{equation}
Write $\kappa_p:=\det_{\rm tor}C_p$ for this based normed determinant
line.  This distinction matters: $C_p$ is a perfect torsion complex,
whereas $\kappa_p$ is an invertible object of the determinant Picard
groupoid.
For a label $n$, use $C_p$ when $n=p^k$ and the zero object when $n$ has at
least two prime divisors.  This realizes $\Lambda(n)$ objectwise; it does
not make the false assertion that $\Z/p$ is idempotent under derived tensor.

For finite ruling divisors $x,y$ as in \eqref{eq:Dpr}, tensor the appropriate
powers of the determinant lines $\kappa_p$.  No derived tensor power of
$\mathbb F_p$ is taken.  The resulting contact line
$\lambda_C(x,y)$ has metric exponent
\begin{equation}\label{eq:Clambda}
 C_\Lambda(x,y)=\sum_p
 (x_{p,1}y_{p,2}+x_{p,2}y_{p,1})\log p.
\end{equation}
If $\delta\lambda_{\rm int}(x,y)$ denotes
$\lambda_{\rm int}(x+y)\otimes
\lambda_{\rm int}(x)^{-1}\otimes
\lambda_{\rm int}(y)^{-1}$, define
\begin{equation}\label{eq:greenline}
 \lambda_G(x,y):=\delta\lambda_{\rm int}(x,y)
                  \otimes\lambda_C(x,y)^{-1}.
\end{equation}

\begin{corollary}[Canonical contact--Green factorization]
\label{cor:determinantcomparison}
There is a canonical isometry
\begin{equation}\label{eq:determinantComparison}
 \delta\lambda_{\rm int}
 \simeq\lambda_C\otimes\lambda_G,
\end{equation}
whose logarithmic metric identity is
\begin{equation}\label{eq:green}
 B_{\rm int}=C_\Lambda+G.
\end{equation}
All symmetry, associativity and biextension interchange diagrams commute.
\end{corollary}

\begin{proof}
Equation \eqref{eq:determinantComparison} is obtained by tensoring
\eqref{eq:greenline} with $\lambda_C$.  Taking negative logarithms of the
distinguished-generator norms and using \eqref{eq:intrinsicRRform} and
\eqref{eq:Clambda} gives \eqref{eq:green}.  Every structural map is built
from the associator, symmetry and inverse in the determinant Picard
groupoid.  Mac~Lane coherence, or direct cancellation of the based lines,
proves all stated diagrams.
\end{proof}

The logical direction is part of the statement: $\lambda_G$ is the unique
based normed line defined by the quotient \eqref{eq:greenline}.  Thus the
corollary is a canonical factorization of the intrinsic section determinant,
not an independent construction of a Green current by a second geometric
process.

On the real ruling quotient
\begin{equation}\label{eq:rulingquotient}
 N^1_{\mathrm{rul}}:=
 \left(\bigoplus_p\R e_{p,1}\oplus\bigoplus_p\R e_{p,2}\right)
       /\ker(d_1,d_2)\simeq\R^2,
\end{equation}
the matrix of $B_{\rm int}$ is
$\left(\begin{smallmatrix}0&1\\1&0\end{smallmatrix}\right)$ and has
signature $(1,1)$.  This is the numerical ruling form.  It is not claimed
to contain the completed mixed correspondence space used in row~(d).

\subsection{Exact correspondences and the correct finiteness theorem}

Extend the nuclear coefficient sheaf in \eqref{eq:threeStructures} and put
\begin{equation}\label{eq:NDN}
 \mathcal N_{DN}:=
 \mathcal O_{\mathcal C}e_1\oplus
 \mathcal O_{\mathcal C}e_2\oplus
 \mathcal O_{\mathcal C}e_\Gamma.
\end{equation}
It is locally free of rank three over $\mathcal O_{\mathcal C}$.  The first
two generators are the rulings.  The third is an explicitly adjoined free
coefficient direction for the correspondence algebra; it is not asserted
to be the class of a mixed Cartier divisor on the spherical square.  On
this third summand define
\begin{equation}\label{eq:nuclearAction}
 \rho_n(ae_\Gamma)=(\delta_n*a)e_\Gamma.
\end{equation}
The estimate $q_k(\delta_n*a)\leq n^kq_k(a)$ proves continuity, and
\eqref{eq:DirichletComposition} gives
$\rho_m\rho_n=\rho_{mn}$.  The diagonal contact is the continuous functional
$\ell$, so
\begin{equation}\label{eq:nuclearExactContact}
 \ell(\delta_n)=\Lambda(n),\qquad
 \ell(\delta_m*\delta_n)=\Lambda(mn),
\end{equation}
including even prime powers; no
reduction of the arithmetic label modulo two is involved.

The cotangent objects themselves admit the compatible enrichment
\[
 \mathcal O_{\mathcal C}\otimes_{\R}V_{m,n}.
\]
They are finite free of rank $r(m)r(n)$, and multiplication by $\delta_k$
commutes with restrictions, zero-extensions and the Cartesian-basis
K\"unneth map.  The augmentation \eqref{eq:augmentation} recovers the real
digit determinant, while $\ell$ retains the independent contact mass.

\begin{theorem}[Finite-rank obstruction]\label{thm:finiteNSnogo}
There is no finite-dimensional real vector space $E$, vectors
$\gamma_n\in E$ for every $n\geq1$, and bilinear form $K$ satisfying
\begin{equation}\label{eq:finiteRankImpossible}
 K(\gamma_m,\gamma_n)=\Lambda(mn)
\end{equation}
for all $m,n$.  Consequently no finitely generated abelian numerical group
equipped with classes $\gamma_n$ and a bilinear form can realize the full
two-variable kernel $\Lambda(mn)$ in this way.
\end{theorem}

\begin{proof}
Choose distinct primes $p_1,\ldots,p_r$.  The matrix of $K$ on the vectors
$\gamma_{p_i}$ is
\[
 (\Lambda(p_ip_j))_{i,j}
   =\operatorname{diag}(\log p_1,\ldots,\log p_r),
\]
because a product of two distinct primes is not a prime power.  It has rank
$r$, hence $\dim E\geq r$.  Since $r$ is arbitrary, $E$ cannot be finite
dimensional.  Tensoring a finitely generated abelian group with $\R$ would
produce such an $E$, proving the last assertion.
\end{proof}

Theorem~\ref{thm:finiteNSnogo} rules out finite rank under the displayed
full-kernel hypothesis; it does not classify every conceivable finite
numerical theory.  Formula \eqref{eq:NDN} supplies one explicit replacement:
three module generators suffice after the infinitely many labels are moved
into the scalar nuclear algebra.  No universality or uniqueness of this
extension of scalars is asserted.

\subsection{Completion theorem and audit}

\begin{theorem}[Constructed intrinsic periodic--Deligne--nuclear square]
\label{thm:rowaDN}
The tuple
\begin{equation}\label{eq:ADN}
 \mathscr A_{DN}=
 (\mathscr Y_{\mathbb S},\mathcal O_{\mathscr Y},
  \operatorname{Perf}_{IDN},\mathcal L,\mathcal F_{\rm per},
  \mathbf H^{\rm int},\lambda_{\rm int},\lambda_{\mathrm{code}},
  \lambda_C,\lambda_G,\mathcal N_{DN})
\end{equation}
is an intrinsic periodic--Deligne--nuclear arithmetic square in the sense
of Definition~\ref{def:strongrowa}.  No carrier, descent, section-functor,
regular-cotangent, K\"unneth, continuous-determinant, reduced-contact,
Green-factorization, composition or stated module-finiteness assertion in
that definition is conditional.
\end{theorem}

\begin{proof}
The literal carrier and noncollapse are
\eqref{eq:sphericalsquare} and Proposition~\ref{prop:sphericalnoncollapse}.
Proposition~\ref{prop:divisorlines} gives invertible divisor objects and
principal descent, and Proposition~\ref{prop:externaleffectivity} computes
their effective cone.  Theorem~\ref{thm:intrinsicPeriodicKunneth} constructs
the internal derived section functor and its K\"unneth equivalence.
Theorem~\ref{thm:regularSectionCotangent} constructs the cotangent of its
regular section moduli, while \eqref{eq:intrinsicContinuousDimension}
proves the coefficient-one quadratic dimension.  Proposition
\ref{prop:intrinsicCategoricalGluing} places these objects, their integral
and nuclear realizations and the contact in one stable symmetric monoidal
perfect category, and Proposition~\ref{prop:intrinsicArithmeticAction}
constructs their compatible arithmetic scalar action.
Theorem~\ref{thm:intrinsicCodeComparison} constructs the
continuous determinant and its independent spherical code certificate.
Proposition~\ref{prop:derivedPrimeContact}, equations
\eqref{eq:DNcontact}--\eqref{eq:contactmass}, and Corollary
\ref{cor:determinantcomparison} give the contact determinant and intrinsic
Green factorization.
Proposition~\ref{prop:nuclearDirichlet}
and \eqref{eq:NDN}--\eqref{eq:nuclearAction} give exact composition,
contact and local freeness of rank three.  These verify the five clauses of
Definition~\ref{def:strongrowa}.
\end{proof}

For clarity, Theorem~\ref{thm:rowaDN} neither asserts nor needs that
all raw bounded sections have quadratic dimension: Theorem
\ref{thm:rawsmashnogo} proves the contrary.  It also neither asserts nor
needs finite abelian numerical rank: Theorem~\ref{thm:finiteNSnogo} proves
that finite rank cannot realize the full bilinear kernel $\Lambda(mn)$
required by the particular contact pairing used in rows~(b) and~(c).
The negabinary envelope remains auxiliary.  The intrinsic cohomology is
the cotangent of the regular moduli of actual periodic coefficient sections
on the square.  No assertion equating it with all raw bounded functions is
made or needed.

% END INLINED SOURCE: row-a-deligne-nuclear


\section{The corner pairing}

\subsection{Correspondence potentials}

Let $\Psi_\lambda$ denote the Frobenius correspondence of slope $\lambda>0$
on the square of the Scaling Site \cite{CC-ScalingSite,CC-RRstrategy}, whose
tropical hypersurface on the real positive chart is $\{y=\lambda x\}$, with
convex potential $\max(y-\lambda x,0)$.  For a test function $f$ one forms
\begin{equation}
 D_f=\int_0^\infty f(\lambda)\Psi_\lambda\,\frac{d\lambda}{\lambda},
 \qquad
 U_f(x,y)=\int f(\lambda)\max(y-\lambda x,0)\,\frac{d\lambda}{\lambda}.
\end{equation}
Writing $r=y/x$, $U_f(x,y)=x\,u_f(r)$, and
\begin{equation}\label{eq:defining}
 u_f''(r)=\frac{f(r)}{r},
\end{equation}
which determines $U_f$ uniquely modulo affine functions of $(x,y)$.  We take
\eqref{eq:defining}, not the superposition integral, as the defining
relation; the integral is one construction method, valid on compactly
supported $f$.

\subsection{The interior carries nothing}

\begin{theorem}\label{thm:interior}
Let $U(x,y)=x\,u(r)$ and $V(x,y)=x\,v(r)$ be homogeneous potentials.  Then
\[
 \nabla^2U=\frac{u''(r)}{x}\,\mathbf v_r\mathbf v_r^{\!\top},
 \qquad \mathbf v_r=(r,-1)^{\!\top},
\]
so every such Hessian has rank one, and the mixed density
\[
 U_{xx}V_{yy}+U_{yy}V_{xx}-2U_{xy}V_{xy}
 =\frac{u''v''}{x^2}\big[r^2+r^2-2r^2\big]=0
\]
vanishes identically.
\end{theorem}

Consequently no intersection pairing on the smooth interior of the square can
be nonzero, and every nonvanishing contribution must be supported by data
lost on the punctured chart: the common corner of the Frobenius rays, the
periodic adelic quotient and its isotropy, and the diagonal renormalization.
This is consistent with, and explains, the divergent $\log\Lambda$ term
recorded in \cite{CC-RRstrategy} for the naive diagonal self-intersection.

\subsection{The renormalized corner trace}

For a test function $g$ put
\begin{equation}\label{eq:Tateinvolution}
 g^\vee(x)=x^{-1}\overline{g(x^{-1})}.
\end{equation}
With the Mellin convention used below this is the Tate involution:
$\widehat{g^\vee}(s)=\overline{\widehat g(1-\bar s)}$.

Let $S$ be a finite set of places containing $\infty$, let
$X_S=\Z_S^\times\backslash\A_S$ and $\mathcal H_S=L^2(X_S)$, and let $\theta$
be the scaling representation of the semilocal idele class group.  With
$P_\Lambda$ the position cutoff, $\widehat P_\Lambda$ its Fourier conjugate
and $R_\Lambda=\widehat P_\Lambda P_\Lambda$, set
\begin{equation}\label{eq:cornertrace}
 \mathfrak T_S(h)=\lim_{\Lambda\to\infty}
 \Big(\Tr_{\mathcal H_S}\big(\theta(h)R_\Lambda\big)-2h(1)\log\Lambda\Big).
\end{equation}
The semilocal trace theorem \cite{Connes1999,CCJacobian} gives
\begin{equation}\label{eq:arith}
 \mathfrak T_S(h)=\sum_{v\in S}\int_{\Q_v^\times}'
 \frac{h(u^{-1})}{|1-u|_v}\,d^\times u,
\end{equation}
the local factor arising as a transverse kernel trace,
$\int_{\Q_v}\delta((u-1)x)\,dx=|1-u|_v^{-1}$.  For $h$ compactly supported
in $[e^{-T},e^{T}]$ only primes with $p^k\le e^T$ contribute, so the sum is
finite and $\mathfrak T_S$ stabilizes in $S$.  Setting
$I_\partial(D_f,D_g):=\mathfrak T(f\star g^\vee)$ one obtains, by Tate's
normalization and the global explicit formula, the value
$I_\partial(D_f,D_g)=N(f\star g^\vee)$.

We stress the logical order.  \eqref{eq:cornertrace} is defined from a
representation and two geometric cutoffs, before \eqref{eq:arith} and before
any mention of the explicit formula; the identification with $N$ is a
theorem, not a definition.

\subsection{The numerical quotient}

$I_\partial$ is Hermitian, and for any bilinear form the induced form on the
quotient by its radical is well defined and nondegenerate.  Hence
\begin{equation}
 \overline I_\partial \text{ on } V:=\{\text{correspondence divisors}\}/
 \operatorname{rad}I_\partial
\end{equation}
exists unconditionally.  This is the analogue of numerical equivalence, and
is available with no further hypothesis --- in particular with no principal
invariance.

\subsection{The signature, and why it is not enough}

Write $\rho'=1-\bar\rho$.

\begin{lemma}\label{lem:mirror}
$\rho\mapsto\rho'$ is an involution of the zero multiset of $\xi$ preserving
multiplicities, and $\rho'=\rho$ if and only if $\Rel\rho=\tfrac12$.
\end{lemma}

So the on-line zeros are exactly the fixed points of the mirror involution
and the off-line zeros exactly its two-cycles.  With
$\widehat{g^\vee}(s)=\overline{\widehat g(1-\bar s)}$ the explicit
formula gives
\begin{equation}\label{eq:evalcoords}
 I_\partial(D_f,D_g)=
 \widehat f(0)\overline{\widehat g(1)}+\widehat f(1)\overline{\widehat g(0)}
 -\sum_\rho m_\rho\,\widehat f(\rho)\overline{\widehat g(\rho')},
\end{equation}
which is the classical Weil scalar product \cite{Lagarias2007}.  The form
splits orthogonally: the two polar coordinates carry
$\left(\begin{smallmatrix}0&1\\1&0\end{smallmatrix}\right)$, of signature
$(1,1)$; each on-line zero contributes a negative line; each off-line mirror
pair contributes a hyperbolic plane.  Hence
\begin{equation}
 n_+(\overline I_\partial)=1+\#P,\qquad
 \#P=\tfrac12\#\{\text{off-line zeros}\},
\end{equation}
so $\RH\iff n_+=1$, the classical Hodge-index shape.

\begin{remark}
This is an equivalence derived from the explicit formula and therefore does
no work; it specifies the target statement without supplying it.  The
inertia count itself is due to Connes and van Suijlekom \cite{CvS}.  What is
new here is only the packaging: the form is an intersection pairing on a
numerical quotient of a divisor group, its interior density vanishes
identically (Theorem \ref{thm:interior}), and its polar block is exactly the
intersection matrix $\left(\begin{smallmatrix}0&1\\1&0\end{smallmatrix}\right)$
of the two rulings of a product surface --- which agrees with the bidegree
coefficient $\deg^+(D)\deg^+(E)$ obtained independently from a tropical
tensor construction on a periodic product.
\end{remark}

\section{Mellin descent and its pairing}

\subsection{The descent forces Mellin monomials}

The quotient by the Frobenius action makes strict invariance unavailable.

\begin{proposition}\label{prop:strict}
There is no nonzero $f\in C_c((0,\infty))$ with $U_{n\cdot f}-U_f$ affine for
every $n\in\N^\times$.
\end{proposition}

Indeed the condition reduces, by \eqref{eq:defining}, to $f(r/n)=f(r)$; in
exponential coordinates this is invariance under translation by $\log n$, and
$\{\log n\}$ generates a dense subgroup of $\R$, so a continuous compactly
supported invariant is constant, hence zero.  The incompatible hypothesis is
compact support.

Relaxing invariance to a character twist is the correct move, and it is
rigid.

\begin{theorem}\label{thm:classification}
Let $f\in C((0,\infty))$ be nonzero and suppose $f(r/n)=\chi(n)f(r)$ for all
$n\in\N^\times$.  Then $\chi$ is totally multiplicative, extends without
obstruction to $\Q^\times_+$, and there is a unique $s\in\R$ with
$\chi(q)=q^{-s}$; moreover
\[
 f(r)=c\,r^{s},\qquad c\in\R^\times .
\]
\end{theorem}

The proof runs through exponential coordinates: continuity of $f$ at a point
where it does not vanish forces $\chi\circ\exp$ continuous at $0$, hence
continuous on the dense group $\Gamma=\log\Q^\times_+$, hence uniquely
extendable to a continuous homomorphism $(\R,+)\to(\R_{>0},\times)$, which is
$a\mapsto e^{-sa}$.  Substituting back, $F(t-a)=e^{-sa}F(t)$ for all real
$a,t$, and setting $a=t$ gives $F(t)=F(0)e^{st}$.

\begin{remark}
There is \emph{no} residual log-periodic factor.  A genuine period would
require the invariance group to be discrete of rank one; here it is dense,
which forces the periodic factor to be constant.  The descent therefore
expels the sections from compact support into a one-parameter graded family
of Mellin monomials, and the grade is the spectral variable.
\end{remark}

\subsection{The graded family is dual to the test class}

Write $f_s(x)=x^{-s}$ for the graded family, so that $f_s(u^{-1})=|u|^{s}$
under the module embedding.

\begin{proposition}\label{prop:nomellin}
The Mellin integral $\int_0^\infty x^{-s}x^{w}\,\frac{dx}{x}
=\int_\R e^{t(w-s)}\,dt$ converges for no $w\in\C$.  In particular $f_s$
lies in no space $A_\delta$ of test functions whose Mellin transform is
analytic on a strip about the critical line.
\end{proposition}

Under the cutoff $[\tfrac1T,T]$ used in the covariance form of the explicit
formula \cite{BombieriLagarias,Lagarias2007}, however,
\begin{equation}
 \int_{1/T}^{T}x^{-s}x^{w}\frac{dx}{x}
 =\frac{2\sinh\big((w-s)\log T\big)}{w-s},
\end{equation}
which on the vertical line $w=s+iu$ is the Dirichlet kernel
$2\sin(u\log T)/u$, of total mass $2\pi$ for every $T$, converging weakly to
$2\pi\delta$.

\begin{corollary}\label{cor:duality}
$\Gr:=\{f_s\}$ is not a subspace of the test class; each $f_s$ defines
instead a measure supported on the vertical line $\Rel w=s$.  The graded
family is Mellin-dual to the test class, and the cutoff of the covariance
form is the pairing between the two sides.
\end{corollary}

This is the structural reason the naive formalization of the preceding descent --- namely
substituting $f_s$ into \eqref{eq:arith} --- is a category error, and it
governs everything in \S4.4--\S4.6.

\subsection{Local terms, in closed form, on the critical strip}

Monomials diagonalise convolution with the Tate involution
\eqref{eq:Tateinvolution}:
\begin{equation}\label{eq:eigen}
 f_s\star g^\vee=c_g(s)\,f_s,
 \qquad
 c_g(s)=\int_0^\infty t^{1-s}\overline{g(t)}\,d^\times t,
\end{equation}
entire in $s$ for compactly supported $g$.  So the problem reduces to a
single monomial.

\begin{theorem}\label{thm:localterms}
Decomposing $\Q_p^\times$ into shells $|u|_p=p^{-n}$ of $d^\times u$-measure
$1$,
\[
 W_p(f_s)=\frac{p^{-s}}{1-p^{-s}}+\frac{p^{s-1}}{1-p^{s-1}}+C_p,
 \qquad
 C_p=\int_{\Z_p^\times}'\frac{d^\times u}{|1-u|_p}
\]
with $C_p$ independent of $s$; the two series converge if and only if
$\Rel s>0$ and $\Rel s<1$ respectively.  At the archimedean place
\[
 W_\infty(f_s)=\int_{\R^\times}'\frac{|u|^{s-1}}{|1-u|}\,du
 =\frac{\pi}{\sin\pi s}+\pi\cot\pi s=\pi\cot\frac{\pi s}{2},
\]
regular on $0<\Rel s<1$, with value $\pi$ at $s=\tfrac12$, and convergent
on exactly the same strip: the integrand is $|u|^{s-1}$ near $0$ and
$|u|^{s-2}$ near $\infty$.
\end{theorem}

\begin{corollary}
Every local term, finite and archimedean alike, converges precisely on the
critical strip $0<\Rel s<1$, with both boundaries excluded.
\end{corollary}

We record, without overstating it, that the strip here arises from purely
local analysis of a monomial, with no input from the zeros or the functional
equation.  The symmetry $s\leftrightarrow 1-s$ interchanging the two series
is built into the kernel $|u|^{s}/|1-u|$; this is a consistency check that
the construction sits in the right place, not a discovery.

\subsection{The $s$-dependent assembly, and the divergent constant}

Write $A(s)=\sum_p p^{-s}/(1-p^{-s})$ and $B(s)=A(1-s)$.  Each is a Dirichlet
series over prime powers, convergent respectively on $\Rel s>1$ and
$\Rel s<0$ --- half-planes disjoint from the strip and exchanged by
$s\mapsto 1-s$.  Their analytic continuations are governed by
$A(s)=\sum_{N\ge1}\frac{\varphi(N)}{N}\log\zeta(Ns)$, whose singularities come
from the pole of $\zeta$ at $1$.

\begin{theorem}\label{thm:locint}
$A+B$ is locally integrable on the open strip, hence defines a distribution
of order zero there.
\end{theorem}

The singular set $\{1/N\}\cup\{1-1/M\}$ is finite on every compact
subinterval and accumulates only at the endpoints; each singularity is
logarithmic, and $\log$ is locally integrable; the tail converges uniformly,
since $\log\zeta(\sigma)\cdot2^{\sigma}\to1$ geometrically.  Pointwise the two
halves reinforce at $s=\tfrac12$, which is the unique fixed point of
$s\mapsto1-s$ and where $B=A$ identically, so no cancellation is possible
there; but pointwise divergence and local integrability are compatible, and
it is the latter that Corollary \ref{cor:duality} requires.

The constant is a different matter.

\begin{theorem}\label{thm:constant}
Each shell $|1-u|_p=p^{-k}$, $k\ge1$, has measure $p^{-k}$ and integrand
$p^{k}$, hence contributes exactly $1$; so $C_p=+\infty$.  Moreover the
$k=0$ shell contributes the finite quantity $(p-2)/(p-1)\to1$, whose sum over
primes grows like $\pi(x)$.  The constant therefore diverges for two
independent reasons, only the first of which is a principal-value artefact.
\end{theorem}

\subsection{Why the constant cannot be removed}

We record the repairs that fail, since each localises the obstruction
further.

\begin{proposition}\label{prop:norepair}
\begin{enumerate}
\item The counterterm $2h(1)\log\Lambda$ of \eqref{eq:cornertrace} is linear
in $\log\Lambda$, whereas the prime sum over the stabilizing set
$\{p\le e^{T}\}$ grows like $\pi(e^{T})$.  The scales are incommensurate.
\item The zeta-regularized value $-\zeta'/\zeta(0)=-\log2\pi$ is finite, but
the identification fails: the $(p-2)/(p-1)$ part is untouched by the
regulator, and the continuation has $s=0$ as an accumulation point of poles.
\item Restricting $g$ by a primitivity condition cannot help: $c_g$ is
entire, and an entire function not identically zero cannot vanish on an
interval.
\item Membership in $\operatorname{rad}I_\partial$ is excluded by an explicit
nonvanishing witness.
\end{enumerate}
\end{proposition}

The unifying reason is an identity constraint.

\begin{theorem}\label{thm:identity}
The finite-place principal-value local term is finite, with no residual
constant, if and only if the test argument vanishes at the identity,
$h(1)=0$.  Every element of the graded family violates this: $f_s(1)=1$ for
every $s\in\C$, because $\chi(1)=1$ is the definition of a quasi-character.
\end{theorem}

So the divergence is not an artefact of the shell splitting, and not a
normalization one may choose away.  It is a group identity.  In parallel,
there is no continuous extension of $\Lambda_g(f):=\mathfrak T(f\star
g^\vee)$ from the test class to individual $f_s$: the functional is
represented by a discrete measure supported on $\{0,1\}\cup\{\rho\}$, and
$\widehat f_s$ is a measure on a vertical line, so the pairing is a product
of two singular objects.  Smearing in $s$ does produce a convergent sum
$\sum_\rho\varphi(\Rel\rho)\overline{\widehat g(\rho')}$, but that sum is a
zero-sourced definition and is inadmissible under the constraint of
\S\ref{ss:constraint}.

\begin{remark}\label{rem:replaces}
Theorem \ref{thm:identity} and the two statements following it settle, in the
negative, the formulation one would naturally attempt first: extend the
corner trace to individual graded elements and identify the result with the
zero-side sum.  That programme is not merely unachieved; it is impossible,
and impossible for a reason internal to the objects --- a quasi-character
takes the value $1$ at the identity, and finiteness of the local term
requires the opposite.  The construction of \S4.6 is therefore not a
weakened substitute but the statement the obstruction itself selects: since
the divergence enters with an $s$-independent coefficient, it is confined to
the total-mass direction and disappears on differences.
\end{remark}

\subsection{The obstruction is confined to the total-mass direction}

The resolution is that the divergence enters with coefficient $f_s(1)=1$,
which by Theorem \ref{thm:identity} is \emph{independent of $s$}.  Write
$\Lambda_g(f_{s,T})=D(T)+L_g(s,T)$ with $D(T)$ the $s$-independent divergent
piece and $L_g$ the assembly of Theorem \ref{thm:locint}.

\begin{lemma}\label{lem:cancel}
For coefficients $\lambda_i$ with $\sum_i\lambda_i=0$,
\[
 \sum_i\lambda_i\big[D(T)+L_g(s_i,T)\big]
 =D(T)\sum_i\lambda_i+\sum_i\lambda_iL_g(s_i,T)
 =\sum_i\lambda_iL_g(s_i,T).
\]
The cancellation is an identity in the coefficients, valid at every finite
$T$, and uses no smoothness of the combination.
\end{lemma}

It remains to know that mass-zero combinations of principal elements exist.
Here the definition of ``principal'' must be the standard equivariant one.

\begin{proposition}\label{prop:semiinv}
$\Div$ is additive over sums of potentials and annihilates affine functions;
it is therefore already logarithmic in the tropical dictionary, in which the
multiplicative operation is addition of piecewise-linear functions, and
$\Div(cU)=c\,\Div(U)$ is the correct power law.  Requiring $\Div(U)$ to be
strictly invariant forces the trivial character and hence weight $s=0$, which
lies on the excluded boundary of the strip and at an accumulation point of
the singular set.  Under the standard convention --- principal divisors are
the divisors of semi-invariant sections --- there are nonzero principal
witnesses $\divv(U_s)$ at every real weight, in particular at every
$s\in(0,1)$.
\end{proposition}

\begin{theorem}[Mellin descent theorem]\label{thm:mellindescent}
Let $\Prin'(\Gr)$ be the subgroup generated by $\{\divv(U_s):s\in\R\}$.  Let
$s_0,s_1\in(0,1)$ be distinct and off the singular set
$\{1/N\}\cup\{1-1/M\}$, and let $g$ be admissible.  Then
$\divv(U_{s_0})-\divv(U_{s_1})\in\Prin'(\Gr)$ has total mass zero, and
\[
 \boxed{\;
 \Lambda_g^{0}\big(\divv(U_{s_0})-\divv(U_{s_1})\big)=L_g(s_0)-L_g(s_1)\;}
\]
is finite, defined without reference to any zero of $\xi$, with the divergent
constant contributing exactly zero at every finite regularization depth.
\end{theorem}

\begin{proof}
Membership holds because $\Prin'$ is by construction a subgroup, as
$\Prin$ always is so that $\operatorname{Pic}=\Div/\Prin$ is defined.  Each
witness has grade-mass $1$, so the difference has mass $0$.  The evaluation
and the exact vanishing of the constant are Lemma \ref{lem:cancel}.  The
singular set is countable, so generic $s_0,s_1$ avoid it and $L_g$ has honest
pointwise values there by Theorem \ref{thm:locint}.
\end{proof}

\begin{corollary}
Principal invariance is a well-formed statement.  The objection that the
local principal subspace is trivial, so that the question could not even be
posed, is answered: it is posed on $\Prin'(\Gr)$, and Theorem
\ref{thm:mellindescent} evaluates the pairing on it.
\end{corollary}

\section{Scope of Theorem \ref{thm:mellindescent}}

Two of the following are limitations of the theorem; the rest fix its
reading.

\begin{itemize}
\item \textbf{Whether principal invariance holds is not addressed here.}
Theorem \ref{thm:mellindescent} makes the question askable.  Answering it is
the principal-difference analysis below.  This is the most likely misreading of the paper and we state it
first.

\item \textbf{The semi-invariant convention is a revision.}  Proposition
\ref{prop:semiinv} adopts the standard equivariant definition in place of a
stricter one.  The revision is justified --- the stricter reading would
trivialize the group of invariant principal divisors already on
$\mathbb P^1$ under a torus action --- but it is a change of definition and
should be read as such.  It is the load-bearing choice of this paper.

\item \textbf{The constant is not evaluated.}  $\sum_pC_p$ remains unknown
under any regularization.  Theorem \ref{thm:mellindescent} does not need it, but
only because it works with differences; any statement requiring an absolute
normalization would still be blocked.

\item \textbf{Nothing here bears on $\RH$.}
\end{itemize}

\section{Principal differences and the radical}

The descent question is whether the pairing of Theorem \ref{thm:mellindescent}
kills principal elements.  Answering it needs one ingredient not yet
available: the finite part of the preceding local sum in closed form.  That
is Theorem \ref{thm:assembly} of the next section, which we use here as a
forward reference; nothing in the correspondence construction depends on the present section, so
there is no circularity.  Granting it, the object of Theorem
\ref{thm:mellindescent} acquires a closed form:
\begin{equation}\label{eq:Lg}
 L_g(s)=c_g(s)\,\Phi(s),
 \qquad
 \Phi(s)=\pi\cot\frac{\pi s}{2}-\frac{\zeta'}{\zeta}(s)
 -\frac{\zeta'}{\zeta}(1-s),
\end{equation}
with $\Phi$ independent of $g$.  By Theorem \ref{thm:mellindescent} the descent
question is therefore whether $\Phi$ is constant.

\begin{lemma}\label{lem:phipsi}
$\Phi(s)=\pi\cot(\pi s/2)+\tfrac12\big[\psi(s/2)+\psi((1-s)/2)\big]-\log\pi$.
\end{lemma}

This follows from logarithmic differentiation of $\xi(s)=\xi(1-s)$, the polar
terms cancelling in the symmetric combination.  It is what makes $\Phi$
computable without reference to the zeros.

\begin{theorem}\label{thm:principalradical}
\begin{enumerate}
\item $\Phi$ is not constant on $(0,1)$; hence principal invariance fails
and the pairing does not descend to $\operatorname{Pic}$.
\item $\Phi\not\equiv0$; in particular $\Phi(\tfrac12)=-2.2305907\ldots\ne0$,
so the pairing is not vacuous.
\item $\operatorname{rad}\Lambda^{0}$ is spanned by the point masses at the
zeros of $\Phi$: writing
$\Lambda_g^{0}(\sum_i\lambda_i\delta_{s_i})=\sum_i\lambda_ic_g(s_i)\Phi(s_i)$
and using linear independence of the Mellin evaluations at distinct $s_i$,
one gets $\lambda_i=0$ unless $\Phi(s_i)=0$.
\item On the real segment $(0,1)$, $\Phi$ has exactly one sign change, at
$s^{\ast}=0.3016923882\ldots$
\end{enumerate}
\end{theorem}

Item (2) is a control that had to be run before building further: the
functional equation gives $(\xi'/\xi)(s)+(\xi'/\xi)(1-s)\equiv0$, so
$\Phi$ vanishing identically was a genuine possibility, and would have made
the radical everything and the construction empty.

That principal invariance fails is not fatal to the programme.  Weil's
argument does not use the pairing on $\operatorname{Pic}$ but on classes
modulo algebraic or numerical equivalence, and the available Hodge-index
statements are likewise formulated on $NS\otimes\Q$.  A pairing descending to
the quotient by a discrete radical is an object of that type.  What is not
established here is whether that quotient coincides with, or maps to, the
numerical quotient of \S3.4; the two are constructed differently and their
comparison --- which is what would join the two halves of the programme ---
is not attempted.  There is no formal contradiction with
Theorem~\ref{thm:rowaDN}: the pairing in this section is the analytic
Mellin--corner candidate on completed mixed correspondences, whereas the
form \eqref{eq:intrinsicRRform} of row~(a) is the intrinsic periodic cotangent determinant on the two
prime-generated rulings.  The nuclear realization of
Section~\ref{sec:rowc} joins their contact data at the level of categorical
trace.  Extending the intrinsic periodic determinant and effectivity from
the ruling lines to the completed mixed classes is still required for
row~(d); it is not a hypothesis of the nuclear Lefschetz identity.


% BEGIN INLINED SOURCE: row-b-witt-lefschetz
\section{Witt correspondences and their local Lefschetz complexes}
\label{sec:correspondences}

The analogue of a Frobenius graph used here is a cohomological
correspondence supported by the diagonal, not an algebraic cycle with a
moving support.  This distinction is structural.  The endomorphisms which
exist over $\Z$ act on the Witt coefficient sheaf, and their local contact
must therefore be extracted from that action.  We construct it below as a
derived intersection of the cyclic Witt orbit with the unit section.  No
finite module is appended in order to prescribe an intersection number.

\subsection{Witt dynamics}

Let
\[
 W_{\Z}:=\bigoplus_{r\geq1}\Z\phi_r
\]
be Haran's global cyclotomic Witt module.  We use the sheaf
$\mathcal W_X:=\underline{W_{\Z}}$ obtained by sheafifying the constant
presheaf on the topological space of the absolute arithmetic curve $X$.
This is an explicitly defined coefficient system on the spherical carrier;
we do not identify Haran's prop-scheme category with the category of
$\Gamma$-ringed spaces used in row~(a).  Haran's Frobenius maps on $W_\Z$
induce stalkwise sheaf endomorphisms
\begin{equation}\label{eq:Frobmult}
 F_mF_n=F_{mn},\qquad F_1=\operatorname{id}_{\mathcal W_X}.
\end{equation}
The constant sections give a distinguished invariant cyclic submodule
\begin{equation}\label{eq:Wittbasis}
 \bigoplus_{r\geq1}\Z\phi_r\ \subseteq\ \mathcal W_X(X).
\end{equation}
For the Hilbert structure defined by Haran's state, the adjoints
$V_n=F_n^*$ satisfy
\begin{equation}\label{eq:versch}
 V_n\phi_r=\phi_{nr},\qquad V_mV_n=V_{mn}
\end{equation}
\cite[(12.35),(12.37)]{Haran2022}.  Haran's lambda characteristic is
\begin{equation}\label{eq:HaranLambdaCharacteristic}
 \chi_T(x):=\operatorname{tr}(\lambda_T(x)),\qquad
 \chi_T(\phi_r)=\Phi_r(T),
\end{equation}
where $\Phi_r$ is the cyclotomic polynomial
\cite[(13.37)]{Haran2022}.  Consequently the polynomial extracted from the
lambda characteristic of the canonical cyclic vector $V_n\phi_1$ is
\begin{equation}\label{eq:dynamicPolynomial}
 P_{\Gamma_n}(T):=\chi_T(V_n\phi_1)=\Phi_n(T).
\end{equation}

\subsection{The dynamic local contact}

Put $R=\Z[T,T^{-1}]$.  For $n\geq1$ define the derived intersection
\begin{equation}\label{eq:derivedWittContact}
 K_n^W:=
 R/(P_{\Gamma_n}(T))
 \mathbin{\otimes_R^{\mathbf L}}R/(T-1).
\end{equation}
This is the intersection of the cyclotomic lambda-characteristic locus
$\operatorname{Spec}R/(P_{\Gamma_n})$ associated with $F_n$ and the unit
section $T=1$.  We do not claim that this affine locus represents a full
geometric orbit of $F_n$.  Resolving the first quotient by
$[R\xrightarrow{P_{\Gamma_n}(T)}R]$ gives the canonical model
\begin{equation}\label{eq:derivedWittContactModel}
 K_n^W\simeq
 [\Z\xrightarrow{\,P_{\Gamma_n}(1)\,}\Z]
\end{equation}
in degrees $-1,0$.  Thus the complex, its differential and its determinant
are all functorially determined by $F_n$, the canonical vector $\phi_1$,
the adjoint, and the lambda characteristic.

\begin{proposition}[Cyclotomic local contact]
\label{prop:cyclotomicFixedContact}
For $n>1$, $K_n^W$ is perfect and concentrated in cohomological degree
zero, where
\begin{equation}\label{eq:contactCohomology}
 P_n:=H^0(K_n^W)\simeq\Z/\Phi_n(1)\Z.
\end{equation}
For the identity put $P_1=\Z$.  Moreover
\begin{equation}\label{eq:cyclotomicAtOne}
 P_n\simeq
 \begin{cases}
  \mathbb F_p,&n=p^k,\\
  0,&n>1\text{ has at least two distinct prime divisors}.
 \end{cases}
\end{equation}
\end{proposition}

\begin{proof}
Equation~\eqref{eq:derivedWittContactModel} proves perfectness.  Since
$\Phi_n(1)\ne0$ for $n>1$, its left differential is injective and its
cokernel is \eqref{eq:contactCohomology}.  To evaluate the scalar, divide
$T^n-1=\prod_{d\mid n}\Phi_d(T)$ by $T-1$ and let $T\to1$:
\[
 n=\prod_{\substack{d\mid n\\d>1}}\Phi_d(1).
\]
Induction gives $\Phi_{p^k}(1)=p$ along a prime-power tower.  If $n$ has
at least two prime divisors, the factors belonging to its proper divisors
already multiply to $n$, so $\Phi_n(1)=1$.  This proves both cases.
For $n=1$, $P_{\Gamma_1}=T-1$ and the reduced degree-zero self-contact is
$H^0([\,\Z\xrightarrow0\Z\,])=\Z$, the required unit.
\end{proof}

The reduced contacts do possess a composition law, but it is not the
derived tensor product of their full perfect resolutions.  Let
$P(a)=\Z/a\Z$ for $a\geq0$, with $P(0)=\Z$, and define
\begin{equation}\label{eq:reducedContactProduct}
 P(a)\star P(b):=
 H^0\!\left(P(a)\mathbin{\otimes_\Z^{\mathbf L}}P(b)\right)
 =P(\gcd(a,b)).
\end{equation}
The equality follows from the two-term free resolution of either cyclic
module.  The positive generator of the ideal $(a,b)$ makes the displayed
identification canonical.  Associativity, symmetry and the unit $P(0)$
follow either from the associativity of ordinary tensor product on the
degree-zero quotients or directly from the corresponding properties of
$\gcd$.

Set $a_1=0$ and $a_n=\Phi_n(1)$ for $n>1$.  Formula
\eqref{eq:cyclotomicAtOne} gives
\begin{equation}\label{eq:cyclotomicGcdLaw}
 a_{mn}=\gcd(a_m,a_n).
\end{equation}
Indeed, if both integers are powers of the same prime, all three values are
that prime; in every other nonunit case the product has mixed prime support
and both sides are $1$.  The cases involving $1$ use $\gcd(0,a)=a$.
Therefore
\begin{equation}\label{eq:contactComposition}
 P_m\star P_n\simeq P_{mn}.
\end{equation}
This is the correct reduced-contact law.

It would be false to replace $\star$ by the full derived tensor product.
For example, a resolution of $\mathbb F_p\otimes_\Z^{\mathbf L}\mathbb
F_p$ has both $\operatorname{Tor}_1^\Z(\mathbb F_p,\mathbb F_p)$ and
$\mathbb F_p\otimes\mathbb F_p$, whereas $K_{p^2}^W$ has only one
cohomology group.  Thus no K\"unneth formula for an ambient cotangent
$H_1$ is being assumed.

\subsection{The correspondence category and its arithmetic family}

We now give a category requiring no unspecified bicategorical
infrastructure.  Let $\operatorname{End}(\mathcal W_X)$ be the discrete
monoidal category whose objects are Witt-sheaf endomorphisms and whose
product is composition.  Define
$\operatorname{End}_{\mathcal O_{\mathcal C}}(\mathcal N_{DN})$ in the
same way for \eqref{eq:NDN}.  Let
$\operatorname{Pic}^{\mathrm{met}}_{\mathrm{ext}}(X)$ be the framed
metrized factorized Picard category of row~(a), and let
$\operatorname{Cont}_{\mathrm{cyc}}$ be the symmetric monoidal category
generated by the pointed cyclic modules $P(a)$ with product $\star$.
Take their product monoidal category
\begin{equation}\label{eq:explicitCorrespondenceCategory}
 \mathsf E_{W,DN}(X):=
 \operatorname{End}(\mathcal W_X)\times
 \operatorname{End}_{\mathcal O_{\mathcal C}}(\mathcal N_{DN})\times
 \operatorname{Pic}^{\mathrm{met}}_{\mathrm{ext}}(X)\times
 \operatorname{Cont}_{\mathrm{cyc}}.
\end{equation}
Associators and units are componentwise, so coherence in the four factors
proves coherence of \eqref{eq:explicitCorrespondenceCategory}.
Its support functor sends every object below to the diagonal
$X\to\mathscr Y_{\mathbb S}$, which is proper over both factors.

For $n=\prod_pp^{v_p(n)}$, let
\[
 D(n)=\sum_pv_p(n)[p].
\]
The line $\mathcal O_X(D(n))$, with its canonical section, is one of the
actual divisor modules of row~(a).  Equivalently, pull
$\mathcal L(D(n),0)$ back along the diagonal.  Give its positive frame the
Arakelov norm
\[
 \|\mathbf1_{D(n)}\|:=\exp(-\deg D(n))=n^{-1};
\]
this explicitly defines the framed metrized object $\mathcal T_n$.
Additivity of Arakelov degree and the tensor law of
Proposition~\ref{prop:divisorlines} give
\begin{equation}\label{eq:Tnmetric}
 \nu_\infty(\mathcal T_n)=\deg D(n)=\log n,\qquad
 \mathcal T_m\otimes\mathcal T_n\simeq\mathcal T_{mn}.
\end{equation}
On the nuclear module define $\rho_n$ by
\eqref{eq:nuclearAction}.  Finally set
\begin{equation}\label{eq:Gamman}
 \Gamma_n=(F_n,\rho_n,\mathcal T_n,P_n)
 \in\mathsf E_{W,DN}(X),
\end{equation}
where the last component is not additional data: it is the output
$H^0(K_n^W)$ of \eqref{eq:derivedWittContact}.
Define the \emph{arithmetic Witt correspondence category}
$\mathsf{Corr}_{W,DN}^{\mathrm{arith}}(X)$ to be the graph-image
subcategory with precisely the objects $\Gamma_n$ and the componentwise
isomorphisms between them.  Its fourth component is therefore constrained
by $P_n=H^0(K_n^W)$ and cannot be varied independently.  Equations
\eqref{eq:Frobmult}, \eqref{eq:nuclearAction}, \eqref{eq:Tnmetric} and
\eqref{eq:contactComposition} prove that this graph-image is closed under
the ambient monoidal product and contains its unit.

\begin{theorem}[Composition and faithfulness]\label{thm:rowbcompose}
There are coherent canonical isomorphisms
\[
 \Gamma_m\circ\Gamma_n\simeq\Gamma_{mn},\qquad
 \Gamma_1\simeq\Delta,
\]
and $\Gamma_m\simeq\Gamma_n$ implies $m=n$.
\end{theorem}

\begin{proof}
The Witt component composes by \eqref{eq:Frobmult}.  The nuclear component
composes because $\delta_m*\delta_n=\delta_{mn}$.  The metric component
composes by \eqref{eq:Tnmetric}, and the reduced contact component by
\eqref{eq:contactComposition}.  Coherence is componentwise.  If two
correspondences are isomorphic, their Witt maps agree; taking adjoints and
applying to $\phi_1$ gives $\phi_m=\phi_n$ by \eqref{eq:versch}, hence
$m=n$ by \eqref{eq:Wittbasis}.  The nuclear action gives a second proof:
$\rho_n(\delta_1e_\Gamma)=\delta_ne_\Gamma$.  The metric
$\nu_\infty(\mathcal T_n)=\log n$ gives a third.
\end{proof}

\subsection{Comparison with the perfect and nuclear structures of row (a)}

For $n>1$, the realification of \eqref{eq:derivedWittContactModel} is
contractible because $\Phi_n(1)\ne0$.  Hence
\begin{equation}\label{eq:DNLefschetzComplex}
 \mathbb L_n=
 \bigl(\mathcal O_{\Z}\otimes_\Z K_n^W,\ 0,\
 (\mathcal O_{\Z}\otimes K_n^W)\otimes_{\mathcal O_{\Z}}
 \mathcal O_{\R}\simeq0\bigr)
 \in\operatorname{Perf}_{DN}(\mathscr Y_{\mathbb S}).
\end{equation}
It is perfect because its integral component is a two-term complex of
finite free modules.  Since its nuclear component is zero,
Proposition~\ref{prop:intrinsicCategoricalGluing} gives a canonical lift to
$\operatorname{Perf}_{IDN}$ with zero periodic component.  Equations
\eqref{eq:derivedWittContactModel} and \eqref{eq:cyclotomicAtOne} give the
canonical equivalence
\begin{equation}\label{eq:contactComparisonA}
 \mathbb L_n\simeq
 \begin{cases}
  C_p,&n=p^k,\\
  0,&n\text{ has at least two prime divisors},
 \end{cases}
\end{equation}
where $C_p$ is exactly the object \eqref{eq:DNcontact} constructed in
row~(a).  This is an objectwise equivalence in
$\operatorname{Perf}_{DN}$ and, after the zero-periodic lift, in
$\operatorname{Perf}_{IDN}$; it is not an equality of numerical masses.

With the standard torsion metric, \eqref{eq:derivedWittContactModel} gives
\begin{equation}\label{eq:contactdegreeagain}
 -\log\|\det_{\mathrm{tor}}\mathbb L_n\|
 =\log\Phi_n(1)=\Lambda(n).
\end{equation}
The first equality is the determinant of the dynamic intersection, and the
second is \eqref{eq:cyclotomicAtOne}.

\begin{theorem}[Local Witt--Lefschetz contact]\label{thm:rowbcontact}
For every $n>1$, the derived unit contact of the cyclotomic
lambda-characteristic locus associated with $F_n$ is
the perfect intrinsic Deligne--nuclear object $\mathbb L_n$.  It is canonically
equivalent to the finite contact of row~(a), its determinant degree is
$\Lambda(n)$, and its reduced contact is compatible with correspondence
composition by \eqref{eq:contactComposition}.
\end{theorem}

\begin{proof}
The dynamic derivation is
\eqref{eq:dynamicPolynomial}--\eqref{eq:derivedWittContactModel};
Proposition~\ref{prop:cyclotomicFixedContact} computes its cohomology;
\eqref{eq:contactComparisonA} proves the comparison in
$\operatorname{Perf}_{DN}$ and hence, by the canonical lift, in
$\operatorname{Perf}_{IDN}$; and \eqref{eq:contactdegreeagain} computes its
determinant.  The last assertion is the proved gcd law
\eqref{eq:cyclotomicGcdLaw}.
\end{proof}

The connection with the nuclear correspondence module is equally
explicit.  On the algebraic Witt basis define
\begin{equation}\label{eq:WittNuclearIntertwiner}
 J:\bigoplus_{r\geq1}\R\phi_r\longrightarrow\mathcal C_{\R},
 \qquad J(\phi_r)=\delta_r.
\end{equation}
It is injective, and
\begin{equation}\label{eq:WittNuclearIntertwining}
 J(V_n\phi_r)=\delta_{nr}=\delta_n*\delta_r
 =\rho_n(J(\phi_r)).
\end{equation}
Thus it intertwines the actual Witt dynamics on its algebraic cyclic domain
with the action on the $e_\Gamma$ summand of $\mathcal N_{DN}$.  Applying
the continuous contact functional gives
\begin{equation}\label{eq:nuclearContactComparison}
 \ell(J(V_n\phi_1))=\ell(\delta_n)=\Lambda(n)
 =-\log\|\det_{\mathrm{tor}}\mathbb L_n\|.
\end{equation}
The Witt orbit, the finite perfect contact and the nuclear action are
therefore joined by explicit morphisms and determinant comparison.

\subsection{Central metric normalization}

This normalization can be made internal to the based metric Picard
groupoid already used in row~(a), without invoking an unconstructed
Poincar\'e duality.  On positive based metric lines define
\[
 \mathfrak d_n(L):=\mathcal T_n\otimes L^{-1}.
\]
The canonical cancellations show
$\mathfrak d_n^2(L)\simeq L$, so this is an involution.  A fixed point
satisfies $L^{\otimes2}\simeq\mathcal T_n$.  It exists by the real metric
power operation $L^{\langle c\rangle}$ used in
Proposition~\ref{prop:continuousRRdet}, and it is unique among positive
based lines because its norm is the positive square root of
$\|\mathcal T_n\|=n^{-1}$.  Consequently
\begin{equation}\label{eq:Tatecharacter}
 w_{1/2}(\Gamma_n)=
 \exp\!\left(-\tfrac12\nu_\infty(\mathcal T_n)\right)=n^{-1/2}
\end{equation}
is the unique self-dual positive metric character.  It is multiplicative
because $\mathcal T_m\otimes\mathcal T_n\simeq\mathcal T_{mn}$, and the
normalized contact coefficient is
\begin{equation}\label{eq:centralLocalCoefficient}
 w_{1/2}(\Gamma_n)
 \bigl(-\log\|\det_{\mathrm{tor}}\mathbb L_n\|\bigr)
 =\frac{\Lambda(n)}{\sqrt n}.
\end{equation}
This is a self-dual metric normalization.  It is not asserted to be an
algebraic half-Tate line or a half-integral Cartier divisor.

\begin{corollary}[Constructed row (b)]\label{cor:rowb}
The family $\{\Gamma_n\}_{n\geq1}$ is a faithful multiplicative family of
properly supported Witt cohomological correspondences on the spherical
square.  Its local contact is derived from $F_n$, has a canonical lift to
$\operatorname{Perf}_{IDN}(\mathscr Y_{\mathbb S})$, agrees canonically
with the contact object of row~(a), and has determinant degree
$\Lambda(n)$.  Its faithful nuclear realization intertwines composition,
and its self-dual metric normalization is $\Lambda(n)/\sqrt n$.  Thus
row~(b) is constructed completely in this cohomological
Deligne--nuclear sense.
\end{corollary}

\begin{proof}
The ambient category and its arithmetic graph-image are
\eqref{eq:explicitCorrespondenceCategory} and the paragraph following
\eqref{eq:Gamman};
Theorem~\ref{thm:rowbcompose} proves composition and faithfulness;
Theorem~\ref{thm:rowbcontact} constructs and computes the contact;
\eqref{eq:WittNuclearIntertwining} proves the nuclear comparison; and
\eqref{eq:Tatecharacter} proves the metric assertion.  The diagonal support
is proper.  None of the definitions uses a zero of $\xi$, the explicit
formula or a positivity assertion.
\end{proof}

The conclusion is intentionally cohomological.  Forgetting the coefficient
maps and frames sends every $\Gamma_n$ to the same diagonal, so no distinct
undecorated Chow cycles are claimed.  The result proved here is instead
that the dynamics composes, its lambda-characteristic locus forces the perfect local
Lefschetz complex, and that complex is the contact already internal to the
Deligne--nuclear spherical square.  Section~\ref{sec:rowc} constructs the
global trace of this family.

% END INLINED SOURCE: row-b-witt-lefschetz


\subsection{The shell functional}

\begin{definition}
For $p$ prime and $k\in\Z\setminus\{0\}$ set
$\Gamma_{p,k}(h):=p^{\min(k,0)}h(p^{k})$, a scaled Dirac evaluation at the
point $p^{k}\in(0,\infty)$.
\end{definition}

\begin{proposition}\label{prop:pairing3}
$\Gamma_{p,k}$ is linear, and on the graded family
\[
 \Gamma_{p,k}(f_s)=p^{-ks}\ (k\ge1),
 \qquad
 \Gamma_{p,-m}(f_s)=p^{m(s-1)}\ (m\ge1).
\]
Resumming the two families recovers the two closed forms of Theorem
\ref{thm:localterms}.
\end{proposition}

\subsection{The normalization, and the Weil coefficient}

Theorem \ref{thm:localterms} was stated with $d^\times u$ normalized by
$\mathrm{vol}(\Z_p^\times)=1$, in which every shell has mass $1$ and no
factor $\log p$ arises.  We adopt instead the adelic normalization, in which
the valuation map pushes $d^\times u$ to the measure assigning $\log p$ to
each step of the value group $p^{\Z}$ --- the normalization making the finite
places commensurate with the archimedean one, where the value group is
$\R_{>0}$ with $dt/t$, and the one in which Weil's local terms are stated.
Then $\Gamma^{\mathrm{Tate}}_{p,k}=(\log p)\Gamma_{p,k}$.

\begin{theorem}\label{thm:weilcoeff}
At the central weight,
\[
 \Gamma^{\mathrm{Tate}}_{p,k}(f_{1/2})=(\log p)\,p^{-k/2}
 =\frac{\Lambda(p^{k})}{\sqrt{p^{k}}},
\]
and the mirror shells give the same value.  Thus the Weil local term
\[
 W_p(h)=(\log p)\sum_{k\ge1}p^{-k/2}
 \bigl[h(p^{k})+h(p^{-k})\bigr]
\]
is recovered term by term.
\end{theorem}

The coincidence of the two shell families at $s=\tfrac12$ is the fixed-point
property of the mirror involution of Lemma \ref{lem:mirror}.

\subsection{The assembly}

\begin{theorem}\label{thm:assembly}
For $\Re s>1$,
\[
 \sum_p\sum_{k\ge1}\Gamma^{\mathrm{Tate}}_{p,k}(f_s)
 =\sum_{n\ge1}\Lambda(n)n^{-s}
 =-\frac{\zeta'}{\zeta}(s),
\]
absolutely convergent, and the mirror family gives $-\zeta'/\zeta(1-s)$ for
$\Re s<0$.
\end{theorem}

The two half-planes are exchanged by $s\mapsto1-s$ and are disjoint from the
strip; on the strip the assembly exists only as the meromorphic
continuation.  This is the same phenomenon as the divergence of the naive
all-places sum, now at the level of the correspondence functionals.

\begin{remark}
Theorem \ref{thm:assembly} by itself is only an identity between a sum of
functionals and a Dirichlet series.  The next section constructs the global
nuclear diagonal trace which promotes it to row~(c).  The archimedean term is
absent here by design and is supplied there by the Poisson quotient trace.
Finiteness sits on $\Gamma_{p,k}$, whose support is a single point, and
equivariance on the sections it is paired against; the two demands are on
opposite sides of the pairing, which is what makes the construction possible
at all.
\end{remark}


% BEGIN INLINED SOURCE: row-c-nuclear-lefschetz
\section{The nuclear Lefschetz character}
\label{sec:rowc}

Row~(b) produced two related, but differently typed, objects: the
correspondence $\Gamma_n$, which composes multiplicatively, and its dynamic
perfect contact $\mathbb L_n$, whose determinant degree is $\Lambda(n)$.
The latter degree is not a multiplicative character.  Thus a global
Lefschetz formula cannot be obtained by pretending that
$\mathbb L_n\mapsto\Lambda(n)$ is a monoidal functor.  The construction in
this section keeps the two types separate.  It realizes $\Gamma_n$
monoidally as a scaling multiplier $U_n$, and then forms the operator-valued
contact
\[
 \sum_{n\geq2}\deg_{\det}(\mathbb L_n)U_n
 =\sum_{n\geq2}\Lambda(n)U_n.
\]
The equality with the logarithmic Euler character of the Zeta operator is a
theorem below.  Meyer's Poisson construction then supplies a summable
super-representation whose supercharacter adds the opposite finite
orientation, the two polar characters and the archimedean term.  This is
the precise cohomological/nuclear meaning of row~(c).  No Cartier or Chow
intersection on the spherical carrier is asserted.

\subsection{Weighted spaces and the trace category}

Identify $\R_+^\times$ with $\R$ by the logarithm and write
$d^*x=dx/x$.  For an interval $I\subseteq\R$, let
\begin{equation}\label{eq:weightedSchwartz}
 \mathcal S_I=\{f:x^sf(x)\in\mathcal S(\R_+^\times)
                  \text{ for every }s\in I\}.
\end{equation}
For compact $I$ this has the intersection Fr\'echet topology, and in
general the projective-limit topology over a countable compact exhaustion.
These are nuclear Fr\'echet spaces.  Following
\cite[\S2]{MeyerZeta}, put
\begin{equation}\label{eq:Meyerspaces}
 \mathcal H_-:=\mathcal S_{(-\infty,\infty)},\qquad
 \mathcal S_>:=\mathcal S_{(1,\infty)},\qquad
 \mathcal S_<:=\mathcal S_{(-\infty,0)},
\end{equation}
and let $\mathcal H_+$ be the even Schwartz functions on $\R$.

Our regular-representation convention is
\begin{equation}\label{eq:dilationconvention}
 (\lambda_t f)(x)=f(t^{-1}x),\qquad
 \lambda(h)f=\int_0^\infty h(t)\lambda_t(f)d^*t.
\end{equation}
Let $\mathsf{NFRep}$ be the exact category of smooth scaling
representations on nuclear Fr\'echet spaces, with closed invariant
embeddings having nuclear Fr\'echet quotient as the admissible
monomorphisms.  Its $\Z/2$-graded enlargement is denoted
$\mathsf{NFRep}^{\mathrm{sup}}$.  A graded object $E=E^{\bar0}\oplus
E^{\bar1}$ is \emph{summable} when the two integrated operators
$\lambda_E(h)$ are nuclear for every $h\in C_c^\infty(\R_+^\times)$; then
\begin{equation}\label{eq:nuclearcharacter}
 \chi_E(h)=\operatorname{Tr}(\lambda_{E^{\bar0}}(h))
           -\operatorname{Tr}(\lambda_{E^{\bar1}}(h))
\end{equation}
is its nuclear supercharacter.  Closed subspaces and quotients of nuclear
Fr\'echet spaces are nuclear, so the objects below belong to this exact
category.  Meyer's summability theorem
\cite[Proposition~5.6]{MeyerZeta} proves the additional trace-class
assertion required by \eqref{eq:nuclearcharacter}.

Let $\mathsf{Mult}(\mathcal S_>)$ be the one-object strict monoidal
$\C$-linear category whose endomorphism algebra consists of the continuous
convolution multipliers of $\mathcal S_>$.  This is the target of the
individual correspondences.  The supercategory
$\mathsf{NFRep}^{\mathrm{sup}}$, rather than the multiplier category, is the
target of the completed trace object.

\subsection{The typed realization of rows (a) and (b)}

Let $\mathsf{Corr}_{\C}^{\mathrm{arith}}(X)$ be the one-object
$\C$-linear category whose endomorphism algebra is the monoid algebra of
the isomorphism classes in the graph-image category of row~(b).  Its
distinguished basis arrows are $[\Gamma_n]$ and
\begin{equation}\label{eq:linearizedcomposition}
 [\Gamma_m][\Gamma_n]=[\Gamma_{mn}].
\end{equation}
This definition uses the actual composition isomorphism of
Theorem~\ref{thm:rowbcompose}; it is not a new freely chosen correspondence
algebra.  Put
\begin{equation}\label{eq:Un}
 U_n:=\lambda_n^{-1},\qquad (U_nf)(x)=f(nx).
\end{equation}
Since $U_mU_n=U_{mn}$, the formula
\begin{equation}\label{eq:Rfunctor}
 \mathcal R([\Gamma_n])=U_n
\end{equation}
extends uniquely to a strict monoidal $\C$-linear functor
\[
 \mathcal R:\mathsf{Corr}_{\C}^{\mathrm{arith}}(X)
       \longrightarrow\mathsf{Mult}(\mathcal S_>).
\]

This realization is connected to both coefficient systems already
constructed.  The intertwiner of row~(b) sends
$V_n\phi_r$ to $\delta_n*\delta_r=\delta_{nr}$ in the nuclear Dirichlet
algebra $\mathcal C_{\R}$.  Define
\begin{equation}\label{eq:DirichletDistributionMap}
 \mathfrak j:\mathcal C_{\R}\longrightarrow
 \mathcal D'(\R_+^\times),\qquad
 \mathfrak j((a_r))=\sum_{r\geq1}a_r\delta_{r^{-1}}.
\end{equation}
The sum is locally finite as a distribution: every compact subset of
$\R_+^\times$ contains only finitely many $r^{-1}$.  On a bounded set of
test functions the supports lie in one compact set, so the corresponding
strong-dual seminorm is bounded by a finite coordinate seminorm, and
$\mathfrak j$ is continuous.  Dirichlet convolution becomes multiplicative
convolution.  Consequently
\begin{equation}\label{eq:fullintertwiner}
 \mathfrak j(J(V_n\phi_r))
 =\delta_{(nr)^{-1}}
 =\delta_{n^{-1}}*\mathfrak j(J(\phi_r)),
\end{equation}
where $J$ is the Witt--nuclear intertwiner
\eqref{eq:WittNuclearIntertwiner}.  Acting on $\mathcal S_>$, convolution by
$\delta_{n^{-1}}$ is precisely $U_n$.  Thus \eqref{eq:Rfunctor} is the
distributional realization of the adjoint Witt dynamics and of the
$e_\Gamma$ summand of $\mathcal N_{DN}$, not merely a map of labels.

The metric component gives the same label intrinsically:
$\nu_\infty(\mathcal T_n)=\log n$.  Hence either the Witt action, the
nuclear point mass or the metric line recovers $n$, and all three determine
the same multiplier $U_n$ under \eqref{eq:Rfunctor}.

The intrinsic row-(a) cohomology is compatible with this realization, with
no identification of unlike types.  At every finite periodic depth its
nuclear frame is
\[
 T_{\mathcal C}(\mathbf H_{r,s}^{\rm int}(x))
 =\mathcal O_{\mathcal C}\otimes_\R V_{r,s}^{\rm per}(x).
\]
Left convolution by $\delta_n$ on the scalar factor commutes with principal
translations, finite-support sums and the Cartesian-basis K\"unneth maps,
because $\mathcal C_{\R}$ is commutative and those maps are
$\mathcal O_{\mathcal C}$-linear.  Under $\mathfrak j$ this scalar action is
convolution by $\delta_{n^{-1}}$, hence the multiplier $U_n$.  Moreover,
for the dynamic contact of row~(b),
\begin{equation}\label{eq:ABCcompatibility}
 \ell(\delta_n)
 =-\log\|\det_{\rm tor}\mathbb L_n\|
 =\Lambda(n).
\end{equation}
Thus rows (a), (b) and (c) share the same nuclear scalar action and the same
derived contact determinant.  The statement is a commuting comparison on
the common nuclear component; it does not falsely promote the completed
nuclear representation to a perfect periodic object.

For later convergence, choose the usual defining seminorms on
$\mathcal S_a$ involving $(\log x)^j(x\partial_x)^kf$.  Since $x\partial_x$
commutes with $U_n$ while
$\log x=\log(nx)-\log n$, for every such seminorm there are $C,N$ with
\begin{equation}\label{eq:Unestimate}
 q_{a,j,k}(U_nf)\leq Cn^{-a}(1+\log n)^N
   \sum_{j'\leq j}q_{a,j',k}(f).
\end{equation}
Thus $\sum_n\Lambda(n)U_n$ converges absolutely in every defining
multiplier seminorm after choosing any $a>1$; polynomial logarithmic factors
do not affect convergence.

\subsection{Dynamic contact and the logarithmic Euler character}

Meyer's Zeta operator is
\begin{equation}\label{eq:MeyerZ}
 Z=\sum_{n\geq1}U_n,\qquad Zf(x)=\sum_{n\geq1}f(nx).
\end{equation}
The coefficient sequence $(1,1,\ldots)$ is not in the rapidly decreasing
algebra $\mathcal C_{\R}$.  Accordingly \eqref{eq:MeyerZ} is not presented
as its image.  Rather, the finite sums
$\mathfrak j(\sum_{n\leq N}\delta_n)$ converge in
$\mathcal D'(\R_+^\times)$ to the locally finite Dirac comb
$\check\zeta=\sum_n\delta_{n^{-1}}$: on every bounded set of test functions
the pairing is eventually constant outside one finite set of indices.
The operator $Z$ is the continuous action of this distribution on
$\mathcal S_>$.
Proposition~3.2 of \cite{MeyerZeta} proves that $Z$ and
$Z^{-1}=\sum_n\mu(n)U_n$ are mutually inverse continuous multipliers of
$\mathcal S_>$, and proves the convergent Euler product
\begin{equation}\label{eq:ZEuler}
 Z=\prod_p(1-U_p)^{-1}.
\end{equation}
No invertibility on $\mathcal H_+$ or on the quotient representations below
is asserted.

For functions or distributions on $\R_+^\times$ put
\begin{equation}\label{eq:partialconvention}
 (\partial f)(x)=(\log x)f(x).
\end{equation}
It is a continuous derivation for multiplicative convolution.  Since
$U_n$ is represented by $\delta_{n^{-1}}$,
\begin{equation}\label{eq:partialsign}
 \partial U_n=-(\log n)U_n.
\end{equation}
The estimates \eqref{eq:Unestimate} justify differentiating
\eqref{eq:ZEuler} in the topology of multipliers of $\mathcal S_>$.  The
derivation rule and \eqref{eq:partialsign} give
\begin{equation}\label{eq:contactEulercharacter}
 Z\partial(Z^{-1})
 =\sum_p(1-U_p)^{-1}(\log p)U_p
 =\sum_p\sum_{k\geq1}(\log p)U_{p^k}
 =\sum_{n\geq2}\Lambda(n)U_n.
\end{equation}

This identity is exactly compatible with the dynamic contact of row~(b).
Indeed, Theorem~\ref{thm:rowbcontact} gives, before any analytic
realization,
\begin{equation}\label{eq:rowccontactcomparison}
 \deg_{\det}(\mathbb L_n)
 :=-\log\|\det_{\mathrm{tor}}\mathbb L_n\|=\Lambda(n).
\end{equation}
Combining \eqref{eq:Rfunctor}, \eqref{eq:contactEulercharacter} and
\eqref{eq:rowccontactcomparison} yields the equality of convergent
operator-valued distributions
\begin{equation}\label{eq:realizedcontact}
 Z\partial(Z^{-1})
 =\sum_{n\geq2}\deg_{\det}(\mathbb L_n)\,
       \mathcal R([\Gamma_n]).
\end{equation}
This is a labelwise determinant realization, not the false assertion that
determinant degree is a monoidal character.

For a convolution kernel put $\tau(f)=f(1)$.  If
$h\in C_c^\infty(\R_+^\times)$, apply the continuous functional
$A\mapsto\tau(h*A)$ to \eqref{eq:realizedcontact}.  Since
$\tau(h*U_n)=h(n)$, this produces
\begin{equation}\label{eq:Wfinplus}
 W_{\mathrm{fin}}^+(h)=
 \sum_p\sum_{k\geq1}(\log p)h(p^k).
\end{equation}
Let
\begin{equation}\label{eq:Janalytic}
 (\mathscr Jh)(x)=x^{-1}h(x^{-1}).
\end{equation}
Conjugating by this involution gives the opposite orientation
\begin{equation}\label{eq:Wfinminus}
 W_{\mathrm{fin}}^-(h)=
 \sum_p\sum_{k\geq1}(\log p)p^{-k}h(p^{-k}).
\end{equation}
Both sums are finite for compactly supported $h$, because such a support is
bounded away from $0$ and $\infty$.

\subsection{The source-defined Poisson quotient}

Use the additive Fourier transform
\begin{equation}\label{eq:additiveFourier}
 (\mathcal Ff)(y)=\int_\R f(x)e^{2\pi ixy}\,dx,
\end{equation}
fixed by the standard self-dual additive character, and set
\begin{equation}\label{eq:Hcap}
 \mathcal H_\cap=
 \{f\in\mathcal H_+:f(0)=\mathcal Ff(0)=0\}.
\end{equation}
Poisson summation gives $Zf=\mathscr JZ\mathcal Ff$ on
$\mathcal H_\cap$.  Meyer constructs a scaling representation
$\mathcal H_\cup$ fitting into the exact sequence
\begin{equation}\label{eq:poleextension}
 0\longrightarrow\mathcal H_-
 \longrightarrow\mathcal H_\cup
 \longrightarrow\C(x^0)\oplus\C(x^1)\longrightarrow0.
\end{equation}
His Theorem~3.3 proves that
$Z:\mathcal H_+\to\mathcal H_\cup$ is continuous, has closed range and is
a topological isomorphism onto that range, and that
$Zf\in\mathcal H_-$ exactly for $f\in\mathcal H_\cap$.  Therefore the
closed quotients
\begin{equation}\label{eq:Poissonquotients}
 \mathcal H_+^0:=\mathcal H_+/\mathcal H_\cap
 \simeq\C(x^0)\oplus\C(x^1),
 \qquad
 \mathcal H_-^0:=\mathcal H_-/Z\mathcal H_\cap
\end{equation}
are nuclear Fr\'echet scaling representations.  Define the actual
superobject
\begin{equation}\label{eq:Lcomp}
 \mathbb L_{\mathrm{comp}}^{\bar0}=\mathcal H_+^0,\qquad
 \mathbb L_{\mathrm{comp}}^{\bar1}=\mathcal H_-^0
 \quad\text{in }\mathsf{NFRep}^{\mathrm{sup}}.
\end{equation}
This replaces the potentially ambiguous phrase ``virtual object'' by a
typed $\Z/2$-graded representation.  This supercategorical packaging is our
definition; the quotient representations and their summability are
Meyer's.  Proposition~5.6 of
\cite{MeyerZeta} proves that it is summable.  Its even part is exactly the
two-dimensional polar object; its odd part is the closed quotient
$\mathcal H_-^0$.
All ingredients in \eqref{eq:Lcomp} are defined from the Dirac comb, the
additive Fourier transform and closed subquotients, before any spectral
description is taken.

The archimedean distribution must retain Meyer's normalization.  For a
compactly supported function $h$ on $\R_+^\times$, extend it by zero at the
origin and let $h_{\mathrm{ev}}(x)=h(|x|)$.  Define
\begin{equation}\label{eq:Winftyprecise}
 W_\infty(h):=-\left\langle\mathcal F(\log|x|),
                 y\longmapsto h_{\mathrm{ev}}(1-y)\right\rangle.
\end{equation}
Equivalently,
\begin{equation}\label{eq:Winftyfinitepart}
 W_\infty(h)=\operatorname{fp}_{\mathcal F}\int_0^\infty
 \left(\frac{h(x)}{|1-x|}+\frac{h(x)}{1+x}\right)dx.
\end{equation}
Here $\operatorname{fp}_{\mathcal F}$ is the Fourier finite-part extension,
not a Cauchy principal value of an even singularity.  Explicitly it is the
limit after subtracting $2h(1)\log(1/\varepsilon)$ from a symmetric cutoff,
plus $c_{\mathcal F}h(1)$, where the unique constant $c_{\mathcal F}$ is
fixed by
$\mathcal F^2(\log|x|)(1)=0$.  This is precisely the normalization in the
proof of \cite[Theorem~5.8]{MeyerZeta}.  The measure in
\eqref{eq:Winftyfinitepart} is $dx$, not $d^*x$.

\subsection{The completed diagonal character}

\begin{definition}\label{def:nuclearintersection}
For $h\in C_c^\infty(\R_+^\times)$, the \emph{nuclear cohomological
diagonal trace} of the completed correspondence family is
\begin{equation}\label{eq:nuclearintersection}
 I_\Delta^{\mathrm{nuc}}(h)
 :=\chi_{\mathbb L_{\mathrm{comp}}}(h).
\end{equation}
This is the supertrace of the source-defined object
\eqref{eq:Lcomp}; it is not defined by the explicit formula.
\end{definition}

\begin{theorem}[Arithmetic nuclear Lefschetz identity]
\label{thm:rowc}
For every $h\in C_c^\infty(\R_+^\times)$,
\begin{equation}\label{eq:rowcidentity}
 \boxed{
 I_\Delta^{\mathrm{nuc}}(h)=
 \sum_p\sum_{k\geq1}\log p\,
 \bigl(h(p^k)+p^{-k}h(p^{-k})\bigr)+W_\infty(h).}
\end{equation}
The first finite orientation is
$\sum\deg_{\det}(\mathbb L_n)\mathcal R([\Gamma_n])$;
the second is its $\mathscr J$-conjugate; the polar object is
$\C(x^0)\oplus\C(x^1)$; and the last term is the Fourier--archimedean trace
\eqref{eq:Winftyprecise}.  The same character has the spectral expression
\begin{equation}\label{eq:rowcspectral}
 I_\Delta^{\mathrm{nuc}}(h)=\widehat h(0)+\widehat h(1)
 -\sum_\rho m_\rho\widehat h(\rho),\qquad
 \widehat h(s)=\int_0^\infty h(x)x^s d^*x,
\end{equation}
where the sum is understood as the nuclear character distribution.  No zero
in \eqref{eq:rowcspectral} enters any definition.
\end{theorem}

\begin{proof}
The contact computation before completion is
\eqref{eq:realizedcontact}; it identifies the value of the positive finite
trace expression in Meyer's calculation with \eqref{eq:Wfinplus}.
Conjugation by
$\mathscr J$ proves \eqref{eq:Wfinminus}.  Proposition~5.6 of
\cite{MeyerZeta} proves that the induced operators on the two closed
quotients \eqref{eq:Poissonquotients} are nuclear and that their trace is
the difference of the corresponding ambient commutator traces.  The third
commutator is the Fourier commutator; the calculation in
\cite[Theorem~5.8 and its proof]{MeyerZeta} identifies its trace with
\eqref{eq:Winftyprecise}.  Adding these three proved traces gives
\eqref{eq:rowcidentity}, with the displayed signs.  Compact support handles
the finite sums and the Fourier finite part handles the unique
archimedean singularity.

For the second computation, Theorem~4.1 and Corollary~4.2 of
\cite{MeyerZeta} identify the transpose spectrum of $\mathcal H_-^0$, with
algebraic multiplicity, with the nontrivial zeros of the completed zeta
function.  The isomorphism in \eqref{eq:Poissonquotients} gives the even
eigencharacters $0$ and $1$.  The nuclear trace theorem, applied in
Meyer's Theorem~5.8, therefore gives \eqref{eq:rowcspectral}: the even polar
characters occur with plus sign and the odd quotient representation with
minus sign.  The construction order is
\[
 \{\mathfrak j(\delta_n)\}_{n\geq1}
 \rightsquigarrow\check\zeta=\sum_{n\geq1}\delta_{n^{-1}},\quad
 Z,\quad \text{Poisson summation},\quad
 \mathcal H_\pm^0,\quad\chi_{\mathbb L_{\mathrm{comp}}};
\]
only the subsequent spectral description introduces the zeros.
\end{proof}

\begin{corollary}[Self-dual metric centralization]
\label{cor:rowccentral}
For $g\in C_c^\infty(\R_+^\times)$ put $h(x)=x^{-1/2}g(x)$.  The finite
part of \eqref{eq:rowcidentity} becomes
\[
 \sum_p\sum_{k\geq1}\frac{\Lambda(p^k)}{\sqrt{p^k}}
 \bigl(g(p^k)+g(p^{-k})\bigr).
\]
For each orientation the coefficient is exactly the unique positive
self-dual metric normalization \eqref{eq:centralLocalCoefficient} of
row~(b).  No algebraic half-Tate line is used.
\end{corollary}

\begin{corollary}[Constructed row (c)]\label{cor:rowccomplete}
Row~(c) is constructed completely under the following explicit contract:
the source is the monoid linearization
$\mathsf{Corr}_{\C}^{\mathrm{arith}}(X)$ of the isomorphism classes in
$\mathsf{Corr}_{W,DN}^{\mathrm{arith}}(X)$; individual correspondences are
realized monoidally in $\mathsf{Mult}(\mathcal S_>)$; their contact weights
are the determinant degrees of the dynamic complexes $\mathbb L_n$; and
the completed diagonal trace is the nuclear supercharacter of
$\mathbb L_{\mathrm{comp}}$ in $\mathsf{NFRep}^{\mathrm{sup}}$.  Within
this cohomological/nuclear category, composition, contact, convergence,
both finite orientations, the polar object, the archimedean finite part and
the two trace expressions are all constructed and proved compatible.
This statement neither requires nor asserts a reverse functor to perfect
Cartier complexes on $\mathscr Y_{\mathbb S}$ or an ordinary Chow
intersection product there.
\end{corollary}

\begin{proof}
The monoidal realization and its compatibility with rows~(a) and~(b) are
\eqref{eq:Rfunctor}--\eqref{eq:ABCcompatibility}.  The dynamic contact
comparison is \eqref{eq:realizedcontact}, including its convergence proof.
The source-defined summable superobject is \eqref{eq:Lcomp}.
Theorem~\ref{thm:rowc} computes its character geometrically and spectrally,
and Corollary~\ref{cor:rowccentral} proves compatibility with the metric
normalization of row~(b).  These are exactly the clauses of the stated
contract.
\end{proof}

% END INLINED SOURCE: row-c-nuclear-lefschetz


\section{Row (d): the exact geometric gate}
\label{sec:rowdgate}

We now determine whether the numerical ruling signature already proved on
$N^1_{\mathrm{rul}}\simeq\mathbb R^2$ can control the nuclear correspondence
pairing.  It cannot: the two forms live on numerical objects of different
rank.  This obstruction follows directly from the prime contacts, without
using the spectral expression \eqref{eq:rowcspectral} or any zero of $\xi$.

\subsection{The form and the primitive target}

Let $\mathcal T=C_c^\infty(\mathbb R_+^\times)$ and define
\begin{equation}\label{eq:Bnuc}
 B_{\mathrm{nuc}}(f,g)=I_\Delta^{\mathrm{nuc}}(f\star g^\vee).
\end{equation}
The Tate involution \eqref{eq:Tateinvolution} and
Theorem~\ref{thm:rowc} make this a Hermitian form.  Its two-ruling primitive
subspace is
\begin{equation}\label{eq:twoprimitive}
 \mathcal T^0=\{f\in\mathcal T:
                  \widehat f(0)=\widehat f(1)=0\}.
\end{equation}
In a genuine intersection category containing ruling classes $F_1,F_2$ with
$F_i^2=0$ and $F_1F_2=1$, the corresponding primitive class would be
\begin{equation}\label{eq:requiredprimitiveclass}
 M_f=D_f-\widehat f(1)F_1-\widehat f(0)F_2,
 \qquad M_fF_1=M_fF_2=0.
\end{equation}
This is the specification of a missing comparison, not a claim that $D_f$
is already a divisor-line class on $\mathscr Y_{\mathbb S}$.

\subsection{An infinite-rank obstruction independent of the zeros}

Use the logarithmic coordinate $t=\log x$, and let $W$ represent the
functional $I_\Delta^{\mathrm{nuc}}$.  By \eqref{eq:rowcidentity}, its finite
part has nonzero atoms at
\begin{equation}\label{eq:primeatoms}
 t=k\log p\quad\hbox{and}\quad t=-k\log p
 \qquad(p\text{ prime},\ k\ge1).
\end{equation}
The Fourier--archimedean finite part is smooth away from $t=0$, so it cannot
cancel these atoms.

\begin{theorem}[Infinite-rank contact obstruction]
\label{thm:infiniterankcontact}
The convolution map
\[
 T_W:C_c^\infty(\mathbb R)\longrightarrow\mathcal D'(\mathbb R),
 \qquad a\longmapsto W*a,
\]
has infinite algebraic rank.  Consequently \eqref{eq:Bnuc} cannot factor
through a finite-dimensional vector space, and in particular cannot factor
through $N^1_{\mathrm{rul}}$.
\end{theorem}

\begin{proof}
Suppose the image lies in a finite-dimensional space $E$.  If $(\eta_j)$ is
a compactly supported approximate identity, then $W*\eta_j\to W$ in
$\mathcal D'(\mathbb R)$.  Finite-dimensional subspaces of a Hausdorff
topological vector space are closed, so $W\in E$.  Applying the same argument
to translates of $\eta_j$ shows that all translates of $W$ lie in $E$.
Their span $E_W$ is therefore finite-dimensional and translation invariant.

Difference quotients converge distributionally, hence
$D(E_W)\subset E_W$ for $D=d/dt$.  Cayley--Hamilton gives a nonzero
polynomial $P$ with $P(D)W=0$.  Factoring $P$ over $\mathbb C$ and solving
the successive equations $(D-\lambda)^m u=0$ shows that every global
distributional solution is an exponential polynomial
$e^{\lambda t}Q(t)$, and therefore smooth.  This contradicts the nonzero
point masses \eqref{eq:primeatoms}.  Thus $T_W$ has infinite rank.

If $B_{\mathrm{nuc}}(f,g)=q(Jf,Jg)$ through a finite-dimensional $E_0$,
then $f\mapsto B_{\mathrm{nuc}}(f,\cdot)$ has rank at most $\dim E_0$.
The Tate involution is bijective, so this is the rank of $T_W$, a
contradiction.  Equation \eqref{eq:rulingquotient} now gives the last assertion.
\end{proof}

The hyperbolic signature calculation on the ruling plane is exact but cannot
be reused as a Hodge theorem for row~(d).  Forgetting the Witt map, nuclear
action, metric frame and derived orbit contact makes all $\Gamma_n$ the
same diagonal, whereas retaining their completed contact form has infinite
rank.

\begin{corollary}[Primitive infinite rank]
\label{cor:primitiveinfiniterank}
The restriction of $T_W$ to $\mathcal T^0$ in
\eqref{eq:twoprimitive} still has infinite algebraic rank.
\end{corollary}

\begin{proof}
The two Mellin evaluations defining $\mathcal T^0$ give it codimension at
most two in $\mathcal T$.  If the restricted map had finite rank, adjoining
an algebraic complement of dimension at most two would make $T_W$ finite
rank on all of $\mathcal T$, contrary to
Theorem~\ref{thm:infiniterankcontact}.
\end{proof}

\subsection{The finite-support construction and its obstruction}

There is a natural attempt to enlarge the ruling plane without using the
explicit formula.  Let $S$ be a finite set of prime powers and put
\begin{equation}\label{eq:finitecandidate}
 E_S=\mathbb RF_1\oplus\mathbb RF_2
     \oplus\bigoplus_{n\in S}\mathbb Re_n,
 \qquad H_S=F_1+F_2,
\end{equation}
where $e_n$ records the decorated correspondence $\Gamma_n$.  For
$S\subset T$ the labelled correspondences give the evident inclusion
$E_S\hookrightarrow E_T$.

The intrinsic periodic ruling form \eqref{eq:intrinsicRRform} factors through the two degrees and has
rank two.  Assigning the new symbols zero ruling degree makes all $e_n$
radical; forgetting their decorations makes all $\Gamma_n$ the same
diagonal.  Either operation loses the mixed sector.

Composition and reduced contact give another form which does retain that
sector:
\begin{equation}\label{eq:finitecontactform}
 K_S(e_m,e_n)
 =\deg_{\mathrm{cont}}(\Gamma_m\circ\Gamma_n,\Delta)
 =\Lambda(mn).
\end{equation}
This is defined before the nuclear completion, using only
Theorems~\ref{thm:rowbcompose} and \ref{thm:rowbcontact}.  Its sign is fixed
by $-\log\|\det_{\mathrm{tor}}\mathbb F_p\|=\log p>0$.

\begin{theorem}[Finite-contact truncation obstruction]
\label{thm:finitecontactobstruction}
Suppose that $S$ contains powers of $r$ distinct primes.  On the
correspondence summand of \eqref{eq:finitecandidate}, $K_S$ is the orthogonal
sum
\[
 \bigoplus_{p\mid S}(\log p)\,\mathbf1_p\mathbf1_p^t
\]
and has inertia $(r,0,|S|-r)$.  With the ruling plane adjoined, the inertia is
$(1+r,1,|S|-r)$.  Restriction by any two ruling-degree equations leaves at
least $r-2$ positive correspondence directions.  The transition inclusions
preserve the forms, so the positive index is unbounded on a cofinal family.
\end{theorem}

\begin{proof}
If $m,n$ are powers of the same prime $p$, then
$\Lambda(mn)=\log p$; if they are supported at distinct primes, it is zero.
This proves the displayed block decomposition.  Each all-one block has one
positive eigenvalue and all remaining eigenvalues zero.  The ruling matrix
$\left(\begin{smallmatrix}0&1\\1&0\end{smallmatrix}\right)$ adds one positive
and one negative eigenvalue.  Finally, a positive $r$-plane meets a subspace
of codimension at most two in dimension at least $r-2$, and remains positive
there.  Formula \eqref{eq:finitecontactform} is unchanged when $S$ is enlarged.
\end{proof}

Thus the two available finite constructions fail in opposite ways:
\begin{equation}\label{eq:finitealternative}
 \text{ruling RR}\Longrightarrow\text{mixed radical},
 \qquad
 \text{contact}\Longrightarrow\text{one positive line per prime}.
\end{equation}
Replacing $K_S$ by $-K_S$ would reverse the already fixed torsion determinant
and fail comparison with row~(c).  What is missing is a global Green term
which couples distinct primes while preserving their local contact, together
with sections, effectivity and a quadratic RR theorem for the $e_n$.  None of
those structures is present on \eqref{eq:finitecandidate}.

\subsection{The Green correction is forced}

The finite part of Theorem~\ref{thm:rowc} defines, on the test algebra,
\begin{equation}\label{eq:Ktest}
 K(f,g)=W_{\mathrm{fin}}(f\star g^\vee),\qquad
 W_{\mathrm{fin}}(h)=\sum_{p,k\ge1}\log p
 \bigl(h(p^k)+p^{-k}h(p^{-k})\bigr).
\end{equation}
This is the completed version of the fixed contact determinants.  There is
no freedom in the complementary term.

\begin{theorem}[Forced Green correction]
\label{thm:forcedgreen}
Let $Q$ be a Hermitian form on the test algebra whose finite contact factor
is \eqref{eq:Ktest} and whose final comparison with row~(c) is exact,
$Q=B_{\mathrm{nuc}}$.  Then
\begin{equation}\label{eq:forcedgreen}
 G:=Q-K=G_\infty,\qquad
 G_\infty(f,g)=W_\infty(f\star g^\vee).
\end{equation}
No different Green correction preserves both the dynamic perfect contacts and the
nuclear Lefschetz identity.
\end{theorem}

\begin{proof}
The cohomological character identity \eqref{eq:rowcidentity} says
$B_{\mathrm{nuc}}=K+G_\infty$.  Subtracting the prescribed $K$ from
$Q=B_{\mathrm{nuc}}$ gives \eqref{eq:forcedgreen}; two such corrections have
difference zero.
\end{proof}

The forced term has a genuine, but insufficient, energy property.  Put
$F(t)=e^{t/2}f(e^t)$.  Mellin transformation of the Fourier finite part
\eqref{eq:Winftyprecise} gives
\begin{equation}\label{eq:archmultiplier}
 G_\infty(f,f)=\frac1{2\pi}\int_{\mathbb R}m_\infty(\tau)
 |\widehat f(\tfrac12+i\tau)|^2\,d\tau,\qquad
 m_\infty(\tau)=\log\pi-\Re\psi(\tfrac14+i\tfrac\tau2).
\end{equation}
Here $m_\infty=-2\Re(\Gamma_{\mathbb R}'/\Gamma_{\mathbb R})$ on the
critical line.  The digamma series gives
\begin{equation}\label{eq:archenergy}
 m_\infty(\tau)-m_\infty(0)
 =-\sum_{n=0}^\infty
 \frac{(\tau/2)^2}
 {(n+1/4)((n+1/4)^2+(\tau/2)^2)}\le0,
\end{equation}
strictly away from zero.  On the other hand
$m_\infty(0)=\log\pi-\psi(1/4)>0$, while the asymptotic
$\psi(z)=\log z+O(z^{-1})$ gives $m_\infty(\tau)\to-\infty$.
Thus $G_\infty$ is indefinite, although subtracting its maximal constant
multiplier gives a negative semidefinite translation energy.

The desired assertion is not the conditional inequality
\eqref{eq:archenergy}, but the sharp domination
\begin{equation}\label{eq:sharpdomination}
 (K+G_\infty)(f,f)\le0\qquad(f\in\mathcal T^0).
\end{equation}
By \eqref{eq:Dcircularity}, \eqref{eq:sharpdomination} is equivalent to
$\RH$.  Theorem~\ref{thm:forcedgreen} prevents making the archimedean term
more negative without changing row~(c).  Hence the local energy does not
close row~(d); a non-circular proof must derive the exact domination from a
new mixed Riemann--Roch/effectivity theory.

\subsection{The mixed-section forcing theorem}

We finally test whether a section functor can be extracted from the objects
already present.  The classical index argument uses only the following
numerical consequences of such a functor.  On a mixed divisor space with
form $q$, polarization $H$ and $\deg D=q(D,H)$, suppose there are finite
nonnegative, linear-equivalence-invariant quantities $h^i$ such that
\begin{enumerate}
\item $h^0(0)=0$ and $h^0(D)>0$ is equivalent to strict effectivity;
\item every nonzero strictly effective class has positive degree;
\item $h^2(D)=h^0(-D)$;
\item for every finite-support class,
\begin{equation}\label{eq:mixedRRminimum}
 h^0(nD)-h^1(nD)+h^2(nD)
 =\frac{n^2}{2}q(D,D)+O_D(n).
\end{equation}
\end{enumerate}
These statements must be compatible with finite-support inclusions and the
chosen completion.  They are weaker than asking for actual section objects.

\begin{theorem}[Mixed-section forcing]
\label{thm:mixedsectionforcing}
Every package satisfying the preceding four conditions has
\[
 q(D,D)\le0\qquad\text{whenever }q(D,H)=0.
\]
Consequently, if its independently constructed form is subsequently
identified with $B_{\mathrm{nuc}}$, its existence implies $\RH$.
\end{theorem}

\begin{proof}
If $D$ is primitive and $q(D,D)>0$, then
\eqref{eq:mixedRRminimum} is positive for all large $n$.  Nonnegativity of
$h^1$ and duality give
\[
 0<h^0(nD)-h^1(nD)+h^0(-nD)
 \le h^0(nD)+h^0(-nD).
\]
Hence $nD$ or $-nD$ is nonzero and strictly effective, and so has positive
degree.  Both have degree zero, a contradiction.  The last assertion is
\eqref{eq:Dcircularity}.
\end{proof}

The theorem is logically sharp at the scalar level.  If $\RH$ is assumed,
the numerical form is Lorentzian; orient its positive cone by $\deg>0$, put
$\chi(D)=q(D,D)/2$, assign $h^0_{\mathrm{num}}(D)=\chi(D)$ in the future
timelike cone, a positive constant on its nonzero lightlike boundary, and
zero elsewhere, then set
\[
 h^2_{\mathrm{num}}(D)=h^0_{\mathrm{num}}(-D),\qquad
 h^1_{\mathrm{num}}=h^0_{\mathrm{num}}+h^2_{\mathrm{num}}-\chi.
\]
These scalar functions satisfy the displayed axioms.  They are discontinuous
at the light cone and are not dimensions of objects; this converse only
shows that the numerical requirement is equivalent in strength to $\RH$.

None of the natural categorical candidates supplies the missing functor.
The Zeta operator on $\mathcal S_>$ is invertible, so its kernel and cokernel
vanish.  The positive Poisson quotient contains only the two polar
characters, while the negative quotient has the zero divisor as transpose
spectrum; using its eigenspaces makes effectivity spectral.  Finite contact
homology gives $\Lambda(n)$ and the multi-positive blocks of
Theorem~\ref{thm:finitecontactobstruction}, not quadratic mixed sections.
Finally, a nuclear trace or determinant is a linear or multiplicative
spectral invariant, not a nonnegative section dimension with an effective
cone; selecting its positive spectral part is precisely the circular step
in \eqref{eq:Dcircularity}.  Thus the Poisson quotient does not currently define
$h^0$ for mixed classes.

\subsection{Audit of the two possible Hodge routes}

The arithmetic Hodge-index theorem of Faltings--Hriljac--Moriwaki and its
adelic extension by Yuan and Zhang \cite[Theorems~1.1 and~1.3]{YuanZhang}
apply to Hermitian or integrable adelic line bundles on regular proper
arithmetic varieties or normal projective varieties.  Their hypotheses
include a generic degree-zero condition, a fixed ample or nef-and-big
polarization, admissible or integrable metrics, finite arithmetic
intersections and boundedness conditions for the equality statement.

Those hypotheses have not been constructed for the mixed classes $D_f$.
The spherical ringed square is not a regular proper arithmetic variety in
the sense assumed by those theorems, and a virtual nuclear Fr\'echet
representation is not an adelic line bundle.  Moreover the cotangent
cohomology of Section~\ref{sec:geometric-a} is defined on the effective
external ruling sector; $N^1_{\mathrm{rul}}$ cannot contain the completed
mixed classes by Theorem~\ref{thm:infiniterankcontact}; the nuclear character
gives a trace but no effective cone or Riemann--Roch theorem for the $D_f$;
and no comparison sends these classes to arithmetic line bundles while
preserving degree, composition, intersection and the radical.
Therefore the published arithmetic index theorems cannot yet be applied.

The intrinsic Rosati candidate in the nuclear category is precisely
\eqref{eq:Bnuc}.  Duality proves Hermitian symmetry but not a sign.  Using
the separate spectral computation \eqref{eq:rowcspectral}, the block
calculation of \S3.5 gives
\begin{equation}\label{eq:Dcircularity}
 B_{\mathrm{nuc}}(f,f)\le0\quad(f\in\mathcal T^0)
 \qquad\Longleftrightarrow\qquad \RH.
\end{equation}
An on-line zero contributes a negative block and an off-line mirror pair an
additional hyperbolic plane.  Thus \eqref{eq:Dcircularity} is an acceptance
test, not a proof: selecting a negative spectral part assumes the assertion
row~(d) is meant to prove.

There is a genuine finite analytic result on this side, and it has a long
history which we state before our own contribution to it.  The localization
of the Weil form was introduced by Yoshida \cite{Yoshida1992}, who proved
that $\RH$ is equivalent to positive definiteness of $Q_W$ on
$C_c^\infty(-a,a)$ for every $a>0$, and that this positivity \emph{does}
hold for all sufficiently small $a$ \cite[Lemma~2]{Yoshida1992}.  Both the
localized criterion and the existence of an initial positive window are
therefore due to him.  Suzuki
\cite[Theorems~1.1--1.4]{SuzukiScrew} realizes each localized Weil form by
the Friedrichs extension of a screw-kernel operator, proves continuity of
its lowest eigenvalue, and reproves positivity, simplicity and evenness for
sufficiently small windows.  Positivity or nondegeneracy for every window is
equivalent to $\RH$.  This does not supply the missing geometric theorem:
the localized operator is constructed from the Weil quadratic form itself,
so it is a truncation of row~(c), not an RR/effectivity construction
independent of it.  The exact analytic gate is extension from the proven
small-window range to all windows without importing the spectral sign.


% BEGIN INLINED SOURCE: row-d-local-analysis
\subsection{The primitive pullback, a certified initial interval, and the next endpoint audit}
\label{subsec:primitiveendpoints}

We record the part of the analytic gate which can be closed without a
global sign assumption.  In logarithmic coordinates put
\(I_T=(-T,T)\), extend functions by zero outside \(I_T\), and write
\[
 M_-F=\int_{-T}^{T}F(t)e^{-t/2}\,dt,
 \qquad
 M_+F=\int_{-T}^{T}F(t)e^{t/2}\,dt .
\]
These are exactly the pullbacks of the two polar Mellin evaluations in
\eqref{eq:twoprimitive}: indeed
\(\widehat f(s)=\int_{-T}^TF(t)e^{(s-1/2)t}\,dt\), so evaluation at
\(s=0,1\) gives \(M_-,M_+\).  Let
\(\mathcal P_T=\ker M_-\cap\ker M_+\).  Under
\(F(t)=e^{t/2}f(e^t)\), the Friedrichs operator representing
\(-B_{\rm nuc}\) on compact support has the form
\begin{equation}\label{eq:localizedprimitiveoperator}
 A_T=G_{\Gamma,T}-m_0I-
 \sum_{2\le n<e^{2T}}\frac{\Lambda(n)}{\sqrt n}
       (S_{\log n}+S_{-\log n}),
 \qquad
 m_0=\log\pi+\gamma+\frac{\pi}{2}+3\log2,
\end{equation}
where \(S_aF(t)=\widetilde F(t+a)\) and \(G_{\Gamma,T}\) is the
zero-extension compression of the complete digamma/Gamma multiplier
\[
 g_\Gamma(\tau)=\operatorname {Re}\psi
       \left(\frac14+\frac{i\tau}{2}\right)-\psi\left(\frac14\right).
\]
The strict inequality in the sum is harmless at a threshold because the
new translated supports meet in a null set.  Formula
\eqref{eq:localizedprimitiveoperator} contains every prime power acting on
the support, with the determinant and metric coefficient already fixed in
\eqref{eq:centralLocalCoefficient}; no prime-only truncation is made.

We shall use the following elementary operator form of Feshbach reduction.

\begin{lemma}[Directed Feshbach criterion]\label{lem:directedFeshbach}
Let \(P\) be finite rank, \(Q=1-P\), and let \(A=A^*\).  If
\begin{equation}\label{eq:directedFeshbachHypotheses}
 QAQ\ge\delta Q,\qquad
 (QAP)^*(QAP)\le H,\qquad
 PAP-\delta^{-1}H>0,
\end{equation}
then \(A>0\).
\end{lemma}

\begin{proof}
For \(x\in P\mathcal H\) and \(y\in Q\mathcal H\), the complement bound
and Cauchy's inequality give
\[
 2\operatorname{Re}\langle QAPx,y\rangle+\delta\|y\|^2
 \ge-\delta^{-1}\|QAPx\|^2.
\]
Consequently
\[
 \langle A(x+y),x+y\rangle
 \ge \langle(PAP-\delta^{-1}H)x,x\rangle.
\]
This is strictly positive when \(x\ne0\); when \(x=0\) it is at least
\(\delta\|y\|^2\), and is therefore strictly positive unless \(y=0\).
\end{proof}

\begin{theorem}[Certified primitive initial interval]\label{thm:certifiedendpoints}
The complete operator \eqref{eq:localizedprimitiveoperator} is strictly
positive on \(\mathcal P_T\) for every \(0<T\le\frac12\log3\), and also at
\[
 T=\log2.
\]
These are full Hilbert-space assertions, not signs of finite Galerkin
matrices.
\end{theorem}

\begin{proof}
We give the enclosures and reduction which make the computer-assisted part
reproducible.  All elementary constants, matrix entries, quadrature nodes
and remainders are evaluated as outward-rounded Arb intervals; floating
point is used only to choose bases and fixed congruence matrices.

At \(T=\frac12\log3\), the only non-null finite contact is the shift
\(\log2\).  A shift-compatible three-piece mesh with local Legendre degree
nine has dimension 760.  The first 80 Gamma-resolvent energies are assembled
exactly from their exponential ODE, and the omitted positive tail is bounded
below by its first Robin term.  The even and odd Robin contributions are,
respectively, greater than
\(0.0011586663336\) and \(0.0006199539287\).  Directed congruence gives a
low-block lower bound \(-0.00050\), while exact local Grams and Taylor--Young
off-cell remainders bound the Schur loss by \(3.961\,10^{-6}\).  Hence the
complete even and odd margins are
\[
 0.0006547054607\quad\hbox{and}\quad0.0001159930558.
\]
Lemma~\ref{lem:directedFeshbach} proves the full-space endpoint.  Extension
by zero, proved after the endpoint calculations below, then proves every
\(0<T\le\frac12\log3\).

\begin{remark}[What is, and what is not, new in Theorem
\ref{thm:certifiedendpoints}]
\label{rem:yoshidacalibration}
That \emph{some} initial window is positive is not new: it is
\cite[Lemma~2]{Yoshida1992}, and it is reproved unconditionally, with an
asymptotic expansion of the lowest eigenvalue as \(a\to0^+\), in
\cite[Theorem~1.4]{SuzukiScrew}.  Theorem \ref{thm:certifiedendpoints} is a
\emph{quantitative} refinement of that qualitative statement: it fixes the
explicit endpoint \(T=\log2\), it is a full Hilbert-space assertion rather
than a sign of finite Galerkin matrices, and every constant is an
outward-rounded interval enclosure.  The distinction matters because the
open part of row~(d) is the passage to \emph{all} windows, where no
quantitative initial range, however large, is by itself evidence.
\end{remark}

At \(T=\log2\), split the interval at every translate by \(\log2\) and
\(\log3\), and use local Legendre polynomials of degrees at most nine.
The first 160 positive Gamma resolvents are assembled from their exact
semiseparable exponential kernels.  On one cell the resolvent is obtained
from
\[
 (b^2-\partial_t^2)u=2bf,\qquad
 u'(l)=bu(l),\quad u'(r)=-bu(r),
\]
whose polynomial particular solution is the finite sum
\(2b^{-1}\sum_{j\ge0}b^{-2j}f^{(2j)}\).  The omitted positive Gamma tail is
bounded below by its first Robin resolvent.  The directed enclosures are
\[
\begin{array}{c|c}
\text{quantity}&\text{certified bound}\\ \hline
 QA_TQ&>1.5396358725,Q\\
 \|QA_TP\|^2&<3.65731342\,10^{-4}\\
 \text{even high Schur gap}&>0.4659\\
 \text{odd complement gap}&>0.0497\\
 \langle A_Tv_2,v_2\rangle&>0.0020336972824697\\
 \langle A_Tv_o,v_o\rangle&>8.13856175479\,10^{-6}.
\end{array}
\]
The remaining even direction has Rayleigh value greater than
\(-1.746475\,10^{-6}\), squared residual below
\(4.441\,10^{-24}\), and complement gap \(0.0015\).  Its shorted positive
Gamma-tail capacity is enclosed below by \(1.27084620358308\), whereas the
adverse threshold is below \(1.2111\).  A final application of
Lemma~\ref{lem:directedFeshbach} proves both parity restrictions positive.
This completes the proof of the theorem.
\end{proof}

\paragraph{The audited next endpoint.}
At \(T=\frac12\log5\), use the normalized Legendre space \(V_{200}\), with
the two equations \(M_-=M_+=0\) imposed by an interval two-by-two graph
solve.  A concentration estimate for the Fourier band \([-150,150]\), its
positive Gamma multiplier gap, and the exact Bessel tail of the two Tate
vectors give
\begin{equation}\label{eq:logfivecomplement}
 QA_TQ>0.218Q,\qquad Q=1-P_{V_{200}}.
\end{equation}
Inside \(V_{200}\) perform the exact nested decomposition
\[
 V_{200}=D_5\oplus S_{163}\oplus Y_{30}.
\]
Interval Schur complements give
\[
\begin{array}{c|c|c}
\text{block}&\text{least centre eigenvalue}&
              \text{directed lower enclosure}\\ \hline
Y_{30}&3.27356009&0.9999999999999995-7.08\,10^{-17}\\
S_{163}\text{ Schur}&0.05844291&1-8.08\,10^{-14}\\
D_5\text{ Schur}&3.11556559\,10^{-12}&1-1.72\,10^{-4}.
\end{array}
\]
It remains to control the coupling of the last five-dimensional graph to
the \emph{whole} complement.  For every polynomial graph column the exact
endpoint identity is
\begin{equation}\label{eq:endpointlogidentity}
 A_TF(t)=-\frac12F(t)\log(T^2-t^2)+U_F(t),
\end{equation}
where \(U_F\) is analytic on each of the seven contact cells.  It follows
from
\[
 q(z)=\frac{ze^{z/2}}{2\sinh z},\qquad
 R_1(z)=\operatorname{atanh}(e^{-z/2})+
        \arctan(e^{-z/2})+\tfrac12\log z,
\]
and \(R_1'=-q(z)/z+1/(2z)\).  The Taylor radius is \(\pi>\log5\).
On \(|z|=R=2.5\), the Bernoulli generating function gives
\[
 M_R\le \frac{e^{3R/2}}2
 \left(1+R+2\zeta(2)\frac{(R/\pi)^2}{1-(R/\pi)^2}\right),
\quad
 \|q-q_{<M}\|_{[0,\log5]}
 \le \frac{M_R(\log5/R)^M}{1-\log5/R}.
\]
We take \(M=140\).  Every remaining integral is then a finite polynomial
combination of
\[
\begin{aligned}
 \int_0^1x^m\log x\,dx&=-a^{-2},&
 \int_0^1x^m\log(1-x)\,dx&=-H_a/a,\\
 \int_0^1x^m\log^2x\,dx&=2a^{-3},&
 \int_0^1x^m\log^2(1-x)\,dx&=(H_a^2+H_a^{(2)})/a,\\
 \int_0^1x^m\log x\log(1-x)\,dx
 &=H_a/a^2-(\zeta(2)-H_a^{(2)})/a,
\end{aligned}
\]
where \(a=m+1\).  Thus no numerical quadrature occurs in the contracted
continuum Gram \(H_X\).  The five analytic-tail errors lie between
\(2.14\,10^{-22}\) and \(5.54\,10^{-21}\); propagation of the serialized
graph-coefficient balls in the stable Legendre basis gives the action
errors
\[
 (1.6086834,1.9477040,2.0560371,3.3822491,5.0816822)\times10^{-8}.
\]
Finally the centre eigenvalues of
\(K_{\rm final}-0.218^{-1}H_X\) are
\[
 2.78052720\,10^{-12},\ 5.47152465\,10^{-10},\
 2.87827286\,10^{-7},\ 2.78898083\,10^{-5},\
 1.99028729\,10^{-3}.
\]
They are diagnostics only.  A frozen Cholesky congruence followed by Arb
Gershgorin gives the rigorous lower endpoints
\begin{equation}\label{eq:logfivegershgorin}
 0.9497247499,\quad0.9891485575,\quad0.9741473730,\quad
 0.9360164888,\quad0.7940881021.
\end{equation}
The last inference in the former version of this calculation was not valid.
The projection in \eqref{eq:logfivecomplement} is the complement of the
\emph{whole} space \(V_{200}\), whereas the five-column Gram above is formed
only after eliminating two finite safe blocks.  Relative to the actual
orthogonal decomposition
\[
 \mathcal P_T=D\oplus S\oplus Q,
\]
the missing terms are \(A_{SQ}\) and the corrected cross
\[
 C_D=A_{DQ}-A_{DS}A_{SS}^{-1}A_{SQ}.
\]
Writing \(\kappa=\|A_{QS}A_{SS}^{-1/2}\|^2\), the valid sufficient test is
\[
 \kappa<0.218,
 \qquad
 K_D-(0.218-\kappa)^{-1}C_DC_D^*>0.
\]
Neither inequality follows from the five-column calculation.  Thus
\eqref{eq:logfivecomplement}, the finite Schur enclosures and
\eqref{eq:logfivegershgorin} remain rigorous separate statements, but they
do not yet prove positivity of the complete primitive space at this
endpoint.

The proof also supplies an exact nesting principle.  If \(T<T'\), extension
by zero sends \(\mathcal P_T\) isometrically into \(\mathcal P_{T'}\): both
Tate integrals are unchanged.  The global distributions in
\eqref{eq:localizedprimitiveoperator} give the same quadratic value, because
translations longer than the smaller support have zero correlation.
Consequently positivity at a right endpoint proves every smaller window to
which it applies.  This observation does not turn the two endpoints of
Theorem~\ref{thm:certifiedendpoints} into an unbounded cofinal sequence.

\begin{corollary}[Certified initial interval]\label{cor:certifiedinitialinterval}
The complete primitive operator is strictly positive on \(\mathcal P_T\)
for every
\[
 0<T\le\log2.
\]
\end{corollary}

\begin{proof}
Apply the nesting principle with \(T'=\log2\) and use
Theorem~\ref{thm:certifiedendpoints} at that endpoint.
\end{proof}

\subsection{Threshold capacity and exact return reduction}
\label{subsec:thresholdcapacity}

Put \(T_N=\frac12\log N\) and
\(w_N=\Lambda(N)/\sqrt N\).  At the instant a prime-power contact is born,
its translated supports have null intersection.  Hence
\begin{equation}\label{eq:thresholdbirth}
 H_{N,T_N}=H_{N-1,T_N}.
\end{equation}
At the end of the cell split the enlarged primitive space into the
transported old core and the born two-sided annulus.  If
\[
 H_{N,T_{N+1}}=\begin{pmatrix}A_N&B_N\\B_N^*&D_N\end{pmatrix},
\]
then exact block elimination gives
\begin{equation}\label{eq:thresholdcapacity}
 H_{N,T_{N+1}}\ge0
 \Longleftrightarrow
 \operatorname{Ran}B_N\subseteq\operatorname{Ran}A_N^{1/2},
 \quad D_N-B_N^*A_N^\dagger B_N\ge0.
\end{equation}
Thus global row~(d) is equivalent to a cofinal family of source-defined
shorted-capacity inequalities, together with the initial positive window.

There is a useful exact reduction of the inverse appearing in
\eqref{eq:thresholdcapacity}.  Factor the old positive reference and its
opposite channel as \(R_0=X_0^*X_0\), \(L_0=Y_0^*Y_0\), and put
\[
 T_0=R_0^{\dagger/2}L_0R_0^{\dagger/2},\qquad D_0=I-T_0.
\]
For the harmonic lift \(H=R_0^\dagger X_0^*X_E\) and born reference
\(S_E\), introduce the two operators from the new cell to the normalized
old space
\[
 H_c=R_0^{1/2}HS_E^{\dagger/2},\qquad
 Q_c=R_0^{\dagger/2}(X_0^*X_E-Y_0^*Y_E)S_E^{\dagger/2}.
\]
For a fixed vector \(e\) in the supported part of the new cell, set
\(h_e=H_ce\), \(q_e=Q_ce\).  Direct block algebra gives the normalized
return load
\begin{equation}\label{eq:returnload}
 u_e=D_0h_e-q_e.
\end{equation}
Under the old-core hypothesis \(0\le T_0\le I\), its scalar return moments
\(c_k(e)=\langle u_e,T_0^{k-1}u_e\rangle\) satisfy
\begin{equation}\label{eq:returntelescoping}
 c_k(e)-c_{k+1}(e)
 =\langle u_e,T_0^{k-1}D_0u_e\rangle\ge0.
\end{equation}
If, in addition, \(u_e\in\operatorname{Ran}D_0^{1/2}\), the spectral
theorem gives the exact sum
\begin{equation}\label{eq:returnsum}
 \sum_{k\ge1}c_k(e)
 =\|D_0^{\dagger/2}(D_0h_e-q_e)\|^2.
\end{equation}
Thus the harmonic contribution telescopes; the only nontrivial inverse
datum is the quadratic form
\(\langle q_e,D_0^\dagger q_e\rangle\), with its explicit cross term.
Because all operators are first compressed by the two-moment projection,
every returned vector remains in \(\mathcal P_T\).  Put
\(F_{k,e}=R_0^{\dagger/2}T_0^{k/2}q_e\).  The defining balanced
factorization then identifies each scalar dissipation as
\begin{equation}\label{eq:returnDissipation}
 \langle q_e,T_0^kq_e\rangle-\langle q_e,T_0^{k+1}q_e\rangle
 =-B_{\rm nuc}(F_{k,e},F_{k,e}).
\end{equation}
It retains the complete Gamma place and every \(p^j\), after the continuous
Chebyshev term has been cancelled by the two exact Tate moments.  Equation
\eqref{eq:returnload} and the telescoping identity
\eqref{eq:returntelescoping} are unconditional algebraic identities under
the old-core hypothesis.  The range condition needed for
\eqref{eq:returnsum} is precisely part of the still-open capacity estimate,
so this reduction does not assume the missing enlarged-cell sign.

Finally, the positive reference itself has a uniform exact coercivity which
does not solve the signed problem.  With zero-extension channels
\(\widehat J_{n,-}F=(\widetilde F(\cdot+\log n)-\widetilde F)/\sqrt2\), put
\[
 \widehat{\mathcal R}_T=\mathcal H_{5/4}
 +\sum_{n\le N}\frac{\Lambda(n)}{\sqrt n}
       \widehat J_{n,-}^*\widehat J_{n,-},
 \qquad N=\lfloor e^{2T}\rfloor .
\]

\begin{proposition}[Shift-chain coercivity]\label{prop:shiftchaincoercivity}
One has
\begin{equation}\label{eq:shiftchaincoercivity}
 \widehat{\mathcal R}_T\ge \alpha_N I,\qquad
 \alpha_N=\sum_{n\le N}\frac{\Lambda(n)}{\sqrt n}
 \left(1-\cos\frac{\pi}{\lceil2T/\log n\rceil+1}\right),
\end{equation}
and
\begin{equation}\label{eq:shiftchainasymptotic}
 \alpha_N\ge \frac{2}{9(2T)^2}
 \sum_{n\le N}\frac{\Lambda(n)(\log n)^2}{\sqrt n}
 \ge\left(\frac49+o(1)\right)\sqrt N,
 \qquad \alpha_N\sim\sqrt N.
\end{equation}
The same bound holds after restriction to \(\mathcal P_T\).
\end{proposition}

\begin{proof}
Disintegrate \((0,2T)\) into chains \(r+j\log n\).  Each chain has at most
\(m_n=\lceil2T/\log n\rceil\) vertices.  The Hermitian part of its nilpotent
unilateral shift has norm \(\cos(\pi/(m_n+1))\), proving
\eqref{eq:shiftchaincoercivity} after summing and discarding the nonnegative
Gamma form.  Since \(m_n+1\le3(2T)/\log n\) and
\(1-\cos x\ge2x^2/\pi^2\) on \([0,\pi/2]\), the first inequality in
\eqref{eq:shiftchainasymptotic} follows.  Keeping only primes, applying the
prime number theorem to
\(\int_2^N(\log x)^2x^{-1/2}\,d\vartheta(x)\), and using
\(2T=\log N+O(N^{-1})\) gives the last asymptotic.  Restriction cannot weaken
a quadratic lower bound.  For the sharp equivalent, fix \(K>0\).  On
\(Ne^{-K}\le n<N\), and for sufficiently large \(N\), the shift chain has
two vertices and its bracket in \eqref{eq:shiftchaincoercivity} is exactly
\(1-\cos(\pi/3)=1/2\).  The prime number theorem and partial summation give
\(\sum_{n\le x}\Lambda(n)n^{-1/2}=2\sqrt x+o(\sqrt x)\), hence
\[
 \liminf_{N\to\infty}\frac{\alpha_N}{\sqrt N}\ge1-e^{-K/2}.
\]
Letting \(K\to\infty\) gives the lower limit one.  Apart from a possible
atom at \(n=N\), of size \(o(\sqrt N)\), every bracket is at most \(1/2\).
The same prime-number-theorem estimate gives the matching upper limit.
\end{proof}

The arithmetic part of the return series also admits a uniform all-depth
bound.  Write \(\Lambda_k=\Lambda^{*k}\) for Dirichlet convolution and set
\[
 V_{N,k}=\sum_{m\le N}\frac{\Lambda_k(m)^2}{m},\qquad
 H_{N,k}=\frac{\Lambda_{2k}(N)}{\sqrt N}.
\]

\begin{proposition}[Uniform Witt-moment summability]
\label{prop:uniformWittmoments}
For every \(\varepsilon>0\) there are \(\delta>0\) and
\(C_\varepsilon<\infty\) such that, uniformly for
\(1\le k\le\delta\log N\),
\begin{equation}\label{eq:uniformWittsmall}
 V_{N,k}\le C_\varepsilon\sqrt{k}\,e^{\varepsilon k}
 \frac{k!}{(2k)!}(\log N)^{2k}.
\end{equation}
For every fixed \(c>0\),
\begin{equation}\label{eq:uniformWittlarge}
 \sum_{k\ge\delta\log N}\frac{V_{N,k}}{(c\log N)^{2k}}=o(1),
 \qquad
 \sup_N\sum_{k\ge1}\frac{V_{N,k}+H_{N,k}}{(c\log N)^{2k}}<\infty.
\end{equation}
Thus every arithmetic Witt depth is uniformly summable.  These estimates
do not assert that the total capacity is at most the unit Schur budget.
\end{proposition}

\begin{proof}
Let \((z_p)_p\) be independent Haar variables on the unit circle and, for
\(\sigma>0\), put
\[
 g_\sigma(z)=\sum_p(\log p)\sum_{j\ge1}
 p^{-j(1/2+\sigma)}z_p^j.
\]
The coefficient of the monomial indexed by \(m\) in \(g_\sigma^k\) is
\(\Lambda_k(m)m^{-1/2-\sigma}\).  Bohr orthogonality and Rankin's
inequality therefore give the exact identity and bound
\begin{equation}\label{eq:BohrRankinWitt}
 \|g_\sigma\|_{2k}^{2k}
 =\sum_{m\ge1}\frac{\Lambda_k(m)^2}{m^{1+2\sigma}},
 \qquad
 V_{N,k}\le N^{2\sigma}\|g_\sigma\|_{2k}^{2k}.
\end{equation}
Split \(g_\sigma=P_\sigma+Q_\sigma\) into the terms with \(j=1\) and
\(j\ge2\).  If
\(\mathcal V(\sigma)=\sum_p(\log p)^2p^{-1-2\sigma}\), expansion of the
Steinhaus moment gives
\[
 \|P_\sigma\|_{2k}^{2k}\le k!\mathcal V(\sigma)^k,
 \qquad
 \mathcal V(\sigma)=\frac{1+o(1)}{4\sigma^2}.
\]
For the proper-power part,
\[
 |Q_{\sigma,p}|\le \frac{(\log p)p^{-1}}{1-p^{-1/2}},
 \qquad
 \sum_p\left(\frac{(\log p)p^{-1}}{1-p^{-1/2}}\right)^2<\infty.
\]
Hoeffding's lemma applied separately to real and imaginary parts hence
implies \(\|Q_\sigma\|_{2k}\le C\sqrt k\), uniformly in \(\sigma\).
Minkowski's inequality consequently yields, for small \(\sigma\),
\begin{equation}\label{eq:Wittmomentmaster}
 \|g_\sigma\|_{2k}^{2k}
 \le k!\mathcal V(\sigma)^k(1+C_1\sigma)^{2k}.
\end{equation}
Take \(\sigma=k/\log N\).  The prime number theorem estimate for
\(\mathcal V\), \eqref{eq:BohrRankinWitt},
\eqref{eq:Wittmomentmaster}, and Stirling's inequalities prove
\eqref{eq:uniformWittsmall}, after choosing \(\delta\) in terms of
\(\varepsilon\).  For \(k\ge\delta\log N\), instead fix a small
\(\sigma_0>0\).  The same moment argument gives
\[
 \frac{V_{N,k}}{(c\log N)^{2k}}
 \le \exp\!\left(2\sigma_0\log N-
                 k\log\log N+O_{c,\delta,\sigma_0}(k)\right),
\]
uniformly up to \(k\le\log N/\log2\); beyond that range
\(\Lambda_k(m)=0\) for \(m\le N\).  Summing proves the first part of
\eqref{eq:uniformWittlarge}.  Finally \(0\le H_{N,k}\le V_{N,k}\), since
\(H_{N,k}\) is the off-diagonal entry of the positive two-sided word Gram
whose two diagonal entries equal \(V_{N,k}\).  The small-depth majorant is
summable because \(k!/(2k)!\) beats every exponential; this proves the
last assertion.
\end{proof}

Both signs in the balanced factorization receive the same zero-extension
boundary energy.  Proposition~\ref{prop:shiftchaincoercivity} therefore
controls the reference inverse but does not prove
\eqref{eq:sharpdomination}.  The remaining global theorem is the uniform
nonnegativity of \eqref{eq:thresholdcapacity}, equivalently a Douglas
factorization of the actual centered atomic--continuous cross \(Q_c\)
through \(D_0^{1/2}\) with the required \emph{unit} shorted-capacity bound.
Proposition~\ref{prop:uniformWittmoments} proves summability of every purely
arithmetic return depth, but neither identifies this centered Douglas
factor nor pays its total unit budget.  This is the precise point at which the local-analysis argument stops; the
Gamma--Fourier--Tate source construction below begins from this capacity
defect.

% END INLINED SOURCE: row-d-local-analysis


\begin{theorem}[Exact geometric row-(d) trilemma]
\label{thm:rowdtrilemma}
A non-circular geometric proof of row~(d) must either:
\begin{enumerate}
\item factor the nuclear form through $N^1_{\mathrm{rul}}$;
\item extend the intrinsic row-(a) periodic section-cohomology theory from the two rulings to the
completed mixed correspondence module, with effectivity, Riemann--Roch,
continuity and an equality theorem; or
\item prove a new Hodge--Rosati theorem intrinsically in the nuclear category
before comparing its form with \eqref{eq:Bnuc}.
\end{enumerate}
The first alternative is impossible by
Theorem~\ref{thm:infiniterankcontact}; the data required by the second and
the theorem required by the third have not been constructed.
\end{theorem}

The minimum positive result would construct an infinite-dimensional space
$N_D$, a class $H$ with $H^2>0$, a degree map and a continuous intersection
form containing \eqref{eq:requiredprimitiveclass}; prove from its own
Riemann--Roch and effectivity theory that $M_f^2\le0$, including the equality
case; and only afterwards prove $M_fM_g=B_{\mathrm{nuc}}(f,g)$.  Such a theorem would close row~(d).  The source construction below gives a
second non-circular route to the same primitive inequality: an independently
positive Gamma--Euler--Tate factorization.  Neither route is completed here.


% BEGIN CLEAN GAMMA--FOURIER--TATE ROW (d)
\subsection{The Gamma--Fourier--Tate boundary construction}
\label{ss:gammafouriertate}

The local analysis above proves the primitive inequality on a nontrivial
initial interval and identifies the obstruction to propagation as a shorted
old/new coupling.  We now construct the arithmetic and archimedean boundary
data entering that coupling.  The construction is independent of the zeros
of $\xi$.  It fixes the amplitude, the moving Tate constraints and the
positive Gamma storage before the remaining estimate is formulated.

\subsubsection{The finite logarithmic boundary port}

Let $\mathcal H_N$ be the Hilbert space with orthonormal basis
$\{\delta_d:1\leq d\leq N\}$.  For $n\leq N$, let
$S_n\delta_d=\delta_{nd}$ when $nd\leq N$ and let it be zero otherwise.
For $t>0$ put
\begin{equation}\label{eq:newdZetaConnection}
 \mathscr Z_t=\sum_{n\leq N}n^{-t/2}S_n,
 \qquad
 \mathscr C_t=\sum_{n\leq N}\Lambda(n)n^{-t/2}S_n,
 \qquad Q\delta_d=(\log d)\delta_d.
\end{equation}
All the sums in this subsection belong to the finite divisibility algebra.
In particular, no convergence or domain assertion is implicit.

\begin{lemma}[Finite logarithmic connection]
\label{lem:newdFiniteConnection}
One has
\begin{align}
 [Q,\mathscr Z_t]&=\mathscr Z_t\mathscr C_t
                  =\mathscr C_t\mathscr Z_t,
                  \label{eq:newdCommutator}\\
 \mathscr Z_t'&=-\tfrac12\mathscr Z_t\mathscr C_t,
                  \label{eq:newdZetaDerivative}\\
 (\mathscr Z_t^{-1})'&=\tfrac12\mathscr C_t\mathscr Z_t^{-1}.
                  \label{eq:newdMobiusDerivative}
\end{align}
\end{lemma}

\begin{proof}
Since $[Q,S_n]=(\log n)S_n$, the coefficient of $S_k$ in the
commutator is $(\log k)k^{-t/2}$.  The coefficient of $S_k$ in
$\mathscr Z_t\mathscr C_t$ is
\[
 k^{-t/2}\sum_{d\mid k}\Lambda(d)=k^{-t/2}\log k.
\]
The algebra is commutative, which proves \eqref{eq:newdCommutator}.
Termwise differentiation proves \eqref{eq:newdZetaDerivative};
differentiating $\mathscr Z_t^{-1}\mathscr Z_t=I$ then proves
\eqref{eq:newdMobiusDerivative}.
\end{proof}

We record the operator identity for which this connection is needed.  Let
$G(t)$ be the positive Gram family at a threshold, let
$v(t)=(-h(t),I)^t$ be its harmonic residual, and let $P_0$ be the old
coordinate projection.  Write $M_0$ for the old block of $G(t_0)$.  The
normal equation is $P_0G(t)v(t)=0$.  Hence differentiation gives
\begin{equation}\label{eq:newdMovingNormal}
 P_0G'(t_0)v(t_0)=M_0h'(t_0).
\end{equation}
Suppose, as in the localized explicit-formula decomposition, that
\[
 G'=R^*\mathscr LR-\tfrac12(QG+GQ),
 \qquad Gv=S\delta_1,
\]
where $S>0$ is the scalar normalization.  With
\[
 W_{\rm H}=\mathscr L^{1/2}Rv,
 \qquad H_0=\mathscr L^{1/2}RP_0^*,
\]
equation \eqref{eq:newdMovingNormal} becomes
\begin{equation}\label{eq:newdFiniteExcess}
 H_0^*W_{\rm H}=M_0h'+B_Q,
 \qquad
 B_Q=\frac S2P_0
 \bigl[-\mathscr C\delta_1+\mathscr Z\mathscr C^*m\bigr],
 \qquad m=\mathscr Z^{-1}\delta_1,
\end{equation}
where the parameter is now fixed and omitted from the notation.

Define the doubled logarithmic boundary port by
\begin{align}
 W_\infty
 &=\sqrt{\frac S2}
   \binom{\mathscr C\delta_1}{\mathscr C^*m},
   \label{eq:newdWinfty}\\
 H_{\infty,0}x
 &=\sqrt{\frac S2}
   \binom{P_0^*x}{-\mathscr Z^*P_0^*x}.
   \label{eq:newdHinfty}
\end{align}

\begin{theorem}[Exact finite Gamma--Tate amplitude]
\label{thm:newdAmplitude}
The port \eqref{eq:newdWinfty}--\eqref{eq:newdHinfty} is defined entirely
by the finite Euler transport and satisfies
\begin{equation}\label{eq:newdAmplitude}
 H_{\infty,0}^*W_\infty=-B_Q.
\end{equation}
Consequently
\begin{equation}\label{eq:newdCompletedDivergence}
 (H_0^*\oplus H_{\infty,0}^*)
 (W_{\rm H}\oplus W_\infty)=M_0h'.
\end{equation}
\end{theorem}

\begin{proof}
Taking the adjoint of \eqref{eq:newdHinfty} and applying it to
\eqref{eq:newdWinfty} gives
\[
 H_{\infty,0}^*W_\infty
 =\frac S2P_0
   (\mathscr C\delta_1-\mathscr Z\mathscr C^*m)
 =-B_Q
\]
by \eqref{eq:newdFiniteExcess}.  Adding this equality to
\eqref{eq:newdFiniteExcess} proves
\eqref{eq:newdCompletedDivergence}.
\end{proof}

The theorem fixes both orientations and the normalization of the boundary
amplitude.  It is an amplitude identity, not an energy inequality: its
finite-dimensional port cannot be identified with the complete
infinite-rank Hodge factor.  The continuum realization below uses it only
for the arithmetic cross-incidence that it computes.

\subsubsection{The moving Tate terminal}

Let $H_0,E$ be Hilbert spaces, let $B_0:H_0\to\C^r$ be bounded and onto,
and let $B_E:E\to\C^r$ be bounded.  Put
\begin{equation}\label{eq:newdMovingData}
 K_0=\ker B_0,
 \quad G_0=B_0B_0^*,
 \quad C=B_0^*G_0^{-1}B_E,
 \quad M=I_E+C^*C=I_E+B_E^*G_0^{-1}B_E.
\end{equation}
The enlarged moment map is $B(h,e)=B_0h+B_Ee$.

\begin{theorem}[Conservative update of a moving constraint]
\label{thm:newdMovingConstraint}
There is an orthogonal decomposition
\begin{equation}\label{eq:newdMovingDecomposition}
 \ker B=(K_0\oplus0)\ \widehat\oplus\
          \{(-Ce,e):e\in E\}.
\end{equation}
Moreover
\begin{equation}\label{eq:newdMovingIsometry}
 W:E\longrightarrow H_0\oplus E,
 \qquad
 Wx=\binom{-CM^{-1/2}x}{M^{-1/2}x}
\end{equation}
is an isometry onto the second summand.  In particular,
\begin{equation}\label{eq:newdMovingConservation}
 \|M^{-1/2}x\|^2+\|CM^{-1/2}x\|^2=\|x\|^2.
\end{equation}
\end{theorem}

\begin{proof}
The surjectivity of $B_0$ gives
$H_0=K_0\widehat\oplus\operatorname{Ran} B_0^*$.  If $h=k+r$ is the corresponding
decomposition, then $B(h,e)=0$ is equivalent to $B_0r=-B_Ee$ and hence to
$r=-B_0^*G_0^{-1}B_Ee=-Ce$.  This proves
\eqref{eq:newdMovingDecomposition}.  Finally,
\[
 \|Wx\|^2
 =\langle M^{-1/2}x,(I+C^*C)M^{-1/2}x\rangle=\|x\|^2,
\]
which proves both remaining assertions.
\end{proof}

For $I_T=(-T,T)$ take the two Tate moments
\begin{equation}\label{eq:newdTateMomentMap}
 B_Tf=
 \binom{\langle f,e^{t/2}\rangle}{\langle f,e^{-t/2}\rangle}.
\end{equation}
Their Gram is
\begin{equation}\label{eq:newdTateGram}
 B_TB_T^*=2
 \begin{pmatrix}\sinh T&T\\T&\sinh T\end{pmatrix},
\end{equation}
with eigenvalues $2(\sinh T\pm T)>0$.  Applying
Theorem~\ref{thm:newdMovingConstraint} to
$L^2(I_{T'})=L^2(I_T)\oplus L^2(I_{T'}\setminus I_T)$ gives an explicit
orthogonal update of the primitive space $\ker B_T$.  At the arithmetic
thresholds
\[
 T_N=\tfrac12\log N,
 \qquad
 \delta_N=T_{N+1}-T_N=\tfrac12\log\frac{N+1}{N},
\]
the two nonzero eigenvalues of $C^*C$ are
\begin{equation}\label{eq:newdTateReturnEigenvalues}
 \kappa_{N,\pm}=
 \frac{\sinh T_{N+1}-\sinh T_N\ \pm\delta_N}
      {\sinh T_N\ \pm T_N}.
\end{equation}
Therefore
\begin{equation}\label{eq:newdTateReturnDecay}
 \|C\|^2=\frac1{2N}+o(N^{-1}),
 \qquad
 \|CM^{-1/2}\|=O(N^{-1/2}),
 \qquad
 \|M^{-1/2}-I\|=O(N^{-1}).
\end{equation}
Thus the motion of the two Tate constraints is an explicit conservative
rank-two terminal; it creates no unidentified projection term.

We next identify the finite incidence column with its physical
old/new boundary action.  Put $I_N=(-T_N,T_N)$ and
\[
 \mathcal E_N=
 (-T_{N+1},-T_N)\cup(T_N,T_{N+1}).
\]
Let $U_n$ be translation by $\log n$ on the zero-extension line and set
\begin{equation}\label{eq:newdPhysicalIncidence}
 \mathcal B_N^{\rm ar}
 =-P_{I_N}\sum_{2\leq n\leq N}
       \frac{\Lambda(n)}{\sqrt n}(U_n+U_n^*)
       \big|_{L^2(\mathcal E_N)}.
\end{equation}
After separating the two newborn wings, the translated intervals give an
isometry
\[
 \mathfrak S_N:
 \ell^2\{2,\ldots,N\}\widehat\otimes L^2(\mathcal E_N)
 \longrightarrow L^2(I_N),
\]
where the $n$th coordinate is the unique orientation entering $I_N$.

\begin{proposition}[Physical incidence fibre]
\label{prop:newdPhysicalIncidence}
If
\[
 \mathbf c_N=\sum_{n=2}^N\frac{\Lambda(n)}{\sqrt n}e_n,
\]
then, for every newborn profile $e$,
\begin{equation}\label{eq:newdPhysicalIncidenceColumn}
 \mathcal B_N^{\rm ar}e=-\mathfrak S_N(\mathbf c_N\otimes e),
\end{equation}
and
\begin{equation}\label{eq:newdPhysicalIncidenceGram}
 (\mathcal B_N^{\rm ar})^*\mathcal B_N^{\rm ar}
 =\left(\sum_{n=2}^N\frac{\Lambda(n)^2}{n}\right)I.
\end{equation}
At word depth $k$, the corresponding label column is
\begin{equation}\label{eq:newdWordColumn}
 (\log N)^{-k}\mathscr C_N^k\delta_1
 =\left\{\frac{\Lambda_k(n)}
                {\sqrt n(\log N)^k}\right\}_{n\leq N}.
\end{equation}
\end{proposition}

\begin{proof}
For a vector on the right newborn wing only $U_n$ enters the old interval;
for a vector on the left wing only $U_n^*$ enters it.  These translated
copies are pairwise disjoint, proving
\eqref{eq:newdPhysicalIncidenceColumn}; isometry gives
\eqref{eq:newdPhysicalIncidenceGram}.  Multiplication of an integer label by
$m$ agrees, whenever the result remains in the threshold fibre, with
physical translation by $\log m$.  Iterating and using Dirichlet convolution
gives \eqref{eq:newdWordColumn}.
\end{proof}

\subsubsection{Euler incidence and rational leakage}

On the right-oriented threshold fibre label the interval
$T_N-\log m+(0,\delta_N)$ by $m\leq N$.  Translation by $\log n$ sends
$m$ to $mn$, while translation in the opposite direction sends it to
$m/n$.  The latter remains in the integer fibre exactly when $n\mid m$.

\begin{proposition}[Divisibility and rational-leakage decomposition]
\label{prop:newdLeakageSplit}
The finite Euler boundary action is the orthogonal sum of its action on the
integer divisibility fibre and an action supported on nonintegral rational
labels.  An integer-labelled cell never overlaps a noninteger-labelled cell
in positive measure.  Two rational cells labelled by $m/n$ and $m'/n'$ can
overlap only when
\begin{equation}\label{eq:newdOverlapCondition}
 \left|\log\frac{m/n}{m'/n'}\right|<\delta_N.
\end{equation}
\end{proposition}

\begin{proof}
The assertions about labels follow directly from translation of the cell
positions.  If $m>kn$, then
\[
 \frac{m}{kn}\geq\frac N{N-1}>\frac{N+1}{N}=e^{2\delta_N};
\]
if $m<kn$, then $kn/m\geq(N+1)/N$, with equality producing at most an
endpoint contact.  Hence integer and noninteger cells are orthogonal.
Two intervals of length $\delta_N$ overlap exactly under
\eqref{eq:newdOverlapCondition}.
\end{proof}

Write $\Lambda_k=\Lambda^{*k}$.  For coprime positive integers $a,b$ with
$b>1$ define
\begin{equation}\label{eq:newdLeakageCoefficient}
 \ell_{N,k}(a,b)=
 \frac1{\sqrt{ab}(\log N)^k}
 \sum_{d\leq N/\max(a,b)}
       \frac{\Lambda(bd)\Lambda_k(ad)}d.
\end{equation}
Let
\[
 \mathscr R_N^{\rm leak}
 =\{a/b:(a,b)=1,\ b>1,\ \max(a,b)\leq N\},
\]
extend $e\in E_N:=L^2(0,\delta_N)$ by zero, and define
\begin{equation}\label{eq:newdLeakageOperator}
 L_{N,k}e=\sum_{a/b\in\mathscr R_N^{\rm leak}}
   \ell_{N,k}(a,b)\,\tau_{\log(a/b)}e,
 \qquad
 \mathbf L_Ne=\bigoplus_{k\geq1}L_{N,k}e.
\end{equation}
Only finitely many depths occur for a fixed $N$.

\begin{proposition}[Exact aggregation on reduced rational labels]
\label{prop:newdRationalAggregation}
The coefficient of the oppositely oriented Euler output on the reduced
rational cell $a/b$ is the sum in
\eqref{eq:newdLeakageCoefficient}.  If
$\kappa_e(u)=\langle\tau_ue,e\rangle$, then
\begin{align}
 \|\mathbf L_Ne\|^2
 ={}&\sum_{k\geq1}
 \sum_{a/b,c/d\in\mathscr R_N^{\rm leak}}
 \ell_{N,k}(a,b)\overline{\ell_{N,k}(c,d)}
 \kappa_e\!\left(\log\frac{ad}{bc}\right).
 \label{eq:newdLeakageGram}
\end{align}
Thus the Gram retains all coherent collisions between equal or nearby
reduced fractions.
\end{proposition}

\begin{proof}
The equality $m/n=a/b$ with $(a,b)=1$ is equivalent to
$m=ad$, $n=bd$ for a unique positive integer $d$.  Substitution of the
Euler weight and of the $k$th word column gives
\eqref{eq:newdLeakageCoefficient}.  Expanding the squared norm in
\eqref{eq:newdLeakageOperator} gives
\eqref{eq:newdLeakageGram}, since the inner product of the cells labelled by
$a/b$ and $c/d$ is $\kappa_e(\log(ad/bc))$.
\end{proof}

\subsubsection{Positive Gamma storage and the two exterior Tate states}

For $a>0$ and $f\in C_c^\infty(\R)$, let $z_{a,f}$ be the causal state
\begin{equation}\label{eq:newdGammaState}
 z_{a,f}'=-az_{a,f}+\sqrt{2a}\,f,
 \qquad z_{a,f}(-\infty)=0,
\end{equation}
and set $g_{a,f}=\sqrt{2a}\,z_{a,f}-f$.

\begin{lemma}[Lossless Gamma state]
\label{lem:newdGammaState}
The state satisfies
\begin{align}
 \frac d{dt}|z_{a,f}|^2&=|f|^2-|g_{a,f}|^2,
 \label{eq:newdGammaConservation}\\
 \frac1{2a}\|f-g_{a,f}\|_2^2
 &=\frac1{a^2}\int_\R|z_{a,f}'(t)|^2\,dt
 =\int_\R\frac{2\tau^2}{a(a^2+\tau^2)}
       |\widehat f(\tau)|^2\frac{d\tau}{2\pi}.
 \label{eq:newdGammaEnergy}
\end{align}
\end{lemma}

\begin{proof}
Expanding $|g_{a,f}|^2$ and using
\eqref{eq:newdGammaState} proves
\eqref{eq:newdGammaConservation}.  Solving
\eqref{eq:newdGammaState} for $f$ gives
$f-g_{a,f}=\sqrt{2/a}\,z_{a,f}'$.  Its Fourier multiplier has squared
modulus $4\tau^2/(a^2+\tau^2)$, proving
\eqref{eq:newdGammaEnergy}.
\end{proof}

Put $a_j=2j+\tfrac12$ and define the unreset Gamma energy
\begin{equation}\label{eq:newdFullGammaEnergy}
 \mathcal G_\Gamma[f]
 =\sum_{j\geq0}\frac1{a_j^2}\int_\R|z_{a_j,f}'(t)|^2\,dt
 =\int_\R g_\Gamma(\tau)|\widehat f(\tau)|^2
      \frac{d\tau}{2\pi},
\end{equation}
where
\begin{equation}\label{eq:newdGammaMultiplier}
 g_\Gamma(\tau)=
 \sum_{j\geq0}\frac{2\tau^2}{a_j(a_j^2+\tau^2)}.
\end{equation}
The causal equations are solved once on the real line; states are not reset
at internal rational-cell endpoints.

For a compactly supported $f$ define
\begin{equation}\label{eq:newdTateMoments}
 M_+(f)=\int_\R e^{t/2}f(t)\,dt,
 \qquad
 M_-(f)=\int_\R e^{-t/2}f(t)\,dt.
\end{equation}
If $I=(L,R)$ contains the support of $f$, the Gram of the two Tate vectors is
\begin{equation}\label{eq:newdExteriorTateGram}
 G_I=
 \begin{pmatrix}
 e^R-e^L&R-L\\
 R-L&e^{-L}-e^{-R}
 \end{pmatrix}>0,
\end{equation}
and the squared norm of the Tate component is
\begin{equation}\label{eq:newdTateTerminalForm}
 \mathcal T_I[f]=
 \begin{pmatrix}\overline{M_+(f)}&\overline{M_-(f)}\end{pmatrix}
 G_I^{-1}
 \binom{M_+(f)}{M_-(f)}.
\end{equation}
The moments in \eqref{eq:newdTateMoments} are exactly the exterior states of
\eqref{eq:newdGammaState} at $a=1/2$: for
$\operatorname{supp}f\subset[L,R]$,
\[
 z_{1/2,f}(t)=e^{-t/2}M_+(f)\quad(t\geq R),
\]
up to the normalization $\sqrt{2a}=1$, and the anticausal state gives
$e^{t/2}M_-(f)$ on $t\leq L$.

Take one interval containing every cell in
\eqref{eq:newdLeakageOperator}, for instance
\begin{equation}\label{eq:newdGlobalLeakageInterval}
 I_N=\left(-\log N,\ \log\frac N2+\delta_N\right).
\end{equation}
The source-defined Gamma--Tate form on the physical leakage column is
\begin{equation}\label{eq:newdGammaTateForm}
 \mathcal E_{\Gamma T,N}(e)=
 \sum_{k\geq1}
 \bigl\{\mathcal G_\Gamma[L_{N,k}e]
       +\mathcal T_{I_N}[L_{N,k}e]\bigr\}.
\end{equation}
Every term in this form is nonnegative and is determined by prime-power
weights, the Gamma parameters and the two Tate moments.  In particular, it
does not use the sign of the primitive Weil form.

\paragraph{The first nontrivial arithmetic cell.}
The threshold $N=3$ gives
\begin{equation}\label{eq:newdFirstCellLength}
 T_3=\tfrac12\log3,
 \qquad T_4=\log2,
 \qquad L_3=\delta_3=\tfrac12\log\frac43.
\end{equation}
At word depth one the only nonzero nonintegral coefficient is
\begin{equation}\label{eq:newdFirstLeakageCoefficient}
 \ell_{3,1}(2,3)
 =\frac{\Lambda(3)\Lambda(2)}{\sqrt6\log3}
 =\frac{\log2}{\sqrt6}.
\end{equation}
There is no deeper word, because $\Lambda_k(m)=0$ for $m\leq3$ and
$k\geq2$.  Consequently the complete rational leakage in this threshold
cell is
\begin{equation}\label{eq:newdFirstLeakageOutput}
 \mathbf L_3e=f_e
 :=\frac{\log2}{\sqrt6}\,
   \tau_{\log(2/3)}e,
 \qquad
 \|f_e\|_2^2=\frac{(\log2)^2}{6}\|e\|_2^2.
\end{equation}

\begin{theorem}[Complete Gamma--Tate payment on one arithmetic cell]
\label{thm:newdFirstCellPayment}
For every $e\in L^2(0,L_3)$, let $f_e$ be given by
\eqref{eq:newdFirstLeakageOutput}.  Then
\begin{equation}\label{eq:newdFirstCellPayment}
 \|\mathbf L_3e\|_2^2
 \leq \mathcal G_\Gamma[f_e]+\mathcal T_{J_3}[f_e],
 \qquad
 J_3=\log(2/3)+(0,L_3).
\end{equation}
More precisely, if $\mathcal G_{0}[f]$ denotes the first Gamma-channel
energy, corresponding to $a_0=1/2$, then
\begin{equation}\label{eq:newdFirstCellStrictPayment}
 \|f_e\|_2^2
 \leq \alpha_3\mathcal G_0[f_e]
       +\eta_3\mathcal T_{J_3}[f_e],
\end{equation}
where
\begin{equation}\label{eq:newdFirstCellConstants}
 \alpha_3=\frac14+\frac{L_3^2}{32}<0.251,
 \qquad
 \eta_3=\frac{1-e^{-L_3}}2
       =\frac{2-\sqrt3}{4}<0.067.
\end{equation}
Thus the Gamma interior energy and the two exterior Tate states pay the
whole leakage simultaneously, with a strict positive remainder:
\begin{align}
 &\mathcal G_\Gamma[f_e]+\mathcal T_{J_3}[f_e]
       -\|\mathbf L_3e\|_2^2                              \notag\\
 &\qquad\geq
 (1-\alpha_3)\mathcal G_0[f_e]
 +(1-\eta_3)\mathcal T_{J_3}[f_e]\geq0.
 \label{eq:newdFirstCellRemainder}
\end{align}
\end{theorem}

\begin{proof}
Translation does not change any of the following estimates, so write
$J_3=(0,L_3)$.  Let $z$ be the causal state of the first Gamma channel,
\[
 z'=-\tfrac12z+f_e,
 \qquad z(0)=0.
\]
The initial value vanishes because $f_e$ is supported on this single cell.
Since $f_e=z'+\tfrac12z$, integration of the cross term gives the exact
identity
\begin{equation}\label{eq:newdFirstCellStateIdentity}
 \|f_e\|_2^2
 =\|z'\|_2^2+\frac14\|z\|_2^2+\frac12|z(L_3)|^2.
\end{equation}
The anchored Cauchy--Schwarz inequality yields
\[
 |z(t)|^2
 \leq t\int_0^t|z'(s)|^2\,ds
 \leq t\|z'\|_2^2,
\]
and hence
\begin{equation}\label{eq:newdFirstCellPoincare}
 \|z\|_2^2\leq\frac{L_3^2}{2}\|z'\|_2^2.
\end{equation}
By \eqref{eq:newdGammaEnergy} at $a=1/2$,
\begin{equation}\label{eq:newdFirstCellGammaState}
 \mathcal G_0[f_e]=4\|z'\|_2^2.
\end{equation}

It remains to pay the endpoint.  The exterior-state formula gives
\[
 z(L_3)=e^{-L_3/2}M_+(f_e).
\]
The Tate terminal is the squared norm of the orthogonal projection onto
$\operatorname{span}\{e^{t/2},e^{-t/2}\}$.  It therefore dominates the
projection onto the first of these one-dimensional subspaces:
\begin{equation}\label{eq:newdFirstCellTateProjection}
 \mathcal T_{J_3}[f_e]
 \geq\frac{|M_+(f_e)|^2}{\|e^{t/2}\|_{L^2(J_3)}^2}
 =\frac{|M_+(f_e)|^2}{e^{L_3}-1}.
\end{equation}
Consequently
\begin{equation}\label{eq:newdFirstCellEndpoint}
 \frac12|z(L_3)|^2
 \leq\frac{1-e^{-L_3}}2\,\mathcal T_{J_3}[f_e]
 =\eta_3\mathcal T_{J_3}[f_e].
\end{equation}
Substitution of \eqref{eq:newdFirstCellPoincare},
\eqref{eq:newdFirstCellGammaState} and
\eqref{eq:newdFirstCellEndpoint} into
\eqref{eq:newdFirstCellStateIdentity} proves
\eqref{eq:newdFirstCellStrictPayment}.  Since the full Gamma energy contains
$\mathcal G_0$ as its first nonnegative summand and both constants in
\eqref{eq:newdFirstCellConstants} are smaller than one,
\eqref{eq:newdFirstCellPayment} and
\eqref{eq:newdFirstCellRemainder} follow.
\end{proof}

This theorem proves the complete one-cell assertion without a sign
assumption and without resetting a state inside the cell.  Its causal state
starts at zero only because, at $N=3$, the whole rational leakage consists
of the single support interval $J_3$.  For a connected cluster at a later
threshold, the outgoing state of one rational cell is the incoming state of
the next; Theorem~\ref{thm:newdFirstCellPayment} cannot be iterated by
setting those intermediate states to zero.

The same argument in fact applies to a whole connected cluster, provided it
is applied once on the exterior hull rather than separately on its
constituent cells.

\begin{lemma}[Gamma--Tate payment on a finite support hull]
\label{lem:newdClusterHull}
Let $f\in L^2(\R)$ be supported in an interval $J=(r,r+D)$.  If
$D<\sqrt{24}$, then
\begin{equation}\label{eq:newdClusterHullPayment}
 \|f\|_2^2
 \leq \alpha(D)\mathcal G_0[f]
       +\eta(D)\mathcal T_J[f]
 \leq \mathcal G_\Gamma[f]+\mathcal T_J[f],
\end{equation}
where
\begin{equation}\label{eq:newdClusterHullConstants}
 \alpha(D)=\frac14+\frac{D^2}{32}<1,
 \qquad
 \eta(D)=\frac{1-e^{-D}}2<1.
\end{equation}
No disjointness or orthogonality assumption is made inside $J$.
\end{lemma}

\begin{proof}
Translate $J$ to $(0,D)$ and solve the first causal Gamma equation once on
the complete hull:
\[
 z'=-\tfrac12z+f,
 \qquad z(0)=0.
\]
There is no condition on $z$ at an internal subdivision.  The proof of
Theorem~\ref{thm:newdFirstCellPayment}, with $D$ in place of $L_3$, gives
\begin{align*}
 \|f\|_2^2
 &=\|z'\|_2^2+\tfrac14\|z\|_2^2+\tfrac12|z(D)|^2,\\
 \|z\|_2^2&\leq\tfrac{D^2}{2}\|z'\|_2^2,
 \qquad
 \mathcal G_0[f]=4\|z'\|_2^2,\\
 \tfrac12|z(D)|^2
 &\leq\tfrac{1-e^{-D}}2\mathcal T_J[f].
\end{align*}
Combining these identities proves the first inequality in
\eqref{eq:newdClusterHullPayment}.  The restriction $D<\sqrt{24}$ is
exactly $\alpha(D)<1$; the second inequality then follows because the full
Gamma form contains $\mathcal G_0$ as a nonnegative summand.
\end{proof}

\paragraph{A three-cell connected cluster.}
Take $N=10$, word depth $k=1$, and
\begin{equation}\label{eq:newdThreeCellLabels}
 q_1=\frac34,
 \qquad q_2=\frac79,
 \qquad q_3=\frac45,
 \qquad
 \delta_{10}=\frac12\log\frac{11}{10}.
\end{equation}
The three coefficients obtained from
\eqref{eq:newdLeakageCoefficient} are
\begin{equation}\label{eq:newdThreeCellCoefficients}
 c_1=\frac{\log2\log3}{\sqrt{12}\log10},
 \qquad
 c_2=\frac{\log3\log7}{\sqrt{63}\log10},
 \qquad
 c_3=\frac{\log5\log2}{\sqrt{20}\log10}.
\end{equation}
They are all nonzero.  Moreover
\begin{equation}\label{eq:newdThreeCellConnections}
 \log\frac{q_2}{q_1}=\log\frac{28}{27}<\delta_{10},
 \qquad
 \log\frac{q_3}{q_2}=\log\frac{36}{35}<\delta_{10},
\end{equation}
so the translated cells overlap in a connected chain.  For
$e\in L^2(0,\delta_{10})$ define their coherent physical output
\begin{equation}\label{eq:newdThreeCellOutput}
 f_{10,e}=\sum_{j=1}^3c_j\tau_{\log q_j}e.
\end{equation}
Its support is contained in the single exterior hull
\begin{equation}\label{eq:newdThreeCellHull}
 J_{10}=\left(\log\frac34,
              \log\frac45+\delta_{10}\right),
 \qquad
 D_{10}=|J_{10}|=\log\frac{16}{15}
                 +\frac12\log\frac{11}{10}.
\end{equation}

\begin{theorem}[Complete payment of a connected three-cell cluster]
\label{thm:newdThreeCellCluster}
For every $e\in L^2(0,\delta_{10})$, including profiles for which the
adjacent translated copies interfere coherently,
\begin{equation}\label{eq:newdThreeCellPayment}
 \|f_{10,e}\|_2^2
 \leq\mathcal G_\Gamma[f_{10,e}]
      +\mathcal T_{J_{10}}[f_{10,e}].
\end{equation}
More precisely,
\begin{equation}\label{eq:newdThreeCellStrictPayment}
 \|f_{10,e}\|_2^2
 \leq\alpha_{10}\mathcal G_0[f_{10,e}]
      +\eta_{10}\mathcal T_{J_{10}}[f_{10,e}],
\end{equation}
where
\begin{align}
 \alpha_{10}
 &=\frac14+\frac1{32}
   \left(\log\frac{16}{15}
        +\frac12\log\frac{11}{10}\right)^2
   <0.251,\notag\\
 \eta_{10}
 &=\frac12\left[1-
   \exp\!\left(-\log\frac{16}{15}
               -\frac12\log\frac{11}{10}\right)\right]
   <0.054.
 \label{eq:newdThreeCellConstants}
\end{align}
Consequently
\begin{align}
 &\mathcal G_\Gamma[f_{10,e}]
  +\mathcal T_{J_{10}}[f_{10,e}]-\|f_{10,e}\|_2^2
 \notag\\
 &\qquad\geq
 (1-\alpha_{10})\mathcal G_0[f_{10,e}]
 +(1-\eta_{10})\mathcal T_{J_{10}}[f_{10,e}]\geq0.
 \label{eq:newdThreeCellRemainder}
\end{align}
\end{theorem}

\begin{proof}
The coefficient formulas follow because, for each pair in
\eqref{eq:newdThreeCellLabels}, the only contributing divisor in
\eqref{eq:newdLeakageCoefficient} is $d=1$.  Equations
\eqref{eq:newdThreeCellConnections} prove connectedness, and the extreme
endpoints give \eqref{eq:newdThreeCellHull}.  Numerically,
\[
 D_{10}=0.1121936110\ldots<\sqrt{24},
 \quad
 \alpha_{10}=0.2503933564\ldots,
 \quad
 \eta_{10}=0.0530644113\ldots.
\]
Apply Lemma~\ref{lem:newdClusterHull} directly to the coherent sum
$f_{10,e}$, before expanding its squared norm.  Therefore every diagonal
term and every interference term
\[
 c_i\overline{c_j}
 \langle\tau_{\log q_i}e,\tau_{\log q_j}e\rangle
\]
is included automatically.  The Gamma state is solved once from the left
endpoint of $J_{10}$ to its right endpoint and is never reset at either
internal cell boundary.  This proves
\eqref{eq:newdThreeCellStrictPayment}; the remaining statements follow as
in Lemma~\ref{lem:newdClusterHull}.
\end{proof}

Thus the earlier warning is not an obstruction to finite clusters: it only
forbids proving a cluster estimate by summing independently reset one-cell
estimates.  The exterior-hull argument proves every connected cluster with
hull length $D<\sqrt{24}$.  The following arithmetic separation lemma shows
that this restriction is automatic for every connected rational-leakage
cluster.  The point is that the cells respect a fixed family of reciprocal
half-integer bands.

\begin{theorem}[Reciprocal half-integer band theorem]
\label{thm:newdHalfIntegerBands}
Fix $N\geq2$ and a word depth $k$.  Form the intervals
\[
 J_q=\log q+(0,\delta_N),
 \qquad
 q=a/b\in\mathscr R_N^{\rm leak},
 \qquad \ell_{N,k}(a,b)>0.
\]
Every connected component of $\bigcup_qJ_q$ is contained in one of the
bands
\begin{equation}\label{eq:newdArithmeticBands}
 \left(\log\frac m2,\log\frac{m+1}{2}\right)
 \quad(m\geq2),
 \qquad\text{or}\qquad
 \left(\log\frac2{m+1},\log\frac2m\right)
 \quad(m\geq2),
\end{equation}
up to the possible inclusion of one boundary cell.  In particular, the
exterior hull of every connected component has length
\begin{equation}\label{eq:newdUniformClusterDiameter}
 D_{N,k}^{\rm comp}
 \leq \log\frac32+\delta_N
 \leq\frac32\log\frac32
 <\sqrt{24}.
\end{equation}
Consequently every connected component $C$ and every profile
$e\in L^2(0,\delta_N)$ satisfy, for
\[
 f_{C,e}=\sum_{q:\,J_q\subset C}
          \ell_{N,k}(q)\tau_{\log q}e,
\]
the complete estimate
\begin{equation}\label{eq:newdEveryConnectedClusterPaid}
 \|f_{C,e}\|_2^2
 \leq\mathcal G_\Gamma[f_{C,e}]
      +\mathcal T_{\operatorname{hull}(C)}[f_{C,e}].
\end{equation}
This assertion is uniform in $N$, $k$, the number of cells in $C$, and
all coherent interference among them.
\end{theorem}

\begin{proof}
We first prove that a rational cell cannot cross the indicated barriers.
Let $q=a/b>1$ be reduced, with $a,b\leq N$, and let $r=m/2>q$ be a
half-integer.  Since $mb-2a$ is a positive integer,
\[
 \log\frac rq
 =\log\frac{mb}{2a}
 \geq\log\left(1+\frac1{2a}\right)
 \geq\log\left(1+\frac1{2N}\right).
\]
But
\begin{equation}\label{eq:newdBarrierMargin}
 \log\left(1+\frac1{2N}\right)
 >\frac12\log\left(1+\frac1N\right)=\delta_N,
\end{equation}
because $(1+1/(2N))^2>1+1/N$.  Hence the right endpoint of $J_q$ lies
strictly to the left of $\log r$.

For $0<q=a/b<r=2/m$, the positive integer $2b-ma$ gives
\[
 \log\frac rq
 =\log\frac{2b}{ma}
 \geq\log\left(1+\frac1{ma}\right).
\]
Here $ma=2b-(2b-ma)\leq2N-1$, and therefore the last quantity is at least
$\log(1+1/(2N-1))$, which is strictly larger than the left side of
\eqref{eq:newdBarrierMargin}, hence larger than $\delta_N$.  Thus no cell
below a reciprocal barrier $2/m$ crosses it either.  Cells starting on a
barrier belong only to the band on its right, so they do not join the two
adjacent open bands.  Notice that this proof uses only $a,b\leq N$; it is
therefore valid at every word depth and does not depend on which of the
coefficients happen to vanish.

The consecutive ratios of the band endpoints are
\[
 \frac{(m+1)/2}{m/2}=1+\frac1m,
 \qquad
 \frac{2/m}{2/(m+1)}=1+\frac1m.
\]
Their logarithms are at most $\log(3/2)$ for $m\geq2$.  Adding one cell
length gives the first inequality in
\eqref{eq:newdUniformClusterDiameter}; the second follows from
$\delta_N\leq\delta_2=\tfrac12\log(3/2)$.  Lemma
\ref{lem:newdClusterHull} now applies once to the complete coherent output
on each connected component and proves
\eqref{eq:newdEveryConnectedClusterPaid}.
\end{proof}

\begin{corollary}[No low-frequency fugitive inside one arithmetic band]
\label{cor:newdBandCoherence}
For a connected component $C$ define its positive exponential polynomial
\[
 P_C(\tau)=\sum_{q:\,J_q\subset C}\ell_{N,k}(q)q^{-i\tau}.
\]
For every $0\leq\omega<\pi/\log(3/2)$ and $|\tau|\leq\omega$ one has
\begin{equation}\label{eq:newdBandCoherence}
 c(\omega)P_C(0)\leq |P_C(\tau)|\leq P_C(0),
 \qquad
 c(\omega)=
 \cos\!\left(\frac{\omega}{2}\log\frac32\right)>0.
\end{equation}
Thus the translates belonging to a single arithmetic band cannot, at any
$N$ or word depth, form a low-frequency fugitive profile by destructive
interference.
\end{corollary}

\begin{proof}
Let $x_q=\log q$ and let $x_C$ be the midpoint of the smallest interval
containing all $x_q$ in $C$.  The band theorem gives
$|x_q-x_C|\leq\tfrac12\log(3/2)$.  Positivity of the coefficients yields
\begin{align*}
 |P_C(\tau)|
 &\geq \Re\!\left(e^{i\tau x_C}P_C(\tau)\right)\\
 &=\sum_{q\in C}\ell_{N,k}(q)
       \cos\bigl(\tau(x_q-x_C)\bigr)\\
 &\geq
 \cos\!\left(\frac{\omega}{2}\log\frac32\right)
 \sum_{q\in C}\ell_{N,k}(q).
\end{align*}
The cosine is positive by the assumed range of $\omega$.  The upper bound
is the triangle inequality.
\end{proof}

For orientation, direct enumeration of the exact nonzero coefficients at
depth one gives the following largest connected hulls.  These figures are
not used in the proof; they display the approach to the sharp first-band
width $\log(3/2)$.
\[
\begin{array}{c|c|c|c}
N&\text{leftmost label}&\text{rightmost label}&
 D_{N,1}^{\rm largest}\\ \hline
10&6/5&4/3&0.1530156\ldots\\
50&33/49&48/49&0.3845948\ldots\\
200&133/199&198/199&0.4004117\ldots\\
1000&665/997&996/997&0.4044600\ldots
\end{array}
\]

The theorem removes the possibility of an arbitrarily long connected
cluster, and Corollary~\ref{cor:newdBandCoherence} removes destructive
low-frequency interference inside each such cluster.  It does not by
itself permit one to add
\eqref{eq:newdEveryConnectedClusterPaid} over disconnected bands: the
unreset Gamma form and the two global Tate moments contain cross-band
terms.  The remaining question is therefore a recombination theorem for
the separated arithmetic bands---equivalently, exclusion of cancellation
between distinct band polynomials---not a further one-cell or
connected-cluster estimate.

\paragraph{Audit of causal recombination across the gaps.}
There is a tempting but incorrect argument which treats decay of the causal
state in a gap as a negative Schur remainder.  The exact sign is the
opposite.  To record this point unambiguously, write
\[
 A_m=z(\inf C_m),\qquad B_m=z(\sup C_m),\qquad
 A_{m+1}=e^{-\Delta_m/2}B_m,qquad A_1=0.
\]
Summing the exact cluster identities gives
\begin{align}
 \frac12\sum_{m=1}^M(|B_m|^2-|A_m|^2)
 &=\frac12|B_M|^2
   +\frac12\sum_{m=1}^{M-1}
      (1-e^{-\Delta_m})|B_m|^2.
 \label{eq:newdGapAuditSign}
\end{align}
Thus the gap term occurs with a plus sign, not a minus sign.  Equivalently,
on a gap the equation $z'=-z/2$ gives
\[
 4\int_{\rm gap}|z'|^2
 =\int_{\rm gap}|z|^2
 =(1-e^{-\Delta_m})|B_m|^2,
\]
and the two positive bulk terms are cancelled exactly by the decrease of
the endpoint storage when the identity is written on the complete hull;
they do not produce an additional negative remainder.  Indeed the global
identity remains
\begin{equation}\label{eq:newdGlobalStateAuditIdentity}
 \|f\|_2^2
 =\frac14\mathcal G_0[f]
  +\frac14\|z\|_2^2+rac12|B_M|^2.
\end{equation}

There is a second independent obstruction to that argument.  The anchored
estimate
\[
 \|z\|_{L^2(C_m)}^2
 \leq\frac{|C_m|^2}{2}\|z'\|_{L^2(C_m)}^2
\]
used in Lemma~\ref{lem:newdClusterHull} requires $A_m=0$ and cannot be
reused for the restriction of the global state to later clusters.  For
example, the homogeneous state $z(t)=Ae^{-(t-L)/2}$ on an interval of
length $D$ has
\[
 \|z'\|_2^2=\frac14\|z\|_2^2,
\]
so the displayed estimate would require $1\leq D^2/8$, which is false for
all the band lengths in \eqref{eq:newdUniformClusterDiameter}.  An
arbitrarily small nonzero source perturbation gives the same failure inside
a genuine later support component.  Therefore causal decay by itself does
not prove the required recombination; a valid proof must control the
incoming states or establish the global low-frequency observability
inequality below.

\paragraph{The single obstruction for arbitrary clusters.}
The preceding results permit an exact separation of the disconnected-band
problem.  Since every summand in \eqref{eq:newdGammaMultiplier} is strictly
increasing for $\tau>0$, while $g_\Gamma(0)=0$ and
$g_\Gamma(\tau)\to\infty$, there is a unique number
$\tau_\Gamma>0$ such that
\begin{equation}\label{eq:newdGammaCutoff}
 g_\Gamma(\tau_\Gamma)=1.
\end{equation}
Numerically, $\tau_\Gamma=0.2860886323\ldots$; the numerical value is not
used below.  Put $\Omega_\Gamma=(-\tau_\Gamma,\tau_\Gamma)$ and define the
positive and negative Gamma weights
\[
 w_-(\tau)=(1-g_\Gamma(\tau))_+,
 \qquad
 w_+(\tau)=(g_\Gamma(\tau)-1)_+.
\]
In fact $\tau_\Gamma<1$, because the $j=0$ summand alone gives
$g_\Gamma(1)>1$, and $1<\pi/\log(3/2)$.  Hence
Corollary~\ref{cor:newdBandCoherence} applies on the whole deficit window
$\Omega_\Gamma$: no individual reciprocal half-integer band has a
low-frequency zero or even an arbitrarily small normalized value there.
For the $k$th rational leakage output write
\begin{equation}\label{eq:newdLeakagePolynomial}
 P_{N,k}(\tau)=
 \sum_{a/b\in\mathscr R_N^{\rm leak}}
 \ell_{N,k}(a,b)e^{-i\tau\log(a/b)},
 \qquad
 \widehat{L_{N,k}e}(\tau)=\widehat e(\tau)P_{N,k}(\tau).
\end{equation}

\begin{proposition}[Exact low-frequency reduction]
\label{prop:newdLowFrequencyReduction}
The uniform Gamma--Tate estimate
\eqref{eq:newdUniformSourceEstimate} is equivalent to
\begin{align}
 &\sum_{k\geq1}\int_{\Omega_\Gamma}
    w_-(\tau)|\widehat e(\tau)P_{N,k}(\tau)|^2
       \frac{d\tau}{2\pi}                                \notag\\
 &\qquad\leq
 \sum_{k\geq1}\mathcal T_{I_N}[L_{N,k}e]
 +\sum_{k\geq1}\int_{\R\setminus\Omega_\Gamma}
    w_+(\tau)|\widehat e(\tau)P_{N,k}(\tau)|^2
       \frac{d\tau}{2\pi}.
 \label{eq:newdLowFrequencyObservability}
\end{align}
Moreover, because all coefficients
$\ell_{N,k}(a,b)$ are nonnegative,
\begin{equation}\label{eq:newdPositiveCoefficientEnvelope}
 |P_{N,k}(\tau)|^2
 \leq P_{N,k}(i/2)P_{N,k}(-i/2)
 \qquad(\tau\in\R).
\end{equation}
Thus every frequency outside the fixed compact band
$\Omega_\Gamma$ is already paid by Gamma.  The only remaining issue is the
two-moment observability of the arithmetic leakage range inside that band.
\end{proposition}

\begin{proof}
For $F_k=L_{N,k}e$, Plancherel and
\eqref{eq:newdFullGammaEnergy} give
\begin{align*}
 &\mathcal G_\Gamma[F_k]+\mathcal T_{I_N}[F_k]-\|F_k\|_2^2\\
 &\quad=\mathcal T_{I_N}[F_k]
 +\int_\R(g_\Gamma(\tau)-1)|\widehat F_k(\tau)|^2
          \frac{d\tau}{2\pi}.
\end{align*}
Splitting the integral at $\Omega_\Gamma$, substituting
\eqref{eq:newdLeakagePolynomial}, and summing over $k$ proves the
equivalence.  For the envelope, write $q=a/b$ and use Cauchy--Schwarz:
\begin{align*}
 |P_{N,k}(\tau)|
 &\leq\sum_q\ell_{N,k}(q)\\
 &=\sum_q
   \sqrt{\ell_{N,k}(q)q^{1/2}}
   \sqrt{\ell_{N,k}(q)q^{-1/2}}\\
 &\leq
 \left(\sum_q\ell_{N,k}(q)q^{1/2}\right)^{1/2}
 \left(\sum_q\ell_{N,k}(q)q^{-1/2}\right)^{1/2}.
\end{align*}
The two factors are $P_{N,k}(i/2)^{1/2}$ and
$P_{N,k}(-i/2)^{1/2}$, proving
\eqref{eq:newdPositiveCoefficientEnvelope}.
\end{proof}

\paragraph{The mean/oscillation Schur decomposition.}
The preceding band calculation suggests a sharper way to formulate the
remaining global problem.  Put
\[
 \mathfrak A_N[e]=\|\mathbf L_Ne\|_2^2,
 \qquad
 \mathfrak Q_N[e]=\mathcal E_{\Gamma T,N}(e),
 \qquad
 \mathfrak D_N=\mathfrak Q_N-\mathfrak A_N,
\]
and decompose the input cell exactly as
\begin{equation}\label{eq:newdMeanZeroSplit}
 E_N=\C u_N\oplus E_N^0,
 \qquad
 u_N=\delta_N^{-1/2}\mathbf1_{(0,\delta_N)},
 \qquad
 E_N^0=\left\{h:\int_0^{\delta_N}h=0\right\}.
\end{equation}
Let $d_N=\mathfrak D_N[u_N]$, let
$b_N(h)=\mathfrak D_N(u_N,h)$, and let $C_N$ be the restriction of
$\mathfrak D_N$ to $E_N^0$.

\begin{proposition}[Exact mean/oscillation recombination criterion]
\label{prop:newdMeanZeroSchur}
The unit Gamma--Tate estimate at threshold $N$ holds if and only if
\begin{equation}\label{eq:newdMeanZeroSchurCriterion}
 C_N\geq0,
 \qquad
 b_N\in\operatorname{Ran}C_N^{1/2},
 \qquad
 d_N-b_N^*C_N^\dagger b_N\geq0.
\end{equation}
Thus the remaining interference problem splits into a coercivity statement
for zero-mean cell profiles and one scalar Schur inequality for the constant
profile after all oscillatory profiles have been eliminated.
\end{proposition}

\begin{proof}
In the orthogonal decomposition \eqref{eq:newdMeanZeroSplit}, the defect
form has block matrix
\[
 \mathfrak D_N=
 \begin{pmatrix}d_N&b_N^*\\ b_N&C_N\end{pmatrix}.
\]
The generalized Schur-complement criterion for a closed Hermitian form is
exactly \eqref{eq:newdMeanZeroSchurCriterion}.  Positivity of
$\mathfrak D_N$ is the unit Gamma--Tate estimate by definition.
\end{proof}

\begin{conjecture}[Logarithmic Schur angle]
\label{conj:newdLogSchurAngle}
For every $N\geq3$ the mean-zero block and constant margin satisfy
\begin{equation}\label{eq:newdLogSchurConjecture}
 C_N\geq0,
 \qquad d_N>0,
 \qquad b_N\in\operatorname{Ran}C_N^{1/2},
 \qquad
 \rho_N:=\frac{b_N^*C_N^\dagger b_N}{d_N}
 \leq\frac1{20\log N}.
\end{equation}
At $N=2$ the rational-leakage output is empty.
\end{conjecture}

\begin{corollary}[Conditional completion of row (d)]
\label{cor:newdConditionalRowD}
Assume Conjecture~\ref{conj:newdLogSchurAngle}.  Then the unit Gamma--Tate
estimate holds at every arithmetic threshold.  Consequently the threshold
Schur complements are nonnegative, the primitive inequality propagates to
all windows, and row~(d) follows.
\end{corollary}

\begin{proof}
For $N\geq3$, equation \eqref{eq:newdLogSchurConjecture} gives
\[
 d_N-b_N^*C_N^\dagger b_N
 \geq d_N\left(1-\frac1{20\log N}\right)>0.
\]
Together with the asserted block positivity and range condition,
Proposition~\ref{prop:newdMeanZeroSchur} gives the unit source estimate.
The exact threshold completion criterion and induction from the certified
initial interval give the stated propagation.  The $N=2$ leakage block is
zero.
\end{proof}

\paragraph{Finite-dimensional reconnaissance (audit only).}
To identify which part of Proposition~\ref{prop:newdMeanZeroSchur} carries
the difficulty, the exact arithmetic coefficients were evaluated on the
space of twelve-bin piecewise-constant profiles on $(0,\delta_N)$.  The
table reports the smallest generalized quotient
$\mathfrak Q_N/\mathfrak A_N$, the same quotient after imposing zero mean,
and the proportion of the constant defect consumed by its Schur correction.
\[
\begin{array}{c|c|c|c}
N&\min \mathfrak Q_N/\mathfrak A_N&
\min_{E_N^0}\mathfrak Q_N/\mathfrak A_N&
b_N^*C_N^{-1}b_N/d_N\\ \hline
10 &5.1436&8.2823&0.0156\\
20 &4.9506&8.9386&0.0175\\
40 &4.8129&9.6305&0.0150\\
80 &4.5263&10.2950&0.0116\\
120&4.4563&10.6917&0.0118
\end{array}
\]
Mesh refinement at $N=80$ from six to sixteen bins changes the first
quotient from $4.5316$ to $4.5263$ and the zero-mean quotient from
$10.3691$ to $10.2807$.  These computations are neither interval-certified
nor a proof for arbitrary profiles.  Their consistent structural signal is
that the prospective minimizer lies in the almost-constant sector, while
the zero-mean sector has a large Gamma margin and the cross-sector Schur
correction is small.  Accordingly, the next analytic target is
\eqref{eq:newdMeanZeroSchurCriterion}, not an unstructured estimate over
all cell profiles at once.

The restriction to the arithmetic range in
\eqref{eq:newdLowFrequencyObservability} is essential.  Indeed, choose
$\varphi\in C_c^\infty(-1/4,1/4)$ with $\|\varphi\|_2=1$ and put
\begin{equation}\label{eq:newdPlateauFamily}
 f_D(t)=D^{-1/2}\varphi(t/D),
 \qquad I_D=(-D/2,D/2).
\end{equation}
Then $\|f_D\|_2=1$, whereas
\begin{equation}\label{eq:newdPlateauEscape}
 \mathcal G_\Gamma[f_D]\longrightarrow0,
 \qquad
 \mathcal T_{I_D}[f_D]\longrightarrow0
 \qquad(D\to\infty).
\end{equation}
For the first limit, substitute
$\widehat f_D(\tau)=D^{1/2}\widehat\varphi(D\tau)$ and use
$g_\Gamma(u/D)\to0$ together with the rapid decay of
$\widehat\varphi$.  For the second, the support of $f_D$ stays a distance
$D/4$ from both endpoints of $I_D$: its two exponential moments grow at
most like $D^{1/2}e^{D/8}$, while the corresponding diagonal Tate Grams
grow like $e^{D/2}$; the off-diagonal entry is only $D$.  Hence the Tate
projection tends to zero exponentially.

Consequently no support-only argument can prove the unit estimate for an
aggregate whose disconnected bands have an arbitrarily long exterior
hull.  The one remaining theorem is precisely that
outputs of the positive arithmetic polynomials
$P_{N,k}$ cannot form the escaping family
\eqref{eq:newdPlateauFamily}; equivalently, they must satisfy the compact
low-frequency observability inequality
\eqref{eq:newdLowFrequencyObservability} uniformly in $N$.

\subsubsection{The exact threshold Schur condition}

Let the two physical analysis factors at one threshold have old/new columns
\[
 X=(X_0,X_E),\qquad Y=(Y_0,Y_E),
\]
and put
\begin{equation}\label{eq:newdPhysicalBlocks}
 R_0=X_0^*X_0,
 \quad L_0=Y_0^*Y_0,
 \quad r=X_0^*X_E,
 \quad \ell=Y_0^*Y_E,
 \quad A_0=R_0-L_0.
\end{equation}
Assume the old primitive block $A_0$ is nonnegative and the reference range
condition $r\in\operatorname{Ran}R_0$ holds.  Define
\begin{align}
 H&=R_0^\dagger r,
 \label{eq:newdReferenceLift}\\
 S_E&=X_E^*X_E-r^*R_0^\dagger r,
 \label{eq:newdReferenceSchur}\\
 Z_E&=Y_E-Y_0H,
 \qquad
 b_E=L_0H-\ell.
 \label{eq:newdOppositeResidual}
\end{align}

\begin{theorem}[Exact threshold completion criterion]
\label{thm:newdThresholdCriterion}
The full old/new signed block $X^*X-Y^*Y$ is nonnegative if and only if
\begin{equation}\label{eq:newdRangeCondition}
 b_E\in\operatorname{Ran}A_0^{1/2}
\end{equation}
and
\begin{equation}\label{eq:newdSchurTarget}
 S_E-Z_E^*Z_E-b_E^*A_0^\dagger b_E\geq0.
\end{equation}
\end{theorem}

\begin{proof}
The equation $R_0H=r$ gives
$X_0^*(X_E-X_0H)=0$ and
\[
 (X_E-X_0H)^*(X_E-X_0H)=S_E.
\]
Hence the reference-harmonic column $v_E=(-H,I)^t$ satisfies
\[
 v_E^*(X^*X-Y^*Y)v_E=S_E-Z_E^*Z_E.
\]
Its old cross is
\[
 (r-\ell)-(R_0-L_0)H=L_0H-\ell=b_E.
\]
Congruence by the invertible triangular change from the coordinate newborn
column to $v_E$ leaves the old block equal to $A_0$ and changes the other
entries to $b_E$ and $S_E-Z_E^*Z_E$.  The generalized Schur-complement
criterion gives precisely \eqref{eq:newdRangeCondition} and
\eqref{eq:newdSchurTarget}.
\end{proof}

This theorem is the fixed global target.  The preceding construction
determines the divergence, the moving Tate cell, the complete positive Gamma
form and the exact rational leakage.  It does not determine the sign in
\eqref{eq:newdSchurTarget}.

The source construction isolates a concrete sufficient estimate.  Let
\[
 \mathscr D_N=
 \{e\in E_N:\mathcal E_{\Gamma T,N}(e)<\infty\}.
\]
It would suffice to prove, with the same unit constant for every arithmetic
threshold,
\begin{equation}\label{eq:newdUniformSourceEstimate}
 \|\mathbf L_Ne\|^2\leq\mathcal E_{\Gamma T,N}(e)
 \qquad(N\geq2,\ e\in\mathscr D_N).
\end{equation}
Indeed, Proposition~\ref{prop:newdLeakageSplit} separates the integer
incidence from the noninteger leakage.  The former is the finite
divisibility current of Theorem~\ref{thm:newdAmplitude}; the latter is
exactly $\mathbf L_Ne$ by
Proposition~\ref{prop:newdRationalAggregation}.  The right side of
\eqref{eq:newdUniformSourceEstimate} is the unreset positive Gamma storage
plus the conservative two-dimensional Tate terminal.  Thus the estimate
would make the joint boundary source contractive with unit normalization;
shorting a passive source network would give
\eqref{eq:newdRangeCondition}--\eqref{eq:newdSchurTarget} at the next
threshold.  Induction from the certified initial interval would then give
the primitive inequality for every $T$.

For clarity, \eqref{eq:newdUniformSourceEstimate} is a sufficient
source-level route, not a definition of the desired Schur complement and not
an additional conclusion of the identities already proved.  Expanded on the
test core $C_c^\infty(0,\delta_N)$, it is the explicit inequality
\begin{align}
 &\sum_{k\geq1}
 \sum_{a/b,c/d\in\mathscr R_N^{\rm leak}}
 \ell_{N,k}(a,b)\overline{\ell_{N,k}(c,d)}
 \kappa_e\!\left(\log\frac{ad}{bc}\right)
 \notag\\
 &\quad\leq
 \sum_{k\geq1}\int_\R g_\Gamma(\tau)
 \left|\widehat e(\tau)
   \sum_{a/b\in\mathscr R_N^{\rm leak}}
       \ell_{N,k}(a,b)e^{-i\tau\log(a/b)}\right|^2
 \frac{d\tau}{2\pi}
 \notag\\
 &\qquad+
 \sum_{k\geq1}
 \begin{pmatrix}
  \overline{M_+(L_{N,k}e)}&\overline{M_-(L_{N,k}e)}
 \end{pmatrix}
 G_{I_N}^{-1}
 \binom{M_+(L_{N,k}e)}{M_-(L_{N,k}e)}.
 \label{eq:newdExpandedSourceEstimate}
\end{align}
All quantities in \eqref{eq:newdExpandedSourceEstimate} have been defined
independently of the desired sign.  The unresolved issue is its uniform
validity, including coherent clusters of nearby rational labels; a cellwise
argument that resets the Gamma state does not prove it.

\begin{center}
\fbox{\parbox{0.91\linewidth}{
\textbf{Conditional closure ledger for row (d).}
The Gamma--Euler--Tate boundary object, rational bands, connected-cluster
payment and exact recombination block are constructed.  Global row~(d)
follows from Conjecture~\ref{conj:newdLogSchurAngle}, which asserts at every
threshold the range condition and the uniform Schur inequality
\[
 S_E-Z_E^*Z_E-b_E^*A_0^\dagger b_E\geq0.
\]
Equivalently, it asserts the unit-constant estimate
\eqref{eq:newdUniformSourceEstimate}.  Conditional on that named conjecture,
Corollary~\ref{cor:newdConditionalRowD} completes row~(d).  The conjecture,
the unconditional global primitive inequality and unconditional row~(d)
are not proved in the present manuscript.
}}
\end{center}

% END CLEAN GAMMA--FOURIER--TATE ROW (d)

\section*{Executable companion}

The source distribution includes the Google-Colab-compatible notebook
\path{certificates/arithmetic_lefschetz_certificates.ipynb} and separate
row-specific runners.  Together they check the exact finite arithmetic,
corrected negabinary admissible lengths, Boolean matrix-code separation,
cotangent ranks, determinant limits, contact determinants,
Dirichlet-convolution identities and the Lorentzian ruling matrix used above.
Every cell is labelled as an exact identity, a finite model or a
non-computational categorical theorem.  Successful execution is evidence for
the finite calculations, not a substitute for the proofs in
Section~\ref{sec:geometric-a}.  The same directory contains a plain-Python runner
which executes all code cells without Jupyter.
For the row-(d) endpoint \(T=\frac12\log5\), the subdirectory
\path{certificates/row_d_log5} contains a SHA-256-pinned frozen-interval
artifact and a standard-library verifier.  It audits the complete contact
list and the strict complement, nested-Schur and final directed-Gershgorin
signs used by the Feshbach proof.  It does not recompute the originating Arb
kernel enclosures; this narrower scope is stated in its README.

\begin{thebibliography}{99}

\bibitem{BermanEnriched}
J.~D.~Berman, \emph{Enriched infinity categories I: enriched presheaves},
arXiv:2008.11323, 2020.

\bibitem{BombieriLagarias}
E.~Bombieri and J.~C.~Lagarias, \emph{Complements to Li's criterion for the
Riemann hypothesis}, J. Number Theory \textbf{77} (1999), 274--287.

\bibitem{CC-ArithmeticSite}
A.~Connes and C.~Consani, \emph{The Arithmetic Site},
C. R. Acad. Sci. Paris \textbf{352} (2014), 971--975.

\bibitem{CC-AbsoluteAlgebra}
A.~Connes and C.~Consani, \emph{Absolute algebra and Segal's Gamma sets},
J. Number Theory \textbf{162} (2016), 518--551.

\bibitem{CC-ScalingSite}
A.~Connes and C.~Consani, \emph{Geometry of the Scaling Site},
Selecta Math. (N.S.) \textbf{23} (2017), 1803--1850.

\bibitem{CC-RRstrategy}
A.~Connes and C.~Consani, \emph{The Riemann--Roch strategy: complex lift of
the Scaling Site}, in: Advances in Noncommutative Geometry, Springer, 2019.

\bibitem{CC-RRSpecZ}
A.~Connes and C.~Consani, \emph{Riemann--Roch for
$\overline{\Spec\Z}$}, Bull. Sci. Math. \textbf{187} (2023), 103293.

\bibitem{CCJacobian}
A.~Connes and C.~Consani, \emph{On the Jacobian of
$\overline{\Spec\Z}$}, arXiv:2602.15941, 2026.

\bibitem{CCWeilArch}
A.~Connes and C.~Consani, \emph{Weil positivity and Trace formula, the
archimedean place}, Selecta Math. (N.S.) \textbf{27} (2021), no.~4, Paper
No.~77; arXiv:2006.13771.

\bibitem{CCM2025}
A.~Connes, C.~Consani and H.~Moscovici, \emph{Zeta zeros and prolate wave
operators}, preprint, 2024.

\bibitem{Connes1999}
A.~Connes, \emph{Trace formula in noncommutative geometry and the zeros of
the Riemann zeta function}, Selecta Math. (N.S.) \textbf{5} (1999), 29--106.

\bibitem{ClausenScholze}
D.~Clausen and P.~Scholze, \emph{Condensed Mathematics}, lecture notes,
2019.

\bibitem{CvS}
A.~Connes and W.~D.~van Suijlekom, \emph{Tolerance relations and operator
systems}, preprint.

\bibitem{Deninger2024}
C.~Deninger, \emph{Dynamical systems for arithmetic schemes},
Indag. Math., 2024.

\bibitem{Grothendieck}
A.~Grothendieck, \emph{Sur une note de Mattuck--Tate},
J. Reine Angew. Math. \textbf{200} (1958), 208--215.

\bibitem{Haran2022}
S.~Haran, \emph{Non-additive geometry and Frobenius correspondences},
arXiv:2209.08536, 2023.

\bibitem{Lagarias2007}
J.~C.~Lagarias, \emph{Li coefficients for automorphic $L$-functions},
Ann. Inst. Fourier \textbf{57} (2007), 1689--1740.

\bibitem{LurieHA}
J.~Lurie, \emph{Higher Algebra}, 2017, especially \S4.8.1.

\bibitem{KnudsenMumford}
F.~Knudsen and D.~Mumford, \emph{The projectivity of the moduli space of
stable curves I: Preliminaries on ``det'' and ``Div''}, Math. Scand.
\textbf{39} (1976), 19--55.

\bibitem{MattuckTate}
A.~Mattuck and J.~Tate, \emph{On the inequality of Castelnuovo--Severi},
Abh. Math. Sem. Univ. Hamburg \textbf{22} (1958), 295--299.

\bibitem{MeyerZeta}
R.~Meyer, \emph{A spectral interpretation for the zeros of the Riemann zeta
function}, in: \emph{Mathematisches Institut, Georg-August-Universit\"at
G\"ottingen: Seminars Winter Term 2004/2005}, Universit\"atsdrucke
G\"ottingen, 2005, 117--137; corrected version arXiv:math/0412277v3 (2013),
especially \S\S2--5.

\bibitem{Morishita2026}
M.~Morishita, \emph{On a relation between Deninger's foliated dynamical
systems and Connes--Consani's adelic spaces}, preprint, 2026.

\bibitem{PietschNuclear}
A.~Pietsch, \emph{Nuclear Locally Convex Spaces}, Ergebnisse der Mathematik
und ihrer Grenzgebiete, vol.~66, Springer, 1972.

\bibitem{Suzuki2025}
M.~Suzuki, \emph{An inverse problem for a class of canonical systems and its
applications}, preprint.

\bibitem{SuzukiScrew}
M.~Suzuki, \emph{Weil's quadratic form via the screw function},
arXiv:2606.09096, 2026.

\bibitem{Weil1948}
A.~Weil, \emph{Sur les courbes alg\'ebriques et les vari\'et\'es qui s'en
d\'eduisent}, Hermann, Paris, 1948.

\bibitem{Weil1952}
A.~Weil, \emph{Sur les ``formules explicites'' de la th\'eorie des nombres
premiers}, Comm. S\'em. Math. Univ. Lund (1952), 252--265.

\bibitem{YuanZhang}
X.~Yuan and S.-W.~Zhang, \emph{The arithmetic Hodge index theorem for
adelic line bundles I: number fields}, Math. Ann. \textbf{367} (2017),
1123--1171.

\bibitem{GeerSchoof}
G.~van der Geer and R.~Schoof, \emph{Effectivity of Arakelov divisors and the
theta divisor of a number field}, Selecta Math. (N.S.) \textbf{6} (2000),
377--398.

\bibitem{He2025}
W.~He, \emph{Numerical cohomology for arithmetic surfaces and applications},
arXiv:2512.01811, 2025.

\bibitem{Yoshida1992}
H.~Yoshida, \emph{On Hermitian forms attached to zeta functions}, in: Zeta
functions in geometry (Tokyo, 1990), 281--325, Adv. Stud. Pure Math.
\textbf{21}, Kinokuniya, Tokyo, 1992.

\end{thebibliography}

\end{document}
